\chapter{The Quot scheme}
\label{chap:Picard_QuotScheme}
% archon:covers AlgebraicJacobian/Picard/QuotScheme.lean AlgebraicJacobian/Picard/GradedHilbertSerre.lean

% STRATEGY NOTE. This chapter packages the Quot-foundations layer of the
% Grothendieck--Altman--Kleiman Quot-scheme construction
% \(\mathrm{Quot}^{\Phi,L}_{E/X/S}\): the per-fiber Hilbert polynomial, the
% Quot functor of \(T\)-flat coherent quotients of \(E_T\) on
% \(X_T = X \times_S T\), and the Grassmannian \emph{scheme} together with
% its representability. The Quot functor underlies the relative Picard
% functor via the Hilbert scheme \(\Hilb_{C/k} = \Quot_{\mathcal{O}_C/C/k}\)
% -- the case \(E = \mathcal{O}_X\) of the Quot functor.
%
% LOAD-BEARING SUB-PREREQUISITE. The Quot construction embeds Quot as a
% locally closed subfunctor of a Grassmannian
% \(\mathrm{Grass}(W \otimes_{\mathcal{O}_S} \mathrm{Sym}^r V, \Phi(r))\).
% Mathlib (at the pinned commit) carries no Grassmannian \emph{scheme} -- only
% a linear-algebra Grassmannian as a finite-rank-quotient variety. The
% scheme-level Grassmannian over an arbitrary noetherian base is therefore an
% in-project sub-build. This chapter introduces it as
% \cref{def:grassmannian_scheme} and states its representability
% (\cref{thm:grassmannian_representable}).

\section{Setup and motivation}
\label{sec:quot_setup}

Throughout this chapter, let \(S\) be a noetherian scheme, \(\pi : X \to S\) a
projective morphism, \(L\) a relatively very ample line bundle on \(X\) over \(S\),
and \(E\) a coherent sheaf on \(X\). The Quot construction packages the
\(T\)-flat coherent quotients \(E_T \twoheadrightarrow \mathcal{F}\) of the
pull-back \(E_T\) to \(X_T = X \times_S T\), modulo the equivalence
\(\langle \mathcal{F}, q\rangle = \langle \mathcal{F}', q'\rangle\) defined by
\(\ker(q) = \ker(q')\). Fixing the Hilbert polynomial
\(\Phi(\lambda) = \chi(\mathcal{F}|_{X_t}(\lambda)) \in \mathbb{Q}[\lambda]\)
(constant on connected \(T\) by flatness)
stratifies the Quot functor into pieces \(\Quot^{\Phi,L}_{E/X/S}\), each of
which Grothendieck proved is representable by a projective \(S\)-scheme.
This chapter develops the foundations of that picture: the Hilbert
polynomial (\cref{def:hilbert_polynomial}), the Quot functor
(\cref{def:quot_functor}), and the Grassmannian scheme
(\cref{def:grassmannian_scheme}) with its representability
(\cref{thm:grassmannian_representable}).

% SOURCE: [Nitsure], "Construction of Hilbert and Quot Schemes", \S 1;
% FGA Explained Ch.~5, AMS MSM 123 (read from
% references/nitsure-hilbert-quot-src/nitsure-hilbert-quot.tex, L287-L500).
% The Hartshorne reference for the polynomial-eventually
% property of Euler characteristics is [Hartshorne], III.5.2 (the scanned
% PDF at references/hartshorne-algebraic-geometry.pdf has no text layer, so
% the verbatim source-quote is taken from
% Nitsure \S 1; Hartshorne III.5.2 is cited in prose as the textbook
% reference).

\section{Hilbert polynomial of a coherent sheaf}
\label{sec:hilbert_polynomial}

\begin{definition}\leanok
[Hilbert polynomial]
  \label{def:hilbert_polynomial}
  \lean{AlgebraicGeometry.Scheme.hilbertPolynomial}
  \uses{thm:hilbertPoly_of_sectionModule}
  % ENCODING. `def:hilbert_polynomial` is formalized via the GRADED HILBERT
  % FUNCTION of the section module ⊕_{m≥0} Γ(X_s, F_s ⊗ L_s^{⊗m})
  % (def:sectionGradedModule), whose polynomiality is extracted in
  % thm:hilbertPoly_of_sectionModule from the graded Hilbert–Serre rationality
  % lemma (lem:gradedHilbertSerre_rational) together with the Mathlib lemma
  % Polynomial.existsUnique_hilbertPoly (lem:hilbertPoly_exists_mathlib). The
  % construction uses only H^0; for m≫0 it agrees with the Euler characteristic
  % χ by Serre vanishing, so it is the classical Hilbert polynomial and no
  % numerical invariant is lost.
  % SOURCE: [Nitsure], \S 1, sub-section "Stratification by Hilbert
  % Polynomials" (read from
  % references/nitsure-hilbert-quot-src/nitsure-hilbert-quot.tex, L453-L478).
  % SOURCE QUOTE:
  % "Let \(X\) be a finite type scheme over
  %  a field \(k\), together with a line bundle \(L\).
  %  Recall that if \(F\) is a coherent sheaf on
  %  \(X\) whose support is proper over \(k\), then
  %  the \textbf{Hilbert polynomial} \({\Phi}\in \Q[\lambda]\) of \(F\) is defined by
  %  the function
  %  \(\){\Phi}(m) = \chi(F(m)) =
  %  \sum_{i=0}^n (-1)^i \dim_k H^i(X, F\otimes L^{\otimes m})\(\)
  %  where the dimensions of the cohomologies are finite because of the
  %  coherence and properness conditions. The fact that \(\chi(F(m))\)
  %  is indeed a polynomial in \(m\) under the above assumption is a special
  %  case of what is known as Snapper's Lemma (see Kleiman [K] for a proof).
  %
  %  Let \(X\to S\) be a finite type morphism of noetherian schemes,
  %  and let \(L\) be a line bundle on \(X\). Let \(\F\)
  %  be any coherent sheaf on
  %  \(X\) whose schematic support is proper over \(S\). Then for each \(s\in S\),
  %  we get a polynomial \({\Phi}_s \in  \Q[\lambda]\)
  %  which is the Hilbert polynomial of
  %  the restriction \(\F_s = F|_{X_s}\) of \(\F\) to the fiber \(X_s\) over \(s\),
  %  calculated with respect to the line bundle \(L_s = L|_{X_s}\).
  %  If \(\F\) is flat over \(S\) then the function \(s\mapsto {\Phi}_s\)
  %  from the set of points of \(S\) to the polynomial ring \(\Q[\lambda]\) is
  %  known to be locally constant on \(S\)."
  \textit{Source: [Nitsure], \S 1 (FGA Explained Ch.~5); cf.\ [Hartshorne],
  III.5.2.}
  Let \(X \to S\) be a finite-type morphism of noetherian schemes and let \(L\)
  be a line bundle on \(X\). For a coherent sheaf \(\mathcal{F}\) on \(X\) whose
  schematic support is proper over \(S\) and a point \(s \in S\), write \(X_s\)
  for the fibre over \(s\), \(\mathcal{F}_s = \mathcal{F}|_{X_s}\) and
  \(L_s = L|_{X_s}\). The \emph{Hilbert polynomial of \(\mathcal{F}\) at \(s\)
  relative to \(L\)} is the unique polynomial
  \(\Phi_{\mathcal{F},s} \in \mathbb{Q}[\lambda]\) that agrees, for all
  \(m \gg 0\), with the graded Hilbert function
  \[
    m \;\longmapsto\; \dim_{\kappa(s)} \Gamma\bigl(X_s,\, \mathcal{F}_s \otimes L_s^{\otimes m}\bigr)
  \]
  of the section module
  \(\bigoplus_{m \ge 0} \Gamma(X_s, \mathcal{F}_s \otimes L_s^{\otimes m})\)
  (\cref{def:sectionGradedModule}), viewed as a finitely generated graded module
  over the homogeneous section ring
  \(\bigoplus_{m \ge 0} \Gamma(X_s, L_s^{\otimes m})\)
  (\cref{def:sectionGradedRing}). Existence and uniqueness of
  \(\Phi_{\mathcal{F},s}\) are the content of
  \cref{thm:hilbertPoly_of_sectionModule}.

  This is the classical Hilbert polynomial. When \(\mathcal{F}_s\) is coherent
  and \(L_s\) is ample on the proper \(\kappa(s)\)-scheme \(X_s\), Serre vanishing
  gives \(H^i(X_s, \mathcal{F}_s \otimes L_s^{\otimes m}) = 0\) for all \(i > 0\)
  and \(m \gg 0\), so for \(m \gg 0\)
  \[
    \dim_{\kappa(s)} \Gamma(X_s, \mathcal{F}_s \otimes L_s^{\otimes m})
    \;=\; \chi(X_s,\, \mathcal{F}_s \otimes L_s^{\otimes m})
    \;=\; \sum_{i \ge 0} (-1)^i \dim_{\kappa(s)} H^i(X_s,\, \mathcal{F}_s \otimes L_s^{\otimes m}),
  \]
  and hence \(\Phi_{\mathcal{F},s}\) coincides with the Euler-characteristic
  Hilbert polynomial of Snapper's Lemma -- no numerical invariant is lost. The
  construction above, however, uses only the global sections \(H^0\): the higher
  cohomology enters solely through the classical statement that the two functions
  agree for \(m \gg 0\), never through the definition of \(\Phi_{\mathcal{F},s}\)
  itself. When \(\mathcal{F}\) is \(S\)-flat the function
  \(s \mapsto \Phi_{\mathcal{F},s}\) is locally constant on \(S\), so \(\Phi\) is
  well-defined on connected components.
\end{definition}

\section{Graded Hilbert polynomial}
\label{sec:graded_hilbert_polynomial}

This section gives the graded-Hilbert-function construction that underlies
\cref{def:hilbert_polynomial}. Working over a single fibre \(X_s\) with its
ample line bundle \(L_s\) and coherent sheaf \(\mathcal{F}_s\), it builds the
section graded ring (\cref{def:sectionGradedRing}) and module
(\cref{def:sectionGradedModule}), records their finite generation
(\cref{lem:sectionGradedModule_fg}), and extracts the Hilbert polynomial
(\cref{thm:hilbertPoly_of_sectionModule}) by combining the graded
Hilbert--Serre rationality of the Hilbert series
(\cref{lem:gradedHilbertSerre_rational}) with Mathlib's polynomial-extraction
lemma (\cref{lem:hilbertPoly_exists_mathlib}). The route uses only \(H^0\)
(global sections), so it is independent of higher coherent cohomology.

\begin{definition}
[Section graded ring]
  \label{def:sectionGradedRing}
  \lean{AlgebraicGeometry.sectionGradedRing}
  \uses{lem:sectionGradedRing_gcommSemiring}
  % NOTE: Formalization is blocked on the absence of tensor products for sheaves of modules in
  % Mathlib at the pinned commit. The degree-\(m\) component \(\Gamma(X_s, L_s^{\otimes m})\)
  % requires a tensor power \(L_s^{\otimes m}\) of a line bundle; this infrastructure (tensor
  % product of sheaves of modules with lax-monoidal global sections) is not yet available, making
  % this definition a prerequisite sub-build.
  % SOURCE: [Nitsure], "Construction of Hilbert and Quot Schemes", \S 1,
  % sub-section "Stratification by Hilbert Polynomials" (read from
  % references/nitsure-hilbert-quot-src/nitsure-hilbert-quot.tex, L453-L464).
  % SOURCE QUOTE:
  % "Let $X$ be a finite type scheme over
  %  a field $k$, together with a line bundle $L$.
  %  Recall that if $F$ is a coherent sheaf on
  %  $X$ whose support is proper over $k$, then
  %  the \textbf{Hilbert polynomial} ${\Phi}\in \Q[\lambda]$ of $F$ is defined by
  %  the function
  %  $${\Phi}(m) = \chi(F(m)) =
  %  \sum_{i=0}^n (-1)^i \dim_k H^i(X, F\otimes L^{\otimes m})$$
  %  where the dimensions of the cohomologies are finite because of the
  %  coherence and properness conditions."
  \textit{Source: [Nitsure], \S 1 (FGA Explained Ch.~5); the section graded
  ring is the homogeneous coordinate ring of [Hartshorne], I.7.}
  Let \(s \in S\) and let \(L_s\) be a line bundle on the fibre \(X_s\), a scheme
  of finite type over the residue field \(\kappa(s)\). The \emph{section graded
  ring} of \(L_s\) is the graded commutative ring
  \[
    R(X_s, L_s) \;=\; \bigoplus_{m \ge 0} \Gamma\bigl(X_s,\, L_s^{\otimes m}\bigr),
  \]
  with degree-\(m\) component \(\Gamma(X_s, L_s^{\otimes m})\) and multiplication
  the tensor product of sections
  \(\Gamma(X_s, L_s^{\otimes m}) \otimes_{\kappa(s)}
  \Gamma(X_s, L_s^{\otimes m'}) \to \Gamma(X_s, L_s^{\otimes (m+m')})\),
  \(\sigma \otimes \tau \mapsto \sigma \cdot \tau\). Its degree-zero component is
  \(\Gamma(X_s, \mathcal{O}_{X_s})\), which contains \(\kappa(s)\). When \(X_s\) is
  proper over \(\kappa(s)\) and \(L_s\) is ample, \(R(X_s, L_s)\) is a finitely
  generated -- hence Noetherian -- graded \(\kappa(s)\)-algebra (a suitable
  Veronese power of \(L_s\) is very ample and embeds \(X_s\) into a projective
  space, exhibiting \(R(X_s, L_s)\) up to a finite extension as a quotient of a
  polynomial ring).
\end{definition}

\begin{definition}
[Section graded module]
  \label{def:sectionGradedModule}
  \lean{AlgebraicGeometry.sectionGradedModule}
  \uses{def:sectionGradedRing, lem:sectionGradedModule_gmodule}
  % SOURCE: [Nitsure], \S 1, sub-section "Stratification by Hilbert
  % Polynomials" (read from
  % references/nitsure-hilbert-quot-src/nitsure-hilbert-quot.tex, L453-L464).
  % SOURCE QUOTE:
  % "Recall that if $F$ is a coherent sheaf on
  %  $X$ whose support is proper over $k$, then
  %  the \textbf{Hilbert polynomial} ${\Phi}\in \Q[\lambda]$ of $F$ is defined by
  %  the function
  %  $${\Phi}(m) = \chi(F(m)) =
  %  \sum_{i=0}^n (-1)^i \dim_k H^i(X, F\otimes L^{\otimes m})$$
  %  where the dimensions of the cohomologies are finite because of the
  %  coherence and properness conditions."
  \textit{Source: [Nitsure], \S 1 (FGA Explained Ch.~5); cf.\ [Hartshorne],
  I.7.}
  Let \(\mathcal{F}_s\) be a coherent sheaf on the fibre \(X_s\) and \(L_s\) a
  line bundle on \(X_s\). The \emph{section graded module} of \(\mathcal{F}_s\)
  twisted by \(L_s\) is the graded \(R(X_s, L_s)\)-module
  (\cref{def:sectionGradedRing})
  \[
    M(X_s, \mathcal{F}_s, L_s)
    \;=\; \bigoplus_{m \ge 0} \Gamma\bigl(X_s,\, \mathcal{F}_s \otimes L_s^{\otimes m}\bigr),
  \]
  with degree-\(m\) component the \(\kappa(s)\)-vector space
  \(\Gamma(X_s, \mathcal{F}_s \otimes L_s^{\otimes m})\) and scalar action given
  by the section tensor product
  \(\Gamma(X_s, L_s^{\otimes m}) \otimes_{\kappa(s)}
  \Gamma(X_s, \mathcal{F}_s \otimes L_s^{\otimes m'})
  \to \Gamma(X_s, \mathcal{F}_s \otimes L_s^{\otimes (m+m')})\). Its graded Hilbert
  function is \(m \mapsto \dim_{\kappa(s)}
  \Gamma(X_s, \mathcal{F}_s \otimes L_s^{\otimes m})\).
\end{definition}

\begin{lemma}
[Finite generation of the section graded module]
  \label{lem:sectionGradedModule_fg}
  \lean{AlgebraicGeometry.sectionGradedModule_fg}
  \uses{def:sectionGradedModule, def:sectionGradedRing}
  % SOURCE: [Nitsure], \S 1, sub-section "Stratification by Hilbert
  % Polynomials" (read from
  % references/nitsure-hilbert-quot-src/nitsure-hilbert-quot.tex, L453-L464).
  % SOURCE QUOTE:
  % "Recall that if $F$ is a coherent sheaf on
  %  $X$ whose support is proper over $k$, then
  %  the \textbf{Hilbert polynomial} ${\Phi}\in \Q[\lambda]$ of $F$ is defined by
  %  the function
  %  $${\Phi}(m) = \chi(F(m)) =
  %  \sum_{i=0}^n (-1)^i \dim_k H^i(X, F\otimes L^{\otimes m})$$
  %  where the dimensions of the cohomologies are finite because of the
  %  coherence and properness conditions."
  \textit{Source: [Nitsure], \S 1 (FGA Explained Ch.~5); Serre's theorem on
  ample line bundles, [Hartshorne], I.7 / II.5.}
  Let \(X_s\) be proper over the field \(\kappa(s)\), let \(L_s\) be an ample line
  bundle on \(X_s\), and let \(\mathcal{F}_s\) be a coherent sheaf on \(X_s\).
  Then the section graded ring \(R(X_s, L_s)\) (\cref{def:sectionGradedRing}) is a
  Noetherian graded ring, and the section graded module
  \(M(X_s, \mathcal{F}_s, L_s)\) (\cref{def:sectionGradedModule}) is a finitely
  generated graded \(R(X_s, L_s)\)-module. In particular each component
  \(\Gamma(X_s, \mathcal{F}_s \otimes L_s^{\otimes m})\) is a finite-dimensional
  \(\kappa(s)\)-vector space.
\end{lemma}

\begin{proof}
  \uses{def:sectionGradedRing, def:sectionGradedModule}
  Since \(L_s\) is ample and \(X_s\) is proper over \(\kappa(s)\), a Veronese
  power \(L_s^{\otimes e}\) is very ample and gives a closed immersion
  \(X_s \hookrightarrow \mathbb{P}^N_{\kappa(s)}\). By Serre's theorem the section
  graded ring \(R(X_s, L_s)\) is a finitely generated graded
  \(\kappa(s)\)-algebra, so it is Noetherian. For the coherent sheaf
  \(\mathcal{F}_s\), the same ampleness yields, again by Serre, that
  \(\bigoplus_{m \ge 0} \Gamma(X_s, \mathcal{F}_s \otimes L_s^{\otimes m})\) is a
  finitely generated graded module over \(R(X_s, L_s)\): the twists
  \(\mathcal{F}_s \otimes L_s^{\otimes m}\) are globally generated for
  \(m \gg 0\), and the finitely many generating sections in low degrees together
  with the multiplication maps generate \(M(X_s, \mathcal{F}_s, L_s)\) over
  \(R(X_s, L_s)\). Finite dimensionality of each component over \(\kappa(s)\)
  follows from coherence and properness (finiteness of cohomology of coherent
  sheaves on a proper scheme over a field).
\end{proof}

\begin{lemma}[Existence and uniqueness of the Hilbert polynomial of a rational
  graded series]
  \label{lem:hilbertPoly_exists_mathlib}
  \lean{Polynomial.existsUnique_hilbertPoly}
  \mathlibok
  \textit{Provided by Mathlib
  (\texttt{Mathlib.RingTheory.Polynomial.HilbertPoly}).}
  Let \(F\) be a field of characteristic zero, let \(p \in F[X]\), and let
  \(d \in \mathbb{N}\). Then there is a unique polynomial \(h \in F[X]\) such that
  for some \(N \in \mathbb{N}\) and all \(n > N\) the coefficient of \(X^n\) in the
  power-series expansion of \(p \cdot (1 - X)^{-d} \in F\llbracket X\rrbracket\)
  equals \(h(n)\):
  \[
    \exists!\, h \in F[X],\ \exists N,\ \forall n > N : \;
    \bigl[ X^n \bigr]\bigl( p \cdot (1-X)^{-d} \bigr) \;=\; h(n).
  \]
  Equivalently, the eventual coefficient function of any rational power series
  with denominator a power of \(1 - X\) is, for \(n \gg 0\), given by a unique
  polynomial \(h\) (Mathlib's \(\mathrm{hilbertPoly}\,p\,d\)). The hypothesis
  \([\mathrm{CharZero}\,F]\) is satisfied by \(F = \mathbb{Q}\), which is the
  instance used to obtain \(\Phi_{\mathcal{F},s} \in \mathbb{Q}[\lambda]\).
\end{lemma}

\begin{lemma}[Additivity of \(\dim_\kappa\) on a short exact sequence
  (rank--nullity)]
  \label{lem:finrank_ses_additive_mathlib}
  \lean{Submodule.finrank_quotient_add_finrank}
  \mathlibok
  \textit{Provided by Mathlib
  (\texttt{Mathlib.LinearAlgebra.Dimension.RankNullity}).}
  Let \(\kappa\) be a field, \(V\) a finite-dimensional \(\kappa\)-vector space and
  \(W \subseteq V\) a subspace. Then
  \[
    \dim_\kappa (V / W) + \dim_\kappa W \;=\; \dim_\kappa V .
  \]
  Equivalently, \(\dim_\kappa = \operatorname{finrank}_\kappa\) is additive on a
  short exact sequence \(0 \to W \to V \to V/W \to 0\) of finite-dimensional
  \(\kappa\)-vector spaces. This is the additivity that converts the degreewise
  short exact sequences of graded components into additive relations among the
  coefficients of the Hilbert series.
\end{lemma}

\begin{lemma}[The rational power series \((1 - X)^{-d}\)]
  \label{lem:invOneSubPow_mathlib}
  \lean{PowerSeries.invOneSubPow}
  \mathlibok
  \textit{Provided by Mathlib
  (\texttt{Mathlib.RingTheory.PowerSeries.WellKnown}).}
  Let \(S\) be a commutative ring and \(d \in \mathbb{N}\). Then \((1 - X)^d\) is a
  unit in the power series ring \(S\llbracket X\rrbracket\), and Mathlib names its
  inverse
  \[
    \operatorname{invOneSubPow}(S, d) \;=\; (1 - X)^{-d}
      \;\in\; S\llbracket X\rrbracket^{\times}.
  \]
  Together with the coercion
  \(\operatorname{toPowerSeries} : F[X] \to F\llbracket X\rrbracket\) of a
  polynomial to its power series, this is the data in which a rational Hilbert
  series \(p \cdot (1 - X)^{-d}\) is expressed; it is the same rational-series shape
  whose eventual coefficients are extracted by
  \cref{lem:hilbertPoly_exists_mathlib}.
\end{lemma}

\begin{lemma}\leanok
[Graded Hilbert--Serre: rationality of the Hilbert series]
  \label{lem:gradedHilbertSerre_rational}
  \lean{AlgebraicGeometry.gradedModule_hilbertSeries_rational}
  \uses{lem:finrank_ses_additive_mathlib, lem:invOneSubPow_mathlib,
    def:ratHilb, lem:ratHilb_ofEventuallyZero, lem:ratHilb_ofDiffEq,
    lem:ratHilb_bump, def:graded_subquotientHilb,
    lem:graded_subquotient_isRatHilb}
  % SOURCE: [Stacks Project], "Commutative Algebra", \S "Noetherian graded
  % rings", Proposition \texttt{proposition-graded-hilbert-polynomial} (Tag
  % 00K1) (read from references/hilbert-serre-algebra.tex, L13893-L13905).
  % SOURCE QUOTE:
  % "Suppose that $S$ is a Noetherian graded ring
  %  and $M$ a finite graded $S$-module. Consider the
  %  function
  %  $$
  %  \mathbf{Z} \longrightarrow K'_0(S_0), \quad
  %  n \longmapsto [M_n]
  %  $$
  %  see Lemma \ref{lemma-graded-module-fg}.
  %  If $S_{+}$ is generated by elements of degree $1$,
  %  then this function is a numerical polynomial."
  \textit{Source: [Stacks Project], "Commutative Algebra", Tag 00K1
  (proposition on graded Hilbert polynomials); cf.\ [Nitsure], \S 1 (Snapper's
  Lemma), [Hartshorne], I.7.}
  Let \(R\) be a Noetherian graded ring whose degree-zero part is a field
  \(\kappa\) and which is generated as a \(\kappa\)-algebra by finitely many
  elements, and let \(M\) be a finitely generated graded \(R\)-module with
  finite-dimensional graded components \(M_n\). Then there exist
  \(p \in \mathbb{Q}[X]\) and \(d \in \mathbb{N}\) and an \(N \in \mathbb{N}\)
  such that for all \(n > N\)
  \[
    \dim_{\kappa} M_n \;=\; \bigl[ X^n \bigr]\bigl( p \cdot (1 - X)^{-d}\bigr);
  \]
  that is, the Hilbert series \(\sum_{n \ge 0} (\dim_\kappa M_n)\, X^n\) is, up to
  a polynomial correction in low degrees, the rational function
  \(p \cdot (1 - X)^{-d}\), expressed in the power-series form
  \(\operatorname{toPowerSeries}(p) \cdot \operatorname{invOneSubPow}(\mathbb{Q}, d)\)
  (\cref{lem:invOneSubPow_mathlib}).

  \emph{Status.} This is the substantive half of the graded Hilbert--Serre
  theorem (rationality of the Hilbert series of a finitely generated graded
  module). It is not covered by Mathlib: \cref{lem:hilbertPoly_exists_mathlib}
  supplies only the converse polynomial-extraction step from a rational series, not
  the production of the rational series from the module. The argument is realised by
  an \emph{ambient subquotient induction} (\cref{lem:graded_subquotient_isRatHilb}):
  the whole induction ranges over pairs of homogeneous submodules \(N' \le N\) of the
  single fixed graded module \(M\), and the Hilbert function is the ambient difference
  \(n \mapsto \dim_\kappa(N \cap M_n) - \dim_\kappa(N' \cap M_n)\). The class of such
  pairs is closed under the kernel and cokernel of multiplication by a degree-one
  element (each presented again as an ambient subquotient pair), so the inductive step
  never forms a graded quotient ring \(R/(x)\), a graded quotient module \(M/xM\), or a
  regrading of a derived carrier. The genuine Mathlib gap is thereby \emph{avoided}
  rather than filled; the residual obligation is the short ambient bookkeeping of the
  subquotient blocks (\cref{def:graded_subquotientHilb},
  \cref{lem:graded_subquotient_ker_coker},
  \cref{lem:graded_subquotient_degreewise_diff},
  \cref{lem:graded_subquotient_finite_transfer},
  \cref{lem:graded_subquotient_isRatHilb}), each of which manipulates only ambient
  families \(n \mapsto N_{\mathrm{aux}} \cap M_n\) of submodules of \(M\).
  % NOTE: The power-series rationality toolkit (\cref{subsec:isRatHilb}) is formalized
  %   and axiom-clean; the present theorem is a closed, axiom-clean Lean declaration
  %   (gradedModule_hilbertSeries_rational). The ambient subquotient route
  %   (\cref{lem:graded_subquotient_isRatHilb}) sidesteps missing graded-quotient API by
  %   working with ambient families of submodules of \(M\).
\end{lemma}

% SOURCE QUOTE PROOF: [Stacks Project], Tag 00K1, proof, inductive step (read
% from references/hilbert-serre-algebra.tex, L13938-L13947).
% "Let $\overline{M} = M/xM$. Note that the map $x : M \to M$
%  is {\it not} a map of graded $S$-modules, since it does
%  not map $M_d$ into $M_d$. Namely, for each $d$ we have the
%  following short exact sequence
%  $$
%  0 \to M_d \xrightarrow{x} M_{d + 1} \to \overline{M}_{d + 1} \to 0
%  $$
%  This proves that $[M_{d + 1}] - [M_d] = [\overline{M}_{d + 1}]$.
%  Hence we win by Lemma \ref{lemma-numerical-polynomial}."
\begin{proof}
  \uses{def:ratHilb, lem:ratHilb_ofEventuallyZero, lem:ratHilb_ofDiffEq,
    lem:ratHilb_bump, lem:graded_degreewise_finrank_diff,
    def:graded_subquotientHilb, lem:graded_subquotient_ker_coker,
    lem:graded_subquotient_degreewise_diff,
    lem:graded_subquotient_finite_transfer, lem:graded_subquotient_isRatHilb,
    lem:graded_homogeneousSubmodule_iSupIndep,
    lem:graded_homogeneousSubmodule_iSup_eq, lem:ideal_homogeneous_span_mathlib,
    lem:isHomogeneousElem_graded_smul_mathlib,
    lem:fg_restrictScalars_of_surjective_mathlib}
    \leanok
  We argue with the abstract graded \(\kappa\)-algebra
  \(R = \bigoplus_{n \ge 0} R_n\), whose degree-zero part \(R_0 = \kappa\) is a
  field, and the finitely generated graded \(R\)-module
  \(M = \bigoplus_{n \ge 0} M_n\) with finite-dimensional components \(M_n\);
  there is no appeal to any geometric realisation of \(R\) or \(M\).

  \emph{Reduction to generation in degree \(1\).} Since \(R\) is Noetherian and
  finitely generated as a \(\kappa\)-algebra, it is generated by finitely many
  homogeneous elements of positive degree. Replacing \(R\) by a Veronese
  subalgebra \(R^{(e)} = \bigoplus_n R_{en}\) for a suitable \(e\) and re-grading
  (and \(M\) by the corresponding Veronese module), we may assume \(R\) is
  generated by its degree-one part, \(\kappa[R_1] = R\). The degree-one part
  \(R_1\) is then a finite-dimensional \(\kappa\)-vector space whose \(\kappa\)-basis
  \(x_1, \dots, x_r\) generates \(R\) as a \(\kappa\)-algebra; \(r = \dim_\kappa R_1\)
  is the number on which we induct.
  % NOTE: The existence of finitely many degree-one generators (the number r the
  % induction is on) is expected to follow from the standing hypotheses --
  % "finitely generated κ-algebra" together with "generated in degree 1" -- so no
  % new hypothesis on gradedModule_hilbertSeries_rational is anticipated. To be
  % confirmed when the prover formalizes the assembly: degree-1 generation gives
  % R₁ as a finite-dimensional κ-vector space spanning R as a κ-algebra, and a
  % κ-basis of R₁ is a finite degree-one generating set.

  \emph{Set-up of the ambient data.} Regard \(M\) as the fixed ambient graded
  \(\kappa\)-module with internal decomposition \(n \mapsto \mathcal{M}_n := M_n\)
  (\cref{def:sectionGradedModule}); the components are finite-dimensional. For each
  basis vector \(x_i \in R_1\) let \(t_i : M \to M\) be multiplication by \(x_i\), a
  \(\kappa\)-linear endomorphism that raises degree by one (it sends \(M_n\) into
  \(M_{n+1}\), \cref{lem:isHomogeneousElem_graded_smul_mathlib}). Since \(R\) is
  commutative the \(t_i\) pairwise commute, and they trivially preserve the pair of
  homogeneous submodules \((N, N') = (\top, \bot)\). Because \(R\) is generated by
  \(x_1, \dots, x_r\) and \(M\) is finitely generated over \(R\), \(M\) is finitely
  generated as a module over the polynomial ring \(\kappa[x_1, \dots, x_r]\) acting
  through the \(t_i\); this is the finiteness condition required below.

  \emph{Invoking the ambient subquotient induction.} Apply
  \cref{lem:graded_subquotient_isRatHilb} to the pair \((N, N') = (\top, \bot)\)
  with the \(r\) commuting degree-\((+1)\) endomorphisms \(t_1, \dots, t_r\). It
  yields \(\mathrm{IsRatHilb}\) of order \(r\) for the ambient subquotient Hilbert
  function of this pair (\cref{def:graded_subquotientHilb}),
  \[
    h(n) \;=\; \dim_\kappa(\top \cap \mathcal{M}_n) - \dim_\kappa(\bot \cap \mathcal{M}_n)
      \;=\; \dim_\kappa M_n - 0 \;=\; \dim_\kappa M_n .
  \]
  Hence \(\mathrm{IsRatHilb}(n \mapsto \dim_\kappa M_n,\, r)\) (\cref{def:ratHilb}):
  there are \(p \in \mathbb{Q}[X]\), \(d := r \in \mathbb{N}\) and \(N \in \mathbb{N}\)
  with
  \[
    \dim_\kappa M_n \;=\; \bigl[X^n\bigr]\bigl(p \cdot (1 - X)^{-d}\bigr)
    \qquad (n > N),
  \]
  which is the claim, with pole order \(d = r\). The internal mechanics of
  \cref{lem:graded_subquotient_isRatHilb} -- the base case via
  \cref{lem:ratHilb_ofEventuallyZero}, and the inductive step assembling the
  kernel/cokernel closure (\cref{lem:graded_subquotient_ker_coker}), the degreewise
  difference identity (\cref{lem:graded_subquotient_degreewise_diff}, built on
  \(D_5\), \cref{lem:graded_degreewise_finrank_diff}) and the finiteness transfer
  (\cref{lem:graded_subquotient_finite_transfer}) into
  \cref{lem:ratHilb_ofDiffEq} and \cref{lem:ratHilb_bump} -- are the Stacks~00K1
  induction realised entirely with ambient homogeneous submodules of \(M\).
\end{proof}

\subsection{The rationality toolkit \texorpdfstring{\(\mathrm{IsRatHilb}\)}{IsRatHilb}}
\label{subsec:isRatHilb}
% NOTE: Declarations in this subsection are currently \texttt{private} in the Lean
% source; \lean{} pins may not resolve until they are moved to a dedicated module.

The proof of \cref{lem:gradedHilbertSerre_rational} runs the Stacks~00K1 induction
on the number of degree-one generators. Each inductive step manipulates the
\emph{shape} of the Hilbert series -- a numerator polynomial over a power of
\((1-X)\) -- rather than the module itself. We isolate that bookkeeping in a
single predicate \(\mathrm{IsRatHilb}\) on functions \(f : \mathbb{N} \to
\mathbb{Q}\), together with the closure lemmas the induction consumes. These are
project-internal helpers realising standard steps of the Stacks~00K1 argument at
the level of power series; Mathlib supplies only the converse extraction
(\cref{lem:hilbertPoly_exists_mathlib}).

\begin{lemma}\leanok
[Coefficients of \((1-X)^{-1}\) times a series are partial sums]
  \label{lem:coeff_invOneSubPow_one_mul}
  \lean{AlgebraicGeometry.coeff_invOneSubPow_one_mul}
  \uses{lem:invOneSubPow_mathlib}
  Let \(F \in \mathbb{Q}\llbracket X\rrbracket\) be a power series. Then for every
  \(n \in \mathbb{N}\) the \(n\)-th coefficient of
  \(\operatorname{invOneSubPow}(\mathbb{Q}, 1)\cdot F = (1-X)^{-1}\cdot F\)
  (\cref{lem:invOneSubPow_mathlib}) is the partial sum of the coefficients of
  \(F\):
  \[
    \bigl[X^n\bigr]\bigl((1-X)^{-1}\cdot F\bigr)
    \;=\; \sum_{k=0}^{n} \bigl[X^k\bigr] F .
  \]
  Equivalently, multiplication by \((1-X)^{-1} = \sum_{m\ge 0} X^m\) is the
  summation operator on coefficient sequences. Proved directly in Lean from the
  identity \(\operatorname{invOneSubPow}(\mathbb{Q},1) = \sum_{m\ge 0} X^m\) and
  the Cauchy product.
\end{lemma}

\begin{lemma}\leanok
[Antidifference of a rational Hilbert series]
  \label{lem:rationalHilbert_antidiff}
  \lean{AlgebraicGeometry.rationalHilbert_antidiff}
  \uses{lem:coeff_invOneSubPow_one_mul, lem:invOneSubPow_mathlib}
  Let \(H, \delta : \mathbb{N} \to \mathbb{Q}\), let \(q \in \mathbb{Q}[X]\) and
  \(e, N \in \mathbb{N}\). Suppose that for all \(n > N\) the sequence \(\delta\)
  is the coefficient sequence of the rational series \(q\cdot(1-X)^{-e}\),
  \[
    \delta(n) \;=\; \bigl[X^n\bigr]\bigl(q\cdot(1-X)^{-e}\bigr),
  \]
  and that \(\delta\) is the first difference of \(H\) beyond \(N\), i.e.
  \(H(n+1) - H(n) = \delta(n+1)\) for all \(n > N\). Then there is a numerator
  polynomial \(p \in \mathbb{Q}[X]\) such that for all \(n > N\)
  \[
    H(n) \;=\; \bigl[X^n\bigr]\bigl(p\cdot(1-X)^{-(e+1)}\bigr).
  \]
  That is, summing a coefficient sequence (anti-differencing) raises the pole
  order by one: the partial-sum series of \(q\cdot(1-X)^{-e}\) equals
  \(p\cdot(1-X)^{-(e+1)}\) up to a polynomial correction in low degrees. This is
  the power-series heart of the Stacks~00K1 inductive step.
\end{lemma}

\begin{proof}
  \uses{lem:coeff_invOneSubPow_one_mul, lem:invOneSubPow_mathlib}
  \leanok
  Write \(F = q\cdot(1-X)^{-e}\). Multiplying by
  \((1-X)^{-1}\) and using the factorisation
  \((1-X)^{-(e+1)} = (1-X)^{-1}\cdot(1-X)^{-e}\)
  (\cref{lem:invOneSubPow_mathlib}), \cref{lem:coeff_invOneSubPow_one_mul} gives
  \(\bigl[X^m\bigr]\bigl(q\cdot(1-X)^{-(e+1)}\bigr) = \sum_{k=0}^{m}[X^k]F\), the
  partial sums of \(F\). On the other hand, telescoping the first-difference
  hypothesis from \(N+1\) expresses \(H(N+1+j)\) as \(H(N+1)\) plus a window sum
  of the coefficients of \(F\). Splitting the partial sum
  \(\sum_{k=0}^{n}[X^k]F\) at \(N+2\) and absorbing the resulting constant
  discrepancy into a numerator \(C\,(1-X)^{e}\) (a constant function is the
  order-\((e+1)\) coefficient sequence of \(C\,(1-X)^{e}\cdot(1-X)^{-(e+1)} =
  C\,(1-X)^{-1}\)), the explicit numerator
  \(p = C\,(1-X)^{e} + q\) realises \(H(n)\) as
  \([X^n]\bigl(p\cdot(1-X)^{-(e+1)}\bigr)\) for all \(n > N\).
\end{proof}

\begin{definition}\leanok
[Rational Hilbert function]
  \label{def:ratHilb}
  \lean{AlgebraicGeometry.IsRatHilb}
  \uses{lem:invOneSubPow_mathlib}
  \textit{Source: Stacks 00K1.}
  For a function \(f : \mathbb{N} \to \mathbb{Q}\) and \(d \in \mathbb{N}\), say
  that \(f\) is a \emph{rational Hilbert function of order \(d\)}, written
  \(\mathrm{IsRatHilb}(f, d)\), when there exist a numerator polynomial
  \(p \in \mathbb{Q}[X]\) and an \(N \in \mathbb{N}\) such that for all \(n > N\)
  \[
    f(n) \;=\; \bigl[X^n\bigr]\bigl(p\cdot(1-X)^{-d}\bigr)
    \;=\; \bigl[X^n\bigr]\bigl(\operatorname{toPowerSeries}(p)\cdot
      \operatorname{invOneSubPow}(\mathbb{Q}, d)\bigr)
  \]
  (\cref{lem:invOneSubPow_mathlib}). This is the ``eventually equal to the
  coefficient sequence of a rational series \(p\cdot(1-X)^{-d}\)'' predicate that
  the graded Hilbert--Serre induction tracks; its conclusion shape is exactly the
  one consumed by \cref{lem:hilbertPoly_exists_mathlib}.
\end{definition}

\begin{lemma}\leanok
[Base case: eventually-zero functions are rational of order \(0\)]
  \label{lem:ratHilb_ofEventuallyZero}
  \lean{AlgebraicGeometry.IsRatHilb.ofEventuallyZero}
  \uses{def:ratHilb}
  If \(f : \mathbb{N} \to \mathbb{Q}\) is eventually zero -- there is
  \(N\) with \(f(n) = 0\) for all \(n > N\) -- then \(\mathrm{IsRatHilb}(f, 0)\)
  (\cref{def:ratHilb}). Indeed \(f\) is then the coefficient sequence of the
  numerator polynomial \(0\) over \((1-X)^{0} = 1\). This is the base case of the
  Stacks~00K1 induction (no degree-one generators). Proved directly in Lean.
\end{lemma}

\begin{lemma}\leanok
[Raising the pole order]
  \label{lem:ratHilb_bump}
  \lean{AlgebraicGeometry.IsRatHilb.bump}
  \uses{def:ratHilb, lem:invOneSubPow_mathlib}
  If \(\mathrm{IsRatHilb}(f, d)\) (\cref{def:ratHilb}) then
  \(\mathrm{IsRatHilb}(f, d+1)\): multiplying the numerator by \((1-X)\) cancels
  one factor of \((1-X)^{-1}\), so the same coefficient sequence is realised at
  pole order \(d+1\). This lets two rational Hilbert functions of different
  orders be brought to a common denominator. Proved directly in Lean using
  \((1-X)\cdot(1-X)^{-(d+1)} = (1-X)^{-d}\) (\cref{lem:invOneSubPow_mathlib}).
\end{lemma}

\begin{lemma}\leanok
[Closure under difference]
  \label{lem:ratHilb_sub}
  \lean{AlgebraicGeometry.IsRatHilb.sub}
  \uses{def:ratHilb}
  If \(\mathrm{IsRatHilb}(f, d)\) and \(\mathrm{IsRatHilb}(g, d)\)
  (\cref{def:ratHilb}) at the same order \(d\), then
  \(\mathrm{IsRatHilb}(n \mapsto f(n) - g(n),\, d)\): the numerator of the
  difference is the difference of the numerators (over the common denominator
  \((1-X)^{-d}\)). Proved directly in Lean.
\end{lemma}

\begin{lemma}\leanok
[Closure under right shift]
  \label{lem:ratHilb_shiftRight}
  \lean{AlgebraicGeometry.IsRatHilb.shiftRight}
  \uses{def:ratHilb}
  If \(\mathrm{IsRatHilb}(f, d)\) (\cref{def:ratHilb}) then
  \(\mathrm{IsRatHilb}(n \mapsto f(n-1),\, d)\): the right shift on coefficient
  sequences is multiplication of the series by \(X\), so the numerator
  \(p\) is replaced by \(X\cdot p\) at the same pole order. Proved directly in
  Lean.
\end{lemma}

\begin{lemma}\leanok
[Antidifference step for the predicate]
  \label{lem:ratHilb_antidiff}
  \lean{AlgebraicGeometry.IsRatHilb.antidiff}
  \uses{def:ratHilb, lem:rationalHilbert_antidiff}
  Let \(H, g : \mathbb{N} \to \mathbb{Q}\) and \(e, N \in \mathbb{N}\). If
  \(\mathrm{IsRatHilb}(g, e)\) (\cref{def:ratHilb}) and \(g\) is the first
  difference of \(H\) beyond \(N\) (\(H(n+1) - H(n) = g(n+1)\) for all
  \(n > N\)), then \(\mathrm{IsRatHilb}(H, e+1)\). This is the predicate-level
  packaging of \cref{lem:rationalHilbert_antidiff}: a finite difference raises
  the pole order by one, the inductive heart of graded Hilbert--Serre.
\end{lemma}

\begin{proof}
  \uses{def:ratHilb, lem:rationalHilbert_antidiff}
  \leanok
  Unfolding \(\mathrm{IsRatHilb}(g, e)\) supplies a numerator \(q\) and a
  threshold \(N_g\) with \(g(n) = [X^n](q\cdot(1-X)^{-e})\) for \(n > N_g\).
  Applying \cref{lem:rationalHilbert_antidiff} with this \(q\) and threshold
  \(\max(N, N_g)\) -- whose two hypotheses are exactly the rational form of \(g\)
  and the first-difference identity for \(H\) -- yields a numerator \(p\) with
  \(H(n) = [X^n](p\cdot(1-X)^{-(e+1)})\) for \(n > \max(N, N_g)\), i.e.
  \(\mathrm{IsRatHilb}(H, e+1)\).
\end{proof}

\begin{lemma}\leanok
[Inductive-step engine from a degreewise short exact sequence]
  \label{lem:ratHilb_ofDiffEq}
  \lean{AlgebraicGeometry.IsRatHilb.ofDiffEq}
  \uses{def:ratHilb, lem:ratHilb_antidiff, lem:ratHilb_sub, lem:ratHilb_shiftRight}
  Let \(h_M, h_C, h_K : \mathbb{N} \to \mathbb{Q}\) and \(d, N \in \mathbb{N}\).
  Suppose \(\mathrm{IsRatHilb}(h_C, d)\) and \(\mathrm{IsRatHilb}(h_K, d)\)
  (\cref{def:ratHilb}), and that for all \(n > N\)
  \[
    h_M(n+1) - h_M(n) \;=\; h_C(n+1) - h_K(n).
  \]
  Then \(\mathrm{IsRatHilb}(h_M, d+1)\). This packages the entire power-series
  side of the Stacks~00K1 inductive step: \(h_M\) is the Hilbert function of
  \(M\), and \(h_C, h_K\) are the Hilbert functions of the cokernel
  \(C = M/xM\) and kernel \(K = \ker(x : M \to M(1))\) of multiplication by a
  degree-one element \(x\); the displayed identity is the alternating-sum
  consequence of the degreewise short exact sequence
  \(0 \to K_n \to M_n \xrightarrow{x} M_{n+1} \to C_{n+1} \to 0\). Given the
  inductive hypothesis that \(h_C, h_K\) are rational of order \(d\), the
  conclusion is rationality of \(h_M\) at order \(d+1\).
\end{lemma}

\begin{proof}
  \uses{lem:ratHilb_antidiff, lem:ratHilb_sub, lem:ratHilb_shiftRight}
  \leanok
  Set \(g(n) = h_C(n) - h_K(n-1)\). By closure under right shift
  (\cref{lem:ratHilb_shiftRight}) \(n \mapsto h_K(n-1)\) is rational of order
  \(d\), and by closure under difference (\cref{lem:ratHilb_sub}) so is \(g\). The
  hypothesis rewrites as \(h_M(n+1) - h_M(n) = g(n+1)\) for \(n > N\) (since
  \(g(n+1) = h_C(n+1) - h_K(n)\)). Applying the antidifference step
  (\cref{lem:ratHilb_antidiff}) to \(H = h_M\) and this \(g\) gives
  \(\mathrm{IsRatHilb}(h_M, d+1)\).
\end{proof}

\begin{theorem}[Hilbert polynomial from the section module]
  \label{thm:hilbertPoly_of_sectionModule}
  \lean{AlgebraicGeometry.Scheme.hilbertPolynomialOfSectionModule}
  \uses{lem:sectionGradedModule_fg, lem:gradedHilbertSerre_rational,
    lem:hilbertPoly_exists_mathlib}
  % SOURCE: [Nitsure], \S 1, sub-section "Stratification by Hilbert
  % Polynomials" (read from
  % references/nitsure-hilbert-quot-src/nitsure-hilbert-quot.tex, L468-L478).
  % SOURCE QUOTE:
  % "Let $X\to S$ be a finite type morphism of noetherian schemes,
  %  and let $L$ be a line bundle on $X$. Let $\F$
  %  be any coherent sheaf on
  %  $X$ whose schematic support is proper over $S$. Then for each $s\in S$,
  %  we get a polynomial ${\Phi}_s \in  \Q[\lambda]$
  %  which is the Hilbert polynomial of
  %  the restriction $\F_s = F|_{X_s}$ of $\F$ to the fiber $X_s$ over $s$,
  %  calculated with respect to the line bundle $L_s = L|_{X_s}$."
  \textit{Source: [Nitsure], \S 1 (FGA Explained Ch.~5).}
  Let \(X_s\) be proper over the field \(\kappa(s)\), let \(L_s\) be an ample line
  bundle on \(X_s\), and let \(\mathcal{F}_s\) be a coherent sheaf on \(X_s\).
  Then there is a unique polynomial \(\Phi_{\mathcal{F},s} \in \mathbb{Q}[\lambda]\)
  such that for all \(m \gg 0\)
  \[
    \Phi_{\mathcal{F},s}(m)
    \;=\; \dim_{\kappa(s)} \Gamma\bigl(X_s,\, \mathcal{F}_s \otimes L_s^{\otimes m}\bigr).
  \]
  This \(\Phi_{\mathcal{F},s}\) is the Hilbert polynomial of
  \cref{def:hilbert_polynomial}.
\end{theorem}

\begin{proof}
  \uses{lem:sectionGradedModule_fg, lem:gradedHilbertSerre_rational,
    lem:hilbertPoly_exists_mathlib}
  By \cref{lem:sectionGradedModule_fg} the section graded module
  \(M(X_s, \mathcal{F}_s, L_s)\) is a finitely generated graded module over the
  Noetherian graded ring \(R(X_s, L_s)\), with finite-dimensional components and
  degree-zero part the field \(\kappa(s)\). Applying the abstract graded
  Hilbert--Serre rationality lemma (\cref{lem:gradedHilbertSerre_rational}) to the
  graded \(\kappa(s)\)-algebra \(R = R(X_s, L_s)\) and the finitely generated
  graded module \(M = M(X_s, \mathcal{F}_s, L_s)\) exhibits the graded Hilbert
  function
  \(m \mapsto \dim_{\kappa(s)} \Gamma(X_s, \mathcal{F}_s \otimes L_s^{\otimes m})\)
  as the eventual coefficient function of a rational series: there are
  \(p \in \mathbb{Q}[X]\) and \(d \in \mathbb{N}\) and an \(N\) with
  \(\dim_{\kappa(s)} \Gamma(X_s, \mathcal{F}_s \otimes L_s^{\otimes m})
  = [X^m]\,(p \cdot (1 - X)^{-d})\) for \(m > N\). Apply
  \cref{lem:hilbertPoly_exists_mathlib} over the characteristic-zero field
  \(F = \mathbb{Q}\) to this pair \((p, d)\): it yields a unique polynomial
  \(h \in \mathbb{Q}[X]\) whose values \(h(n)\) eventually equal
  \([X^n]\,(p \cdot (1 - X)^{-d})\). Setting
  \(\Phi_{\mathcal{F},s} := h\) (with the variable renamed \(\lambda\)) gives a
  polynomial with \(\Phi_{\mathcal{F},s}(m) = \dim_{\kappa(s)}
  \Gamma(X_s, \mathcal{F}_s \otimes L_s^{\otimes m})\) for \(m \gg 0\). Uniqueness
  is the \(\exists!\) clause of \cref{lem:hilbertPoly_exists_mathlib}: any two
  polynomials agreeing with the Hilbert function for \(m \gg 0\) agree at
  infinitely many points, hence are equal.
\end{proof}

\subsection{The ambient subquotient induction for the Stacks~00K1 step}
\label{subsec:gradedModuleApi}

The power-series side of \cref{lem:gradedHilbertSerre_rational} is fully isolated
in the \(\mathrm{IsRatHilb}\) toolkit (\cref{subsec:isRatHilb}). What remains is
the \emph{graded-module} side of the Stacks~00K1 induction. Rather than build the
kernel and cokernel of multiplication by a degree-one element as graded modules
over a smaller quotient ring \(R/(x)\) -- which would force a graded structure on a
derived carrier -- we range the induction over \emph{pairs of homogeneous
submodules of a single fixed ambient graded module} \(M\). A pair \((N, N')\) with
\(N' \le N\) homogeneous submodules of \(M\) stands for the subquotient \(N/N'\); its
ambient Hilbert function is the difference \(n \mapsto \dim_\kappa(N \cap M_n) -
\dim_\kappa(N' \cap M_n)\). The class of such pairs, equipped with finitely many
pairwise-commuting degree-one endomorphisms preserving \(N\) and \(N'\), is
\emph{closed} under the kernel and cokernel of one of those endomorphisms, both
again presented as ambient subquotient pairs. Consequently no graded quotient ring,
graded quotient module, or regrading of a subtype/quotient carrier is ever formed:
every object that appears is an ambient family \(n \mapsto N_{\mathrm{aux}} \cap M_n\)
of submodules of the fixed \(M\).

This subsection records that machinery as a thin layer over existing Mathlib
scaffold: Mathlib already provides the notion of a homogeneous submodule
(\cref{lem:submodule_isHomogeneous_mathlib}), the homogeneity of a span of
homogeneous elements (\cref{lem:ideal_homogeneous_span_mathlib}), the degree-shift
of multiplication by a homogeneous element
(\cref{lem:isHomogeneousElem_graded_smul_mathlib}), the rank--nullity identity
(\cref{lem:finrank_range_add_finrank_ker_mathlib}), and the finite-generation
transfer along a surjection (\cref{lem:fg_restrictScalars_of_surjective_mathlib}).
What is supplied here is the internal direct-sum decomposition of a homogeneous
submodule as an ambient family (the split \(G_1\)), the subquotient Hilbert data,
the kernel/cokernel closure, the degreewise difference identity, the ambient
finiteness transfer, and the induction itself.

\emph{Setting for this subsection.} Let \(R = \bigoplus_{n \ge 0} R_n\) be a graded
ring whose degree-zero part \(R_0 = \kappa\) is a field and which is generated as a
\(\kappa\)-algebra in degree one, and let \(M = \bigoplus_{n \ge 0} M_n\) be a graded
\(\kappa\)-vector space (the fixed ambient module) with finite-dimensional
components \(M_n\). All submodules \(N, N', \dots\) below are homogeneous submodules
of \(M\), and the relevant endomorphisms are \(\kappa\)-linear maps \(M \to M\)
raising degree by one -- typically multiplication by a homogeneous element
\(x \in R_1\).

\subsubsection*{Mathlib dependency anchors}

\begin{lemma}[Homogeneous submodule]
  \label{lem:submodule_isHomogeneous_mathlib}
  \lean{Submodule.IsHomogeneous}
  \mathlibok
  \textit{Provided by Mathlib
  (\texttt{Mathlib.RingTheory.GradedAlgebra.Homogeneous.Submodule}).}
  Let \(R = \bigoplus_n R_n\) be a graded ring and \(M = \bigoplus_n M_n\) an
  internally graded \(R\)-module. A submodule \(p \subseteq M\) is
  \emph{homogeneous} when it is closed under taking homogeneous components:
  \(m \in p\) implies each graded piece \(\mathrm{proj}_n(m) \in p\). Equivalently
  \(p = \bigoplus_n (p \cap M_n)\) as a set. Mathlib bundles this into a homogeneous
  submodule type providing order and extensionality structure; the ambient internal
  direct-sum decomposition \(p = \bigoplus_n (M_n \cap p)\) is supplied by the split
  \(G_1\) (\cref{lem:graded_homogeneousSubmodule_iSupIndep},
  \cref{lem:graded_homogeneousSubmodule_iSup_eq}).
\end{lemma}

\begin{lemma}[A span of homogeneous elements is homogeneous]
  \label{lem:ideal_homogeneous_span_mathlib}
  \lean{Ideal.homogeneous_span}
  \mathlibok
  \textit{Provided by Mathlib
  (\texttt{Mathlib.RingTheory.GradedAlgebra.Homogeneous.Ideal}).}
  Let \(R = \bigoplus_n R_n\) be a graded ring and \(s \subseteq R\) a set each of
  whose elements is homogeneous. Then the ideal \(\mathrm{span}\,s\) is a
  homogeneous ideal. In particular, for \(x \in R_1\) the principal ideal \((x)\)
  is homogeneous, hence \(xM\) and the auxiliary submodule \(N' + x \cdot N\) are
  homogeneous; this homogeneity feeds the kernel/cokernel closure
  (\cref{lem:graded_subquotient_ker_coker}).
\end{lemma}

\begin{lemma}[Rank--nullity]
  \label{lem:finrank_range_add_finrank_ker_mathlib}
  \lean{LinearMap.finrank_range_add_finrank_ker}
  \mathlibok
  \textit{Provided by Mathlib
  (\texttt{Mathlib.LinearAlgebra.FiniteDimensional.Lemmas}).}
  Let \(\kappa\) be a field, \(V\) a finite-dimensional \(\kappa\)-vector space and
  \(\varphi : V \to W\) a \(\kappa\)-linear map. Then
  \[
    \dim_\kappa(\operatorname{range}\varphi) + \dim_\kappa(\ker\varphi)
      \;=\; \dim_\kappa V .
  \]
  This is the rank--nullity identity underlying the degreewise difference identity
  \(D_5\) (\cref{lem:graded_degreewise_finrank_diff}).
\end{lemma}

\begin{lemma}[Finite generation transfers along a surjection]
  \label{lem:fg_restrictScalars_of_surjective_mathlib}
  \lean{Submodule.FG.restrictScalars_of_surjective}
  \mathlibok
  \textit{Provided by Mathlib (\texttt{Mathlib.RingTheory.Finiteness.Basic}).}
  Let \(\varphi : R \to A\) be a surjective ring homomorphism and \(M\) an
  \(A\)-module that is finitely generated when viewed as an \(R\)-module via
  \(\varphi\). Then \(M\) is finitely generated as an \(A\)-module. Applied to the
  surjection \(\kappa[t_0, \dots, t_{r-1}] \twoheadrightarrow \kappa[t_0, \dots,
  t_{r-2}]\) killing the endomorphism \(x = t_{r-1}\) that annihilates the two
  subquotients, this transfers their finite generation down one generator
  (\cref{lem:graded_subquotient_finite_transfer}).
\end{lemma}

\begin{lemma}[Multiplication by a homogeneous element shifts degree]
  \label{lem:isHomogeneousElem_graded_smul_mathlib}
  \lean{SetLike.IsHomogeneousElem.graded_smul}
  \mathlibok
  \textit{Provided by Mathlib (\texttt{Mathlib.Algebra.GradedMulAction}).}
  Let \(\mathcal{A}\) be a graded family acting on a graded family \(\mathcal{M}\)
  via a graded scalar multiplication. If \(a \in R\) is homogeneous of degree
  \(i\) -- in particular \(x \in R_1\) -- then for \(m \in M_n\) one has
  \(a \cdot m \in M_{i+n}\). This is the degree-shifting property that makes \(xM\)
  a homogeneous submodule (with degree-\((n+1)\) part \(x \cdot M_n\)) and that
  underlies the ambient identity \(x \cdot N \cap M_{n+1} = x \cdot (N \cap M_n)\)
  used in the subquotient closure (\cref{lem:graded_subquotient_ker_coker}) and
  difference identity (\cref{lem:graded_subquotient_degreewise_diff}).
\end{lemma}

\begin{lemma}[Commutative subalgebra generated by commuting elements]
  \label{lem:adjoinCommRingOfComm_mathlib}
  \lean{Algebra.isMulCommutative_adjoin}
  \mathlibok
  \textit{Provided by Mathlib
  (\texttt{Mathlib.Algebra.Algebra.Subalgebra.Lattice}).}
  Let \(R\) be a commutative semiring and \(A\) an \(R\)-algebra, possibly
  noncommutative. If \(s \subseteq A\) is a set of pairwise-commuting elements
  (\(a b = b a\) for all \(a, b \in s\)), then the \(R\)-subalgebra
  \(\mathrm{adjoin}_R(s)\) generated by \(s\) has commutative multiplication; the
  commutative-ring structure on it follows. Applied to the \(r\) pairwise-commuting degree-raising endomorphisms
  \(x_0, \dots, x_{r-1} \in \operatorname{End}_\kappa(M)\), it yields a
  \emph{commutative} \(\kappa\)-subalgebra \(A = \mathrm{adjoin}_\kappa\{x_0, \dots,
  x_{r-1}\} \subseteq \operatorname{End}_\kappa(M)\); the polynomial action of
  \(\kappa[t_0, \dots, t_{r-1}]\) on \(M\) factors through this commutative \(A\)
  (\cref{lem:graded_subquotient_finite_transfer}).
\end{lemma}

\begin{lemma}[Finiteness transfers along a surjective linear map]
  \label{lem:module_finite_of_surjective_mathlib}
  \lean{Module.Finite.of_surjective}
  \mathlibok
  \textit{Provided by Mathlib (\texttt{Mathlib.RingTheory.Finiteness.Defs}).}
  Let \(A\) be a ring, \(M, N\) two \(A\)-modules and \(\varphi : M \to N\) a
  surjective \(A\)-linear map. If \(M\) is a finite \(A\)-module, then so is \(N\).
  This is the finiteness-transfer step at the heart of the ambient finiteness
  transfer (\cref{lem:graded_subquotient_finite_transfer}) and of the
  numerator/denominator finiteness lemmas
  (\cref{lem:graded_polyQuot_finite_of_le_numerator},
  \cref{lem:graded_polyQuot_finite_of_le_denominator}).
\end{lemma}

\begin{lemma}[Finiteness descends along an injective linear map into a Noetherian module]
  \label{lem:module_finite_of_injective_mathlib}
  \lean{Module.Finite.of_injective}
  \mathlibok
  \textit{Provided by Mathlib (\texttt{Mathlib.RingTheory.Noetherian.Defs}).}
  Let \(A\) be a ring, \(M, N\) two \(A\)-modules with \(N\) Noetherian (in
  particular, finite over a Noetherian ring \(A\)), and \(\varphi : M \to N\) an
  injective \(A\)-linear map. Then \(M\) is a finite \(A\)-module. This is the
  submodule-finiteness step in \cref{lem:graded_polyQuot_finite_of_le_numerator}.
\end{lemma}

\begin{lemma}[The lift of a linear map through a quotient]
  \label{lem:submodule_liftQ_mathlib}
  \lean{Submodule.liftQ}
  \mathlibok
  \textit{Provided by Mathlib (\texttt{Mathlib.LinearAlgebra.Quotient.Basic}).}
  Let \(A\) be a ring, \(M, N\) two \(A\)-modules, \(P \subseteq M\) a submodule and
  \(\varphi : M \to N\) an \(A\)-linear map with \(P \subseteq \ker\varphi\). Then
  \(\varphi\) factors uniquely through a linear map \(M/P \to N\) on the quotient.
  This realises the descent of a (\(\sigma\)-semilinear) map to the subquotient
  quotient in the ambient finiteness transfer
  (\cref{lem:graded_subquotient_finite_transfer}) and the denominator lemma
  (\cref{lem:graded_polyQuot_finite_of_le_denominator}).
\end{lemma}

\begin{lemma}[Only finitely many members of an independent family are nonzero]
  \label{lem:finite_ne_bot_of_iSupIndep_mathlib}
  \lean{Submodule.finite_ne_bot_of_iSupIndep}
  \mathlibok
  \textit{Provided by Mathlib
  (\texttt{Mathlib.LinearAlgebra.Dimension.Finite}).}
  Let \(\kappa\) be a field (more generally a division ring) and \(V\) a
  finite-dimensional \(\kappa\)-vector space. If \((p_i)_{i \in \iota}\) is an
  \(\mathrm{iSupIndep}\) family of subspaces of \(V\), then the set
  \(\{\,i : p_i \ne \bot\,\}\) is finite. This is the finiteness input that turns
  degreewise independence into eventual vanishing in the base case
  (\cref{lem:graded_subquotient_base_eventuallyZero}).
\end{lemma}

\begin{lemma}[The polynomial ring in finitely many variables is Noetherian]
  \label{lem:mvPolynomial_isNoetherianRing_fin_mathlib}
  \lean{MvPolynomial.isNoetherianRing_fin}
  \mathlibok
  \textit{Provided by Mathlib
  (\texttt{Mathlib.RingTheory.Polynomial.Basic}).}
  Let \(R\) be a Noetherian commutative ring and \(r \in \mathbb{N}\). Then the
  multivariate polynomial ring \(\operatorname{MvPolynomial}(\{0, \dots, r-1\}, R)\)
  is a Noetherian ring (Hilbert's basis theorem in finitely many variables). With
  \(R = \kappa\) a field this makes \(\kappa[t_0, \dots, t_{r-1}]\) Noetherian, the
  base ring over which the subquotient finiteness is measured
  (\cref{lem:graded_polyQuot_finite_of_le_numerator}).
\end{lemma}

\begin{lemma}[A finite module over a Noetherian ring is Noetherian]
  \label{lem:isNoetherian_of_isNoetherianRing_of_finite_mathlib}
  \lean{isNoetherian_of_isNoetherianRing_of_finite}
  \mathlibok
  \textit{Provided by Mathlib (\texttt{Mathlib.RingTheory.Noetherian.Defs}).}
  Let \(A\) be a Noetherian ring and \(M\) a finite \(A\)-module. Then \(M\) is a
  Noetherian \(A\)-module. Combined with
  \cref{lem:mvPolynomial_isNoetherianRing_fin_mathlib} this makes a finite
  \(\kappa[t_0, \dots, t_{r-1}]\)-module Noetherian, so its submodules are again
  finite (\cref{lem:graded_polyQuot_finite_of_le_numerator}).
\end{lemma}

\subsubsection*{The internal grading of a homogeneous submodule (split $G_1$)}

\begin{lemma}\leanok
[$G_1$ (independence) -- the family $n \mapsto M_n \cap p$ is independent]
  \label{lem:graded_homogeneousSubmodule_iSupIndep}
  \lean{AlgebraicGeometry.GradedModule.homogeneousSubmodule_inf_iSupIndep}
  \uses{lem:submodule_isHomogeneous_mathlib}
  Let \(M = \bigoplus_n M_n\) be an internally graded \(\kappa\)-module and let
  \(p \subseteq M\) be a homogeneous submodule
  (\cref{lem:submodule_isHomogeneous_mathlib}). Then the family
  \(n \mapsto M_n \cap p\) of \(\kappa\)-subspaces of \(M\) is independent: for each
  \(n\),
  \[
    (M_n \cap p) \;\cap\; \Bigl(\bigsqcup_{m \ne n} (M_m \cap p)\Bigr) \;=\; 0 .
  \]
  The whole statement lives in the ambient module \(M\): the subspaces
  \(M_n \cap p\) are submodules of \(M\), never of a derived carrier \(\smash{\hat p}\).
\end{lemma}

\begin{proof}
  \uses{lem:submodule_isHomogeneous_mathlib}
  \leanok
  The components \(M_n\) of the internal decomposition of \(M\) are independent. The
  subfamily \(n \mapsto M_n \cap p\) consists of submodules with \(M_n \cap p \le M_n\),
  so independence is inherited: a relation among the \(M_n \cap p\) is in particular a
  relation among the \(M_n\), forcing each term to vanish.
\end{proof}

\begin{lemma}\leanok
[$G_1$ (supremum) -- $\bigsqcup_n (M_n \cap p) = p$ for $p$ homogeneous]
  \label{lem:graded_homogeneousSubmodule_iSup_eq}
  \lean{AlgebraicGeometry.GradedModule.homogeneousSubmodule_iSup_inf_eq}
  \uses{lem:submodule_isHomogeneous_mathlib}
  Let \(p \subseteq M\) be a homogeneous submodule
  (\cref{lem:submodule_isHomogeneous_mathlib}). Then
  \[
    \bigsqcup_{n} (M_n \cap p) \;=\; p .
  \]
  Together with the independence of
  \cref{lem:graded_homogeneousSubmodule_iSupIndep}, this is the ambient internal
  direct sum \(p = \bigoplus_n (M_n \cap p)\), stated entirely inside \(M\) without
  requiring a direct-sum decomposition of the subtype \(\smash{\hat p}\).
\end{lemma}

\begin{proof}
  \uses{lem:submodule_isHomogeneous_mathlib}
  \leanok
  The inclusion \(\bigsqcup_n (M_n \cap p) \le p\) is clear, since each
  \(M_n \cap p \le p\). For the reverse inclusion, take \(m \in p\) and write
  \(m = \sum_n \mathrm{proj}_n(m)\) with \(\mathrm{proj}_n(m) \in M_n\) using the
  internal decomposition of \(M\). Homogeneity of \(p\) gives
  \(\mathrm{proj}_n(m) \in p\), hence \(\mathrm{proj}_n(m) \in M_n \cap p\); thus
  \(m\) is a finite sum of elements of the \(M_n \cap p\), i.e.
  \(m \in \bigsqcup_n (M_n \cap p)\).
\end{proof}

\subsubsection*{The ambient subquotient class and its Hilbert function}

\begin{definition}\leanok
[Subquotient Hilbert data]
  \label{def:graded_subquotientHilb}
  \lean{AlgebraicGeometry.GradedModule.subquotientHilb}
  \lean{AlgebraicGeometry.GradedModule.SubquotientDatum}
  \lean{AlgebraicGeometry.GradedModule.SubquotientDatum.hilb}
  \uses{lem:graded_homogeneousSubmodule_iSupIndep,
    lem:graded_homogeneousSubmodule_iSup_eq}
  Fix a graded \(\kappa\)-module \(\mathcal{M}\), i.e.\ \(M\) together with an
  internal decomposition \(M = \bigoplus_n \mathcal{M}_n\) whose components
  \(\mathcal{M}_n\) are finite-dimensional over \(\kappa\). A \emph{subquotient
  datum of length \(r\)} consists of:
  \begin{itemize}
    \item a pair of homogeneous submodules \(N' \le N \subseteq M\); and
    \item a family \(t : \{0, \dots, r-1\} \to (M \to_\kappa M)\) of pairwise
      commuting \(\kappa\)-linear endomorphisms, each raising degree by one
      (\(t_i(\mathcal{M}_n) \subseteq \mathcal{M}_{n+1}\)), each preserving \(N\)
      and \(N'\) (\(t_i N \subseteq N\), \(t_i N' \subseteq N'\)); together with
    \item a finiteness condition: the subquotient \(N/N'\) is a finite module over the
      polynomial ring \(\kappa[t_0, \dots, t_{r-1}]\) acting through the \(t_i\) (via the
      ambient \(t\)-stable carriers \(\mathrm{polySubmodule}\,N\),
      \(\mathrm{polySubmodule}\,N'\) of \cref{def:graded_polySubmodule}).
  \end{itemize}
  The pair \((N, N')\) represents the subquotient \(N/N'\). Its \emph{ambient
  Hilbert function} is
  \[
    \mathrm{hilb}(n) \;=\; \Bigl(\dim_\kappa (N \cap \mathcal{M}_n)
      \;-\; \dim_\kappa (N' \cap \mathcal{M}_n)\Bigr) \;\in\; \mathbb{Q},
  \]
  the difference of the (finite) \(\kappa\)-dimensions of the ambient intersections,
  computed in \(\mathbb{Z}\) and cast to \(\mathbb{Q}\). Since \(N' \le N\) one has
  \(N' \cap \mathcal{M}_n \le N \cap \mathcal{M}_n\), so
  \(\mathrm{hilb}(n) = \dim_\kappa\bigl((N \cap \mathcal{M}_n)/(N' \cap \mathcal{M}_n)\bigr)\)
  is the dimension of the degree-\(n\) piece of \(N/N'\). The split \(G_1\)
  (\cref{lem:graded_homogeneousSubmodule_iSupIndep},
  \cref{lem:graded_homogeneousSubmodule_iSup_eq}) guarantees that these ambient
  intersections assemble the homogeneous submodules \(N, N'\), so the data depends
  only on the ambient families \(n \mapsto N \cap \mathcal{M}_n\),
  \(n \mapsto N' \cap \mathcal{M}_n\).
\end{definition}

\subsubsection*{Ambient homogeneity calculus}

The kernel/cokernel closure and the degreewise difference identity rest on a small
calculus of how a degree-raising endomorphism interacts with the internal grading and
with the lattice of homogeneous submodules. Throughout, \(M = \bigoplus_n \mathcal{M}_n\)
is the fixed internally graded \(\kappa\)-module and \(x : M \to M\) is a
\(\kappa\)-linear endomorphism raising degree by one. None of these statements forms a
derived carrier: every object is a submodule of the fixed \(M\).

\begin{definition}\leanok
[Degree-raising endomorphism]
  \label{def:graded_raisesDegree}
  \lean{AlgebraicGeometry.GradedModule.RaisesDegree}
  Let \(M = \bigoplus_n \mathcal{M}_n\) be an internally graded \(\kappa\)-module. A
  \(\kappa\)-linear endomorphism \(x : M \to M\) \emph{raises degree by one} when it
  carries each graded piece into the next: \(x(\mathcal{M}_n) \subseteq
  \mathcal{M}_{n+1}\) for every \(n\). This is the grading-ring-free form of
  ``multiplication by a degree-one element'': it abstracts each of the \(r\)
  commuting degree-one generators \(x \in R_1\) acting on \(M\) into a single
  endomorphism, with no reference to the grading ring \(R\).
\end{definition}

\begin{lemma}\leanok
[Membership form of degree-raising]
  \label{lem:graded_raisesDegree_mem}
  \lean{AlgebraicGeometry.GradedModule.RaisesDegree.mem}
  \uses{def:graded_raisesDegree}
  If \(x\) raises degree by one (\cref{def:graded_raisesDegree}) and \(m \in
  \mathcal{M}_n\), then \(x m \in \mathcal{M}_{n+1}\).
\end{lemma}

\begin{proof}
  \uses{def:graded_raisesDegree}
  \leanok
  Immediate from the definition: \(x m \in x(\mathcal{M}_n) \subseteq
  \mathcal{M}_{n+1}\).
\end{proof}

\begin{lemma}\leanok
[Degree-shift commutation]
  \label{lem:graded_decompose_raisesDegree}
  \lean{AlgebraicGeometry.GradedModule.decompose_raisesDegree}
  \uses{def:graded_raisesDegree, lem:graded_raisesDegree_mem}
  Let \(x\) raise degree by one. Then for every \(m \in M\) and every \(i\) the
  degree-\((i+1)\) homogeneous component of \(x m\) equals \(x\) applied to the
  degree-\(i\) component of \(m\):
  \[
    \mathrm{proj}_{i+1}(x m) \;=\; x\bigl(\mathrm{proj}_i(m)\bigr) .
  \]
  This load-bearing commutation says a degree-raising endomorphism shifts the
  internal decomposition by exactly one slot; it is what makes preimages and images
  of homogeneous submodules under \(x\) homogeneous.
\end{lemma}

\begin{proof}
  \uses{def:graded_raisesDegree, lem:graded_raisesDegree_mem}
  \leanok
  Decompose \(m = \sum_j \mathrm{proj}_j(m)\) into its finitely many homogeneous
  components and apply \(x\): \(x m = \sum_j x(\mathrm{proj}_j(m))\) with
  \(x(\mathrm{proj}_j(m)) \in \mathcal{M}_{j+1}\) by
  \cref{lem:graded_raisesDegree_mem}. Reading off the degree-\((i+1)\) component, all
  terms with \(j+1 \ne i+1\) drop out, leaving \(x(\mathrm{proj}_i(m))\).
\end{proof}

\begin{lemma}\leanok
[Degree-raising kills the bottom]
  \label{lem:graded_decompose_raisesDegree_zero}
  \lean{AlgebraicGeometry.GradedModule.decompose_raisesDegree_zero}
  \uses{def:graded_raisesDegree}
  Let \(x\) raise degree by one. Then the degree-\(0\) component of \(x m\) vanishes
  for every \(m\): \(\mathrm{proj}_0(x m) = 0\).
\end{lemma}

\begin{proof}
  \uses{def:graded_raisesDegree}
  \leanok
  Writing \(x m = \sum_j x(\mathrm{proj}_j(m))\) with each summand in
  \(\mathcal{M}_{j+1}\), no summand has degree \(0\) (every \(j+1 \ge 1\)), so the
  degree-\(0\) component is \(0\).
\end{proof}

\begin{lemma}\leanok
[Preimage of a homogeneous submodule is homogeneous]
  \label{lem:graded_comap_isHomogeneous}
  \lean{AlgebraicGeometry.GradedModule.comap_isHomogeneous}
  \uses{lem:submodule_isHomogeneous_mathlib, lem:graded_decompose_raisesDegree}
  Let \(x\) raise degree by one and let \(N' \subseteq M\) be a homogeneous submodule
  (\cref{lem:submodule_isHomogeneous_mathlib}). Then the preimage \(x^{-1}(N') =
  \{m : x m \in N'\}\) is a homogeneous submodule of \(M\).
\end{lemma}

\begin{proof}
  \uses{lem:submodule_isHomogeneous_mathlib, lem:graded_decompose_raisesDegree}
  \leanok
  Let \(m \in x^{-1}(N')\) and fix a degree \(i\). By the degree-shift commutation
  (\cref{lem:graded_decompose_raisesDegree}), \(x(\mathrm{proj}_i(m)) =
  \mathrm{proj}_{i+1}(x m)\), which lies in \(N'\) because \(x m \in N'\) and \(N'\)
  is homogeneous. Hence \(\mathrm{proj}_i(m) \in x^{-1}(N')\), so \(x^{-1}(N')\) is
  homogeneous.
\end{proof}

\begin{lemma}\leanok
[Image of a homogeneous submodule is homogeneous]
  \label{lem:graded_map_isHomogeneous}
  \lean{AlgebraicGeometry.GradedModule.map_isHomogeneous}
  \uses{lem:submodule_isHomogeneous_mathlib, lem:graded_decompose_raisesDegree,
    lem:graded_decompose_raisesDegree_zero}
  Let \(x\) raise degree by one and let \(N \subseteq M\) be a homogeneous submodule.
  Then the image \(x \cdot N = \{x m : m \in N\}\) is a homogeneous submodule of
  \(M\), with degree-\((n+1)\) piece \(x(N \cap \mathcal{M}_n)\) and zero
  degree-\(0\) piece.
\end{lemma}

\begin{proof}
  \uses{lem:submodule_isHomogeneous_mathlib, lem:graded_decompose_raisesDegree,
    lem:graded_decompose_raisesDegree_zero}
    \leanok
  Take \(x m \in x \cdot N\) with \(m \in N\) and fix a degree. In degree \(0\) the
  component \(\mathrm{proj}_0(x m) = 0 \in x \cdot N\)
  (\cref{lem:graded_decompose_raisesDegree_zero}). In degree \(i+1\),
  \(\mathrm{proj}_{i+1}(x m) = x(\mathrm{proj}_i(m))\)
  (\cref{lem:graded_decompose_raisesDegree}); since \(N\) is homogeneous,
  \(\mathrm{proj}_i(m) \in N\), so this component lies in \(x \cdot N\). Thus every
  homogeneous component of \(x m\) lies in \(x \cdot N\), which is therefore
  homogeneous.
\end{proof}

\begin{lemma}\leanok
[Intersection of homogeneous submodules is homogeneous]
  \label{lem:graded_inf_isHomogeneous}
  \lean{AlgebraicGeometry.GradedModule.inf_isHomogeneous}
  \uses{lem:submodule_isHomogeneous_mathlib}
  If \(p, q \subseteq M\) are homogeneous submodules
  (\cref{lem:submodule_isHomogeneous_mathlib}) then so is their intersection \(p
  \cap q\). Mathlib provides no lattice-closure lemmas for homogeneous submodules, so
  this is recorded here.
\end{lemma}

\begin{proof}
  \uses{lem:submodule_isHomogeneous_mathlib}
  \leanok
  For \(m \in p \cap q\) and any degree \(i\), homogeneity of \(p\) and of \(q\)
  gives \(\mathrm{proj}_i(m) \in p\) and \(\mathrm{proj}_i(m) \in q\), hence
  \(\mathrm{proj}_i(m) \in p \cap q\).
\end{proof}

\begin{lemma}\leanok
[Sum of homogeneous submodules is homogeneous]
  \label{lem:graded_sup_isHomogeneous}
  \lean{AlgebraicGeometry.GradedModule.sup_isHomogeneous}
  \uses{lem:submodule_isHomogeneous_mathlib}
  If \(p, q \subseteq M\) are homogeneous submodules then so is their sum \(p + q\).
\end{lemma}

\begin{proof}
  \uses{lem:submodule_isHomogeneous_mathlib}
  \leanok
  Write an element of \(p + q\) as \(a + b\) with \(a \in p\), \(b \in q\). For any
  degree \(i\), \(\mathrm{proj}_i(a+b) = \mathrm{proj}_i(a) + \mathrm{proj}_i(b)\)
  with \(\mathrm{proj}_i(a) \in p\) and \(\mathrm{proj}_i(b) \in q\) by homogeneity,
  so \(\mathrm{proj}_i(a+b) \in p + q\).
\end{proof}

\begin{lemma}\leanok
[Ambient image identity]
  \label{lem:graded_map_inf_degree_eq}
  \lean{AlgebraicGeometry.GradedModule.map_inf_degree_eq}
  \uses{lem:graded_decompose_raisesDegree, lem:graded_map_isHomogeneous}
  Let \(x\) raise degree by one and let \(N \subseteq M\) be homogeneous. Then the
  degree-\((n+1)\) part of the image \(x \cdot N\) is exactly \(x\) of the
  degree-\(n\) part of \(N\):
  \[
    (x \cdot N) \cap \mathcal{M}_{n+1} \;=\; x \cdot (N \cap \mathcal{M}_n) .
  \]
\end{lemma}

\begin{proof}
  \uses{lem:graded_decompose_raisesDegree, lem:graded_map_isHomogeneous}
  \leanok
  \(\supseteq\): for \(m \in N \cap \mathcal{M}_n\), \(x m \in x \cdot N\) and \(x m
  \in \mathcal{M}_{n+1}\) since \(x\) raises degree. \(\subseteq\): an element \(x m
  \in x \cdot N\) lying in \(\mathcal{M}_{n+1}\) equals its own degree-\((n+1)\)
  component \(\mathrm{proj}_{n+1}(x m) = x(\mathrm{proj}_n(m))\)
  (\cref{lem:graded_decompose_raisesDegree}); homogeneity of \(N\) puts
  \(\mathrm{proj}_n(m) \in N \cap \mathcal{M}_n\), so \(x m \in x \cdot (N \cap
  \mathcal{M}_n)\). The image \(x \cdot N\) is homogeneous by
  \cref{lem:graded_map_isHomogeneous}.
\end{proof}

\begin{lemma}\leanok
[Ambient distributive law]
  \label{lem:graded_sup_inf_degree_eq}
  \lean{AlgebraicGeometry.GradedModule.sup_inf_degree_eq}
  \uses{lem:graded_inf_isHomogeneous, lem:graded_sup_isHomogeneous,
    lem:submodule_isHomogeneous_mathlib}
  For homogeneous submodules \(P, Q \subseteq M\) and any degree \(k\), intersecting
  the sum with a graded piece distributes:
  \[
    (P + Q) \cap \mathcal{M}_k \;=\; (P \cap \mathcal{M}_k) + (Q \cap \mathcal{M}_k) .
  \]
\end{lemma}

\begin{proof}
  \uses{lem:graded_inf_isHomogeneous, lem:graded_sup_isHomogeneous,
    lem:submodule_isHomogeneous_mathlib}
    \leanok
  \(\supseteq\) is the general lattice inequality. For \(\subseteq\), take \(z = p +
  q \in (P + Q) \cap \mathcal{M}_k\) with \(p \in P\), \(q \in Q\). Since \(z \in
  \mathcal{M}_k\), \(z = \mathrm{proj}_k(z) = \mathrm{proj}_k(p) +
  \mathrm{proj}_k(q)\); homogeneity of \(P\) and \(Q\)
  (\cref{lem:submodule_isHomogeneous_mathlib}) puts \(\mathrm{proj}_k(p) \in P \cap
  \mathcal{M}_k\) and \(\mathrm{proj}_k(q) \in Q \cap \mathcal{M}_k\), so \(z \in (P
  \cap \mathcal{M}_k) + (Q \cap \mathcal{M}_k)\). The intersection and sum stay in
  the homogeneous lattice by \cref{lem:graded_inf_isHomogeneous} and
  \cref{lem:graded_sup_isHomogeneous}.
\end{proof}

\subsubsection*{Closure under kernel and cokernel}

The kernel/cokernel closure of the subquotient class is realised by eight ambient
component facts: for a degree-raising endomorphism \(x\) and a homogeneous pair
\(N' \le N\), the kernel pair is \((N \cap x^{-1}(N'),\, N')\) and the cokernel pair is
\((N,\, N' + x \cdot N)\). The lemmas below record that both new pairs are homogeneous,
nest correctly, are annihilated by \(x\), and are preserved by any endomorphism \(t\)
commuting with \(x\) that preserves the original pair.

\begin{lemma}\leanok
[Homogeneity of the kernel pair's lower module]
  \label{lem:graded_ker_isHomogeneous}
  \lean{AlgebraicGeometry.GradedModule.ker_isHomogeneous}
  \uses{lem:graded_inf_isHomogeneous, lem:graded_comap_isHomogeneous}
  Let \(x\) raise degree by one (\cref{def:graded_raisesDegree}) and let \(N, N'\) be
  homogeneous submodules of \(M\). Then the kernel pair's lower module
  \(N \cap x^{-1}(N')\) is homogeneous. Indeed it is the intersection
  (\cref{lem:graded_inf_isHomogeneous}) of the homogeneous \(N\) with the homogeneous
  preimage \(x^{-1}(N')\) (\cref{lem:graded_comap_isHomogeneous}).
\end{lemma}

\begin{lemma}\leanok
[Homogeneity of the cokernel pair's upper module]
  \label{lem:graded_coker_isHomogeneous}
  \lean{AlgebraicGeometry.GradedModule.coker_isHomogeneous}
  \uses{lem:graded_sup_isHomogeneous, lem:graded_map_isHomogeneous}
  Under the same hypotheses, the cokernel pair's upper module \(N' + x \cdot N\) is
  homogeneous: it is the sum (\cref{lem:graded_sup_isHomogeneous}) of the homogeneous
  \(N'\) with the homogeneous image \(x \cdot N\)
  (\cref{lem:graded_map_isHomogeneous}).
\end{lemma}

\begin{lemma}\leanok
[The kernel pair nests]
  \label{lem:graded_ker_le}
  \lean{AlgebraicGeometry.GradedModule.ker_le}
  If \(N' \le N\) and \(x\) preserves \(N'\) (\(x \cdot N' \subseteq N'\)), then
  \(N' \le N \cap x^{-1}(N')\). Indeed \(N' \le N\) by hypothesis, and \(x \cdot N'
  \subseteq N'\) is exactly \(N' \subseteq x^{-1}(N')\); so \(N'\) is below the
  intersection. Hence the kernel pair \((N \cap x^{-1}(N'),\, N')\) is a valid
  subquotient pair (\cref{def:graded_subquotientHilb}).
\end{lemma}

\begin{lemma}\leanok
[The cokernel pair nests]
  \label{lem:graded_coker_le}
  \lean{AlgebraicGeometry.GradedModule.coker_le}
  If \(N' \le N\) and \(x\) preserves \(N\) (\(x \cdot N \subseteq N\)), then
  \(N' + x \cdot N \le N\): both summands lie in \(N\). Hence the cokernel pair
  \((N,\, N' + x \cdot N)\) is a valid subquotient pair
  (\cref{def:graded_subquotientHilb}).
\end{lemma}

\begin{lemma}\leanok
[$x$ annihilates the kernel subquotient]
  \label{lem:graded_ker_annihilate}
  \lean{AlgebraicGeometry.GradedModule.ker_annihilate}
  The endomorphism \(x\) carries the kernel pair's lower module into \(N'\):
  \(x \cdot (N \cap x^{-1}(N')) \subseteq N'\). Indeed \(N \cap x^{-1}(N') \subseteq
  x^{-1}(N')\), and applying \(x\) to a preimage of \(N'\) lands in \(N'\). Thus \(x\)
  acts as zero on the kernel subquotient.
\end{lemma}

\begin{lemma}\leanok
[$x$ annihilates the cokernel subquotient]
  \label{lem:graded_coker_annihilate}
  \lean{AlgebraicGeometry.GradedModule.coker_annihilate}
  The endomorphism \(x\) carries \(N\) into the cokernel pair's upper module:
  \(x \cdot N \subseteq N' + x \cdot N\), trivially (the right summand). Thus \(x\)
  acts as zero on the cokernel subquotient \(N/(N' + x \cdot N)\).
\end{lemma}

\begin{lemma}\leanok
[Commuting endomorphisms preserve the preimage]
  \label{lem:graded_comap_map_le_of_commute}
  \lean{AlgebraicGeometry.GradedModule.comap_map_le_of_commute}
  Let \(t\) commute with \(x\) and preserve \(N'\) (\(t \cdot N' \subseteq N'\)). Then
  \(t\) preserves the preimage: \(t \cdot x^{-1}(N') \subseteq x^{-1}(N')\). Indeed for
  \(m\) with \(x m \in N'\) one has \(x(t m) = t(x m) \in t \cdot N' \subseteq N'\), so
  \(t m \in x^{-1}(N')\). This is what shows the remaining \(t_i\) preserve the kernel
  pair's lower module \(N \cap x^{-1}(N')\).
\end{lemma}

\begin{lemma}\leanok
[Commuting endomorphisms preserve the image]
  \label{lem:graded_map_map_le_of_commute}
  \lean{AlgebraicGeometry.GradedModule.map_map_le_of_commute}
  Let \(t\) commute with \(x\) and preserve \(N\) (\(t \cdot N \subseteq N\)). Then
  \(t\) preserves the image: \(t \cdot (x \cdot N) \subseteq x \cdot N\). Indeed for
  \(m \in N\) one has \(t(x m) = x(t m)\) with \(t m \in N\), so \(t(x m) \in x \cdot
  N\). This is what shows the remaining \(t_i\) preserve the cokernel pair's upper
  module \(N' + x \cdot N\).
\end{lemma}

\begin{lemma}\leanok
[Kernel/cokernel closure of the subquotient class]
  \label{lem:graded_subquotient_ker_coker}
  \lean{AlgebraicGeometry.GradedModule.SubquotientDatum.ker}
  % This is a narrative grouping lemma: its content is realised in Lean by the eight
  % ambient component facts cited in \uses{} below (no single Lean declaration),
  % together with the SubquotientDatum.ker/.coker constructors that bundle them into a
  % length-(r-1) datum. The `\lean` pin names the `.ker` constructor — the single decl most
  % representative of the bundled closure — so this dependency-bearing grouping node (which is
  % `\uses`'d by the keystone rationality induction) stays wired into the graph rather than
  % orphaning the component facts it `\uses`. This is a deliberate duplicate of the
  % `def:graded_subquotientDatum_ker` pin.
  \uses{lem:graded_ker_isHomogeneous, lem:graded_coker_isHomogeneous,
    lem:graded_ker_le, lem:graded_coker_le, lem:graded_ker_annihilate,
    lem:graded_coker_annihilate, lem:graded_comap_map_le_of_commute,
    lem:graded_map_map_le_of_commute,
    def:graded_subquotientHilb, lem:graded_homogeneousSubmodule_iSupIndep,
    lem:graded_homogeneousSubmodule_iSup_eq, lem:ideal_homogeneous_span_mathlib,
    lem:isHomogeneousElem_graded_smul_mathlib, lem:graded_comap_isHomogeneous,
    lem:graded_map_isHomogeneous, lem:graded_inf_isHomogeneous,
    lem:graded_sup_isHomogeneous}
  Let \((N, N')\) be a subquotient datum (\cref{def:graded_subquotientHilb}) with
  endomorphisms \(t_0, \dots, t_{r-1}\), and set \(x = t_{r-1}\). Form the two new
  pairs of homogeneous submodules of \(M\)
  \[
    K \;=\; \bigl(\,N \cap x^{-1}(N'),\ N'\,\bigr),
    \qquad
    C \;=\; \bigl(\,N,\ N' + x \cdot N\,\bigr),
  \]
  where \(x^{-1}(N') = \{\,m \in M : x \cdot m \in N'\,\}\) is the preimage and
  \(x \cdot N\) is the image of \(N\) under \(x\). Then:
  \begin{enumerate}
    \item \(N \cap x^{-1}(N')\) and \(N' + x \cdot N\) are homogeneous submodules of
      \(M\), with \(N' \le N \cap x^{-1}(N')\) and \(N' + x \cdot N \le N\), so both
      \(K\) and \(C\) are subquotient pairs;
    \item the remaining endomorphisms \(t_0, \dots, t_{r-2}\) pairwise commute,
      raise degree by one, and preserve both members of \(K\) and of \(C\); and
    \item \(x\) annihilates both subquotients: \(x \cdot (N \cap x^{-1}(N')) \subseteq N'\)
      and \(x \cdot N \subseteq N' + x \cdot N\).
  \end{enumerate}
  Consequently \(K\) and \(C\) are again subquotient data, now of length \(r-1\)
  (the endomorphisms \(t_0, \dots, t_{r-2}\)), and \(K = \ker x\) and \(C =
  \operatorname{coker} x\) on the subquotient \(N/N'\): \(K\) represents
  \(\ker\bigl(x : N/N' \to N/N'\bigr)\) and \(C\) represents
  \(\operatorname{coker}\bigl(x : N/N' \to N/N'\bigr)\). No quotient or subtype
  carrier is constructed; \(K\) and \(C\) are ambient pairs of homogeneous
  submodules of the fixed \(M\).
\end{lemma}

\begin{proof}
  \uses{lem:graded_ker_isHomogeneous, lem:graded_coker_isHomogeneous,
    lem:graded_ker_le, lem:graded_coker_le, lem:graded_ker_annihilate,
    lem:graded_coker_annihilate, lem:graded_comap_map_le_of_commute,
    lem:graded_map_map_le_of_commute,
    lem:graded_homogeneousSubmodule_iSupIndep,
    lem:graded_homogeneousSubmodule_iSup_eq, lem:ideal_homogeneous_span_mathlib,
    lem:isHomogeneousElem_graded_smul_mathlib, lem:graded_comap_isHomogeneous,
    lem:graded_map_isHomogeneous, lem:graded_inf_isHomogeneous,
    lem:graded_sup_isHomogeneous}
    \leanok
  \emph{Homogeneity.} The preimage \(x^{-1}(N')\) is homogeneous: if
  \(m = \sum_n \mathrm{proj}_n(m)\) with \(x \cdot m \in N'\), then since \(x\) raises
  degree by one (\cref{lem:isHomogeneousElem_graded_smul_mathlib}) the components
  \(x \cdot \mathrm{proj}_n(m) \in \mathcal{M}_{n+1}\) lie in distinct degrees, and
  homogeneity of \(N'\) (split \(G_1\),
  \cref{lem:graded_homogeneousSubmodule_iSup_eq}) forces each
  \(x \cdot \mathrm{proj}_n(m) \in N'\), i.e.\ \(\mathrm{proj}_n(m) \in x^{-1}(N')\).
  Thus \(N \cap x^{-1}(N')\) is an intersection of homogeneous submodules, hence
  homogeneous. The image \(x \cdot N\) is homogeneous with degree-\((n+1)\) piece
  \(x \cdot (N \cap \mathcal{M}_n)\) by
  \cref{lem:isHomogeneousElem_graded_smul_mathlib}, and a sum of homogeneous
  submodules is homogeneous (\cref{lem:ideal_homogeneous_span_mathlib} applied at the
  module level), so \(N' + x \cdot N\) is homogeneous. The inclusions
  \(N' \le N \cap x^{-1}(N')\) (as \(x \cdot N' \subseteq N'\)) and
  \(N' + x \cdot N \le N\) (as \(N' \le N\) and \(x \cdot N \subseteq N\)) are
  immediate. In the language of the ambient homogeneity calculus, this paragraph is
  exactly: \(x^{-1}(N')\) is homogeneous (\cref{lem:graded_comap_isHomogeneous}),
  \(x \cdot N\) is homogeneous (\cref{lem:graded_map_isHomogeneous}), and the
  homogeneous lattice is closed under intersection
  (\cref{lem:graded_inf_isHomogeneous}) and sum
  (\cref{lem:graded_sup_isHomogeneous}).

  \emph{Preservation.} Each \(t_i\) with \(i \le r-2\) commutes with \(x\) and
  preserves \(N, N'\); hence it preserves \(x^{-1}(N')\) (if \(x t_i m = t_i x m \in
  N'\) whenever \(x m \in N'\)), so it preserves \(N \cap x^{-1}(N')\), and it
  preserves \(x \cdot N\) (as \(t_i x N = x t_i N \subseteq x N\)), so it preserves
  \(N' + x \cdot N\). Degree-raising and commutativity are inherited.

  \emph{Annihilation by \(x\).} For \(m \in N \cap x^{-1}(N')\) one has
  \(x \cdot m \in N'\) by definition, so \(x\) carries the first member of \(K\) into
  its second; and \(x \cdot N \subseteq N' + x \cdot N\) trivially. Thus \(x\) acts as
  zero on \(N/N'\) restricted to \(K\) and on the cokernel quotient \(C\).

  \emph{Identification with \(\ker x\) and \(\operatorname{coker} x\).} On the
  subquotient \(N/N'\), the class of \(m \in N\) is killed by \(x\) iff
  \(x \cdot m \in N'\), i.e.\ iff \(m \in N \cap x^{-1}(N')\); so \(\ker(x : N/N' \to
  N/N')\) is exactly \((N \cap x^{-1}(N'))/N'\), represented by \(K\). The image of
  \(x\) on \(N/N'\) is \((N' + x \cdot N)/N'\), so the cokernel
  \((N/N')/\operatorname{im} x = N/(N' + x \cdot N)\) is represented by \(C\).
\end{proof}

\subsubsection*{The subquotient degreewise difference identity ($D_6$)}

\begin{lemma}

\leanok
[preimage dimension equals ambient-intersection dimension]
  \label{lem:graded_finrank_comap_subtype}
  \lean{AlgebraicGeometry.GradedModule.finrank_comap_subtype}
  For submodules \(p, S\) of a \(\kappa\)-vector space \(M\), the dimension of the
  preimage of \(S\) under the inclusion \(p \hookrightarrow M\) equals the dimension of
  the ambient intersection: \(\dim_\kappa(\iota_p^{-1} S) = \dim_\kappa(p \cap S)\),
  where \(\iota_p : p \to M\) is the canonical injection. Project-local helper consumed
  by the two kernel-dimension computations of \cref{lem:graded_subquotient_degreewise_diff}.
\end{lemma}
\begin{proof}

\leanok
  The inclusion \(\iota_p\) is injective, so it restricts to a \(\kappa\)-linear
  isomorphism from \(\iota_p^{-1} S\) onto its image \(\iota_p(\iota_p^{-1} S) = p \cap S\)
  (image of a comap under an injective map). An isomorphism preserves dimension, giving
  the claimed equality.
\end{proof}

\begin{lemma}\leanok
[$D_6$ -- subquotient degreewise difference]
  \label{lem:graded_subquotient_degreewise_diff}
  \lean{AlgebraicGeometry.GradedModule.subquotient_degreewise_diff}
  \uses{def:graded_subquotientHilb, lem:graded_subquotient_ker_coker,
    lem:graded_degreewise_finrank_diff, lem:graded_homogeneousSubmodule_iSupIndep,
    lem:graded_homogeneousSubmodule_iSup_eq, lem:graded_finrank_comap_subtype,
    lem:isHomogeneousElem_graded_smul_mathlib, lem:graded_map_inf_degree_eq,
    lem:graded_sup_inf_degree_eq}
  % NOTE: The proof factors the two kernel-dimension computations through the
  %   helper \cref{lem:graded_finrank_comap_subtype} (preimage dimension = ambient
  %   intersection dimension).
  % SOURCE: [Stacks Project], "Commutative Algebra", \S "Noetherian graded
  % rings", Proposition \texttt{proposition-graded-hilbert-polynomial} (Tag
  % 00K1), proof (read from references/hilbert-serre-algebra.tex, L13939-L13947).
  \textit{Source: [Stacks Project], Tag 00K1 (degreewise short exact sequence).}
  Let \((N, N')\) be a subquotient datum (\cref{def:graded_subquotientHilb}) with
  Hilbert function \(h_M\), and let \(x = t_{r-1}\). Let \(h_K\) and \(h_C\) be the
  Hilbert functions of the kernel and cokernel subquotients
  \(K = (N \cap x^{-1}(N'),\, N')\) and \(C = (N,\, N' + x \cdot N)\) of
  \cref{lem:graded_subquotient_ker_coker}. Then for every \(n\)
  \[
    h_M(n+1) - h_M(n) \;=\; h_C(n+1) - h_K(n).
  \]
  Writing \(Q_n = (N \cap \mathcal{M}_n)/(N' \cap \mathcal{M}_n)\) for the
  degree-\(n\) piece of \(N/N'\) (so \(\dim_\kappa Q_n = h_M(n)\)), multiplication by
  \(x\) is a \(\kappa\)-linear map \(x : Q_n \to Q_{n+1}\) with
  \(\ker(x : Q_n \to Q_{n+1})\) the degree-\(n\) piece of \(K\) and
  \(Q_{n+1}/\operatorname{im} x\) the degree-\((n+1)\) piece of \(C\); the identity is
  then the degreewise rank--nullity difference \(D_5\)
  (\cref{lem:graded_degreewise_finrank_diff}). Specialising to \((N, N') = (\top,
  \bot)\) recovers the classical \(0 \to \mathcal{M}_n \xrightarrow{x}
  \mathcal{M}_{n+1} \to \mathcal{M}_{n+1}/x\mathcal{M}_n \to 0\).
\end{lemma}

% SOURCE QUOTE PROOF: [Stacks Project], Tag 00K1, proof, inductive step (read
% from references/hilbert-serre-algebra.tex, L13939-L13947).
% "Let $\overline{M} = M/xM$. Note that the map $x : M \to M$
%  is {\it not} a map of graded $S$-modules, since it does
%  not map $M_d$ into $M_d$. Namely, for each $d$ we have the
%  following short exact sequence
%  $$
%  0 \to M_d \xrightarrow{x} M_{d + 1} \to \overline{M}_{d + 1} \to 0
%  $$
%  This proves that $[M_{d + 1}] - [M_d] = [\overline{M}_{d + 1}]$.
%  Hence we win by Lemma \ref{lemma-numerical-polynomial}."
\begin{proof}
  \uses{lem:graded_subquotient_ker_coker, lem:graded_degreewise_finrank_diff,
    lem:graded_homogeneousSubmodule_iSup_eq,
    lem:isHomogeneousElem_graded_smul_mathlib, lem:graded_map_inf_degree_eq,
    lem:graded_sup_inf_degree_eq}
    \leanok
  Since \(N' \le N\) are homogeneous, the split \(G_1\)
  (\cref{lem:graded_homogeneousSubmodule_iSup_eq}) identifies the degree-\(n\) piece
  of \(N/N'\) with \(Q_n = (N \cap \mathcal{M}_n)/(N' \cap \mathcal{M}_n)\), a
  finite-dimensional \(\kappa\)-vector space of dimension \(h_M(n)\). Because \(x\)
  raises degree by one and preserves \(N\) and \(N'\)
  (\cref{lem:isHomogeneousElem_graded_smul_mathlib}), it induces a \(\kappa\)-linear
  map \(x : Q_n \to Q_{n+1}\).

  \emph{Kernel.} A class \([m] \in Q_n\) (with \(m \in N \cap \mathcal{M}_n\)) lies in
  \(\ker(x : Q_n \to Q_{n+1})\) iff \(x \cdot m \in N' \cap \mathcal{M}_{n+1}\), i.e.\
  iff \(m \in (N \cap x^{-1}(N')) \cap \mathcal{M}_n\). Hence \(\ker(x : Q_n \to
  Q_{n+1}) = \bigl((N \cap x^{-1}(N')) \cap \mathcal{M}_n\bigr)/(N' \cap \mathcal{M}_n)\)
  is exactly the degree-\(n\) piece of \(K\), of dimension \(h_K(n)\)
  (\cref{lem:graded_subquotient_ker_coker}).

  \emph{Cokernel.} The image of \(x : Q_n \to Q_{n+1}\) is
  \(\bigl(x \cdot (N \cap \mathcal{M}_n) + (N' \cap \mathcal{M}_{n+1})\bigr)/(N' \cap
  \mathcal{M}_{n+1})\). Using the ambient identity \((N' + x \cdot N) \cap
  \mathcal{M}_{n+1} = (N' \cap \mathcal{M}_{n+1}) + x \cdot (N \cap \mathcal{M}_n)\)
  -- which is the ambient distributive law (\cref{lem:graded_sup_inf_degree_eq})
  applied to \(P = N'\), \(Q = x \cdot N\), combined with the ambient image identity
  \((x \cdot N) \cap \mathcal{M}_{n+1} = x \cdot (N \cap \mathcal{M}_n)\)
  (\cref{lem:graded_map_inf_degree_eq}) -- the cokernel \(Q_{n+1}/\operatorname{im} x = (N \cap
  \mathcal{M}_{n+1})/\bigl((N' + x \cdot N) \cap \mathcal{M}_{n+1}\bigr)\) is the
  degree-\((n+1)\) piece of \(C\), of dimension \(h_C(n+1)\).

  \emph{Difference.} Apply the degreewise rank--nullity difference \(D_5\)
  (\cref{lem:graded_degreewise_finrank_diff}) to \(\varphi = (x : Q_n \to Q_{n+1})\),
  with \(V = Q_n\), \(W = Q_{n+1}\):
  \[
    h_M(n+1) - h_M(n)
      = \dim_\kappa Q_{n+1} - \dim_\kappa Q_n
      = \dim_\kappa(Q_{n+1}/\operatorname{im}\varphi) - \dim_\kappa(\ker\varphi)
      = h_C(n+1) - h_K(n).
  \]
  This is the \(\mathrm{hdiff}\) hypothesis consumed by \cref{lem:ratHilb_ofDiffEq}.
\end{proof}

\subsubsection*{The free polynomial-ring module structure}

The finiteness condition of a subquotient datum (\cref{def:graded_subquotientHilb}) and
its transfer down one generator (\cref{lem:graded_subquotient_finite_transfer}) are
phrased over the \emph{free} polynomial ring \(\kappa[t_0, \dots, t_{r-1}] =
\operatorname{MvPolynomial}(\{0, \dots, r-1\}, \kappa)\). The following blocks build the
action of that ring on the ambient module \(M\) from the \(r\) pairwise-commuting
degree-raising endomorphisms \(t_i\), keeping every carrier an ambient submodule of
\(M\). The free polynomial ring -- rather than the subalgebra
\(\mathrm{adjoin}_\kappa\{t_0, \dots, t_{r-1}\} \subseteq \operatorname{End}_\kappa(M)\)
-- is used precisely so the inductive transfer has the ring surjection
\(\kappa[t_0, \dots, t_{r-1}] \twoheadrightarrow \kappa[t_0, \dots, t_{r-2}]\) available
(\cref{lem:graded_subquotient_finite_transfer}).

\begin{definition}\leanok
[Polynomial evaluation ring homomorphism]
  \label{def:graded_polyEndHom}
  \lean{AlgebraicGeometry.GradedModule.polyEndHom}
  \uses{lem:adjoinCommRingOfComm_mathlib}
  Let \(t : \{0, \dots, r-1\} \to \operatorname{End}_\kappa(M)\) be a family of
  pairwise-commuting \(\kappa\)-linear endomorphisms. The \emph{polynomial evaluation
  ring homomorphism}
  \[
    \mathrm{polyEndHom}(t) : \kappa[t_0, \dots, t_{r-1}] \longrightarrow
      \operatorname{End}_\kappa(M)
  \]
  sends each variable \(X_i\) to \(t_i\). It is constructed by evaluating into the
  \emph{commutative} \(\kappa\)-subalgebra \(A = \mathrm{adjoin}_\kappa(\{t_i\})
  \subseteq \operatorname{End}_\kappa(M)\) -- commutative by
  \cref{lem:adjoinCommRingOfComm_mathlib} because the generators commute -- and then
  including \(A \hookrightarrow \operatorname{End}_\kappa(M)\); evaluation into the
  noncommutative \(\operatorname{End}_\kappa(M)\) directly is not an algebra
  homomorphism, which is why the commutative \(A\) is interposed.
\end{definition}

\begin{lemma}\leanok
[Evaluation sends $X_i$ to $t_i$]
  \label{lem:graded_polyEndHom_X}
  \lean{AlgebraicGeometry.GradedModule.polyEndHom_X}
  \uses{def:graded_polyEndHom}
  \(\mathrm{polyEndHom}(t)(X_i) = t_i\) for each \(i\). Immediate from the construction
  (\cref{def:graded_polyEndHom}), as evaluation sends \(X_i\) to the chosen image
  \(t_i\).
\end{lemma}

\begin{lemma}\leanok
[Evaluation sends a constant to a scalar]
  \label{lem:graded_polyEndHom_C}
  \lean{AlgebraicGeometry.GradedModule.polyEndHom_C}
  \uses{def:graded_polyEndHom}
  \(\mathrm{polyEndHom}(t)(C\,c) = c \cdot \mathrm{id}_M\) for each \(c \in \kappa\):
  the evaluation homomorphism (\cref{def:graded_polyEndHom}) is \(\kappa\)-linear and
  carries the constant polynomial \(C\,c\) to the scalar endomorphism \(c \cdot 1\).
\end{lemma}

\begin{definition}\leanok
[Polynomial-ring module structure on $M$]
  \label{def:graded_polyModule}
  \lean{AlgebraicGeometry.GradedModule.polyModule}
  \uses{def:graded_polyEndHom}
  The \emph{polynomial-module structure} is the \(\kappa[t_0, \dots, t_{r-1}]\)-module
  structure on \(M\) obtained by restricting scalars along
  \(\mathrm{polyEndHom}(t)\) (\cref{def:graded_polyEndHom}): a polynomial \(p\) acts on
  \(m \in M\) by \(p \cdot m = \mathrm{polyEndHom}(t)(p)(m)\). This is the module over the
  free polynomial ring on which the ambient subquotient induction's finiteness condition
  is measured.
\end{definition}

\begin{lemma}\leanok
[$X_i$ acts as $t_i$]
  \label{lem:graded_polyModule_X_smul}
  \lean{AlgebraicGeometry.GradedModule.polyModule_X_smul}
  \uses{def:graded_polyModule, lem:graded_polyEndHom_X}
  In the polynomial-module structure (\cref{def:graded_polyModule}), \(X_i \cdot m =
  t_i(m)\) for all \(m \in M\). Immediate from
  \(\mathrm{polyEndHom}(t)(X_i) = t_i\) (\cref{lem:graded_polyEndHom_X}).
\end{lemma}

\begin{lemma}\leanok
[A constant acts as a scalar]
  \label{lem:graded_polyModule_C_smul}
  \lean{AlgebraicGeometry.GradedModule.polyModule_C_smul}
  \uses{def:graded_polyModule, lem:graded_polyEndHom_C}
  In the polynomial-module structure (\cref{def:graded_polyModule}), \((C\,c) \cdot m =
  c \cdot m\) for all \(c \in \kappa\), \(m \in M\). Immediate from
  \(\mathrm{polyEndHom}(t)(C\,c) = c \cdot \mathrm{id}_M\)
  (\cref{lem:graded_polyEndHom_C}).
\end{lemma}

\begin{lemma}\leanok
[$\kappa$-scalar tower]
  \label{lem:graded_polyModule_isScalarTower}
  \lean{AlgebraicGeometry.GradedModule.polyModule_isScalarTower}
  \uses{def:graded_polyModule, lem:graded_polyEndHom_C}
  The polynomial-module structure (\cref{def:graded_polyModule}) is compatible with the
  original \(\kappa\)-action: \(\mathrm{IsScalarTower}\,\kappa\,\kappa[t_0, \dots,
  t_{r-1}]\,M\). Indeed the algebra map \(\kappa \to \kappa[t_0, \dots, t_{r-1}]\)
  followed by the polynomial action recovers the \(\kappa\)-action, since a scalar
  \(c \in \kappa\) maps to the constant \(C\,c\), which acts as \(c\)
  (\cref{lem:graded_polyEndHom_C}).
\end{lemma}

\begin{definition}\leanok
[Polynomial-ring submodule from a $t$-stable submodule]
  \label{def:graded_polySubmodule}
  \lean{AlgebraicGeometry.GradedModule.polySubmodule}
  \uses{def:graded_polyModule, lem:graded_polyModule_X_smul,
    lem:graded_polyModule_C_smul}
  Let \(N \subseteq M\) be a \(\kappa\)-submodule that is stable under each \(t_i\)
  (\(t_i \cdot N \subseteq N\)). Then \(N\), with the \emph{same} ambient carrier, is a
  \(\kappa[t_0, \dots, t_{r-1}]\)-submodule of \(M\) for the polynomial-module structure
  (\cref{def:graded_polyModule}). Stability under the action of an arbitrary polynomial
  follows by induction on monomials: constants act as scalars
  (\cref{lem:graded_polyModule_C_smul}) and multiplication by \(X_i\) acts as \(t_i\)
  (\cref{lem:graded_polyModule_X_smul}), both of which preserve \(N\).
\end{definition}

\begin{lemma}\leanok
[The polynomial submodule has the same carrier]
  \label{lem:graded_polySubmodule_coe}
  \lean{AlgebraicGeometry.GradedModule.polySubmodule_coe}
  \uses{def:graded_polySubmodule}
  The underlying set of \(\mathrm{polySubmodule}(t, N)\)
  (\cref{def:graded_polySubmodule}) equals \(N\): the lift to a
  \(\kappa[t_0, \dots, t_{r-1}]\)-submodule does not change the carrier, so every object
  in the induction remains an ambient submodule of the fixed \(M\).
\end{lemma}

\subsubsection*{Finiteness-transfer infrastructure}

The ambient finiteness transfer (\cref{lem:graded_subquotient_finite_transfer})
factors the polynomial action of the larger ring through a surjection onto the
smaller one. The following blocks build that surjection and the semilinearity
identity it rests on, and record the two elementary finiteness manipulations
(shrinking a numerator, enlarging a denominator) used to feed the kernel and
cokernel constructors.

\begin{definition}

\leanok
[The last-variable surjection of polynomial rings]
  \label{lem:graded_lastVarAlgHom}
  \lean{AlgebraicGeometry.GradedModule.lastVarAlgHom}
  \lean{AlgebraicGeometry.GradedModule.lastVarAlgHom_X_castSucc}
  \lean{AlgebraicGeometry.GradedModule.lastVarAlgHom_X_last}
  \lean{AlgebraicGeometry.GradedModule.lastVarAlgHom_C}
  \lean{AlgebraicGeometry.GradedModule.lastVarAlgHom_rename_castSucc}
  \lean{AlgebraicGeometry.GradedModule.lastVarAlgHom_surjective}
  \lean{AlgebraicGeometry.GradedModule.lastVarAlgHom_ringHomSurjective}
  \uses{lem:fg_restrictScalars_of_surjective_mathlib}
  For a field \(\kappa\) and \(r \in \mathbb{N}\), the \emph{last-variable
  homomorphism} is the \(\kappa\)-algebra map
  \[
    \mathrm{lastVarAlgHom} : \kappa[t_0, \dots, t_r]
      = \operatorname{MvPolynomial}(\{0,\dots,r\}, \kappa)
      \;\longrightarrow\;
      \operatorname{MvPolynomial}(\{0,\dots,r-1\}, \kappa) = \kappa[t_0, \dots, t_{r-1}]
  \]
  obtained by evaluation: it sends the last variable \(X_r\) to \(0\) and each
  earlier variable \(X_i\) (for \(0 \le i < r\)) to \(X_i\), and fixes constants.
  It is a left inverse of the algebra map
  \(\kappa[t_0, \dots, t_{r-1}] \to \kappa[t_0, \dots, t_r]\)
  induced by the canonical inclusion \(\{0, \dots, r-1\} \hookrightarrow \{0, \dots, r\}\)
  on variable indices, and is therefore surjective. This is the concrete
  ring surjection \(\kappa[t_0, \dots, t_r] \twoheadrightarrow
  \kappa[t_0, \dots, t_{r-1}]\) along which finiteness is transferred
  (\cref{lem:fg_restrictScalars_of_surjective_mathlib}).
\end{definition}

\begin{lemma}

\leanok
[A $t$-stable submodule is closed under the polynomial action]
  \label{lem:graded_polyEndHom_mem_of_stable}
  \lean{AlgebraicGeometry.GradedModule.polyEndHom_mem_of_stable}
  \uses{def:graded_polyEndHom, lem:graded_polyEndHom_X, lem:graded_polyEndHom_C}
  Let \(t_0, \dots, t_{r-1}\) be pairwise-commuting \(\kappa\)-linear endomorphisms
  of \(M\) and let \(P' \subseteq M\) be a \(\kappa\)-submodule stable under each
  \(t_i\) (\(t_i \cdot P' \subseteq P'\)). Then \(P'\) is closed under the action of
  every polynomial: for all \(p \in \kappa[t_0, \dots, t_{r-1}]\) and \(m \in P'\),
  \(\mathrm{polyEndHom}(t)(p)(m) \in P'\) (\cref{def:graded_polyEndHom}). The proof
  is induction on \(p\): a constant acts as a scalar (\cref{lem:graded_polyEndHom_C}),
  multiplication by a variable acts as the stabilising \(t_i\)
  (\cref{lem:graded_polyEndHom_X}), and sums are handled additively.
\end{lemma}

\begin{lemma}

\leanok
[Mod-$P'$ semilinearity of the polynomial action]
  \label{lem:graded_polyEndHom_lastVar_sub_mem}
  \lean{AlgebraicGeometry.GradedModule.polyEndHom_lastVar_sub_mem}
  \uses{def:graded_polyEndHom, lem:graded_polyEndHom_X, lem:graded_polyEndHom_C,
    lem:graded_polyEndHom_mem_of_stable, lem:graded_lastVarAlgHom}
  Let \(t_0, \dots, t_r\) be pairwise-commuting \(\kappa\)-linear endomorphisms of
  \(M\), and let \(P' \le P\) be \(\kappa\)-submodules each stable under every
  \(t_i\), with the last endomorphism \(x = t_r\) carrying \(P\) into \(P'\)
  (\(x \cdot P \subseteq P'\)). Then for every polynomial
  \(s \in \kappa[t_0, \dots, t_r]\) and every \(m \in P\),
  \[
    \mathrm{polyEndHom}(t)(s)(m)
      \;-\; \mathrm{polyEndHom}(t_0, \dots, t_{r-1})
        \bigl(\mathrm{lastVarAlgHom}(s)\bigr)(m)
      \;\in\; P' .
  \]
  In words: acting by \(s\) through the full family \(t\) agrees, modulo \(P'\), with
  first projecting away the last variable via \(\mathrm{lastVarAlgHom}\)
  (\cref{lem:graded_lastVarAlgHom}) and acting through the reduced family
  \((t_0, \dots, t_{r-1})\). This is the semilinearity identity that makes the
  \(\sigma\)-semilinear quotient map of \cref{lem:graded_subquotient_finite_transfer}
  well-defined. The proof is induction on \(s\): the case of an earlier variable
  \(s = X_i\) (\(i < r\)) is the inductive hypothesis after \(\mathrm{lastVarAlgHom}\)
  fixes it; the case of the last variable \(s = X_r\) annihilates the reduced term,
  and the full term lands in \(P'\) because \(x = t_r\) carries \(P\) into the
  \(t\)-stable \(P'\) (\cref{lem:graded_polyEndHom_mem_of_stable}).
\end{lemma}

\begin{lemma}

\leanok
[Shrinking the numerator preserves finiteness]
  \label{lem:graded_polyQuot_finite_of_le_numerator}
  \lean{AlgebraicGeometry.GradedModule.polyQuot_finite_of_le_numerator}
  \uses{def:graded_polyModule, def:graded_polySubmodule,
    lem:module_finite_of_injective_mathlib,
    lem:mvPolynomial_isNoetherianRing_fin_mathlib,
    lem:isNoetherian_of_isNoetherianRing_of_finite_mathlib}
  Let \(N_1 \le N_2\) and \(P'\) be \(\kappa\)-submodules of \(M\), each stable under
  the commuting family \(t_0, \dots, t_{r-1}\), and suppose the subquotient
  \(N_2/P'\) is finite over \(\kappa[t_0, \dots, t_{r-1}]\)
  (\cref{def:graded_polyModule}, \cref{def:graded_polySubmodule}). Then \(N_1/P'\) is
  also finite over \(\kappa[t_0, \dots, t_{r-1}]\). Indeed the inclusion
  \(N_1 \subseteq N_2\) induces an injective \(\kappa[t_0, \dots, t_{r-1}]\)-linear
  map \(N_1/P' \hookrightarrow N_2/P'\); since
  \(\kappa[t_0, \dots, t_{r-1}]\) is Noetherian
  (\cref{lem:mvPolynomial_isNoetherianRing_fin_mathlib}), the finite module
  \(N_2/P'\) is Noetherian
  (\cref{lem:isNoetherian_of_isNoetherianRing_of_finite_mathlib}), so its submodule
  \(N_1/P'\) is finite (\cref{lem:module_finite_of_injective_mathlib}). This supplies
  the finiteness witness of the kernel constructor
  (\cref{def:graded_subquotientDatum_ker}).
\end{lemma}

\begin{lemma}

\leanok
[Enlarging the denominator preserves finiteness]
  \label{lem:graded_polyQuot_finite_of_le_denominator}
  \lean{AlgebraicGeometry.GradedModule.polyQuot_finite_of_le_denominator}
  \uses{def:graded_polyModule, def:graded_polySubmodule,
    lem:module_finite_of_surjective_mathlib, lem:submodule_liftQ_mathlib}
  Let \(N\) and \(P_1 \le P_2\) be \(\kappa\)-submodules of \(M\), each stable under
  the commuting family \(t_0, \dots, t_{r-1}\), and suppose the subquotient
  \(N/P_1\) is finite over \(\kappa[t_0, \dots, t_{r-1}]\)
  (\cref{def:graded_polyModule}, \cref{def:graded_polySubmodule}). Then \(N/P_2\) is
  also finite over \(\kappa[t_0, \dots, t_{r-1}]\). Indeed the inclusion
  \(P_1 \subseteq P_2\) gives a surjection \(N/P_1 \twoheadrightarrow N/P_2\)
  (\cref{lem:submodule_liftQ_mathlib}), and finiteness transfers along it
  (\cref{lem:module_finite_of_surjective_mathlib}). This supplies the
  finiteness witness of the cokernel constructor
  (\cref{def:graded_subquotientDatum_coker}).
\end{lemma}

\subsubsection*{Ambient finiteness transfer}

\begin{lemma}\leanok
[Abstract single-pair finiteness transfer down one variable]
  \label{lem:graded_subquotient_finite_transfer}
  \lean{AlgebraicGeometry.GradedModule.subquotient_finite_transfer}
  \uses{def:graded_polyModule, def:graded_polySubmodule,
    lem:graded_lastVarAlgHom, lem:graded_polyEndHom_lastVar_sub_mem,
    lem:module_finite_of_surjective_mathlib, lem:submodule_liftQ_mathlib,
    lem:fg_restrictScalars_of_surjective_mathlib}
  Let \(t_0, \dots, t_r\) be pairwise-commuting \(\kappa\)-linear endomorphisms of
  \(M\), and let \(P' \le P\) be \(\kappa\)-submodules of \(M\), each stable under
  every \(t_i\). Suppose the last endomorphism \(x = t_r\) carries \(P\) into \(P'\)
  (\(x \cdot P \subseteq P'\)), and that the subquotient \(P/P'\) is finite over the
  free polynomial ring \(\kappa[t_0, \dots, t_r]\) for the polynomial-module structure
  of \cref{def:graded_polyModule} (with carriers the ambient \(t\)-stable submodules
  \(\mathrm{polySubmodule}\,P\), \(\mathrm{polySubmodule}\,P'\),
  \cref{def:graded_polySubmodule}). Then \(P/P'\) is finite over
  \(\kappa[t_0, \dots, t_{r-1}]\) for the reduced family \((t_0, \dots, t_{r-1})\).

  This is the abstract single-pair transfer: it does not mention the kernel/cokernel
  pairs, but is the engine applied to each of them by the kernel and cokernel
  constructors (\cref{def:graded_subquotientDatum_ker},
  \cref{def:graded_subquotientDatum_coker}). Because the last variable
  annihilates the subquotient, the \(\kappa[t_0, \dots, t_r]\)-action factors through
  the last-variable surjection \(\kappa[t_0, \dots, t_r] \twoheadrightarrow
  \kappa[t_0, \dots, t_{r-1}]\), which is what makes the transfer drop one variable.
\end{lemma}

\begin{proof}
  \uses{def:graded_polyModule, def:graded_polySubmodule,
    lem:graded_lastVarAlgHom, lem:graded_polyEndHom_lastVar_sub_mem,
    lem:module_finite_of_surjective_mathlib, lem:submodule_liftQ_mathlib}
    \leanok
  Write \(\sigma = \mathrm{lastVarAlgHom} : \kappa[t_0, \dots, t_r]
  \twoheadrightarrow \kappa[t_0, \dots, t_{r-1}]\) for the surjective ring
  homomorphism sending the last variable to \(0\) and each earlier variable \(X_i\)
  (\(i < r\)) to \(X_i\) (\cref{lem:graded_lastVarAlgHom}). Consider the map on numerators
  \[
    g : \mathrm{polySubmodule}\,P \;\longrightarrow\;
      \bigl(\mathrm{polySubmodule}\,P\bigr)\big/
        \bigl(\mathrm{polySubmodule}\,P'\bigr), \qquad y \longmapsto [\,y\,],
  \]
  where the source carries the \(\kappa[t_0, \dots, t_r]\)-module structure and the
  target the reduced \(\kappa[t_0, \dots, t_{r-1}]\)-module structure (both with the
  same ambient carrier \(P\), \cref{def:graded_polySubmodule}). The mod-\(P'\)
  semilinearity identity (\cref{lem:graded_polyEndHom_lastVar_sub_mem}) -- acting by
  \(s\) through the full family and acting by \(\sigma(s)\) through the reduced family
  agree modulo \(P'\), precisely because \(x = t_r\) carries \(P\) into \(P'\) --
  says exactly that \(g\) is \(\sigma\)-semilinear. The map \(g\) is surjective (it is
  the quotient projection on the common carrier), and it kills the denominator
  \(\mathrm{polySubmodule}\,P'\), so it descends through that submodule
  (\cref{lem:submodule_liftQ_mathlib}) to a surjective
  \(\kappa[t_0, \dots, t_{r-1}]\)-linear map
  \[
    \bigl(\mathrm{polySubmodule}\,P\bigr)\big/\bigl(\mathrm{polySubmodule}\,P'\bigr)
      \;\twoheadrightarrow\;
    \bigl(\mathrm{polySubmodule}\,P\bigr)\big/\bigl(\mathrm{polySubmodule}\,P'\bigr)
  \]
  from the \(\kappa[t_0, \dots, t_r]\)-finite subquotient onto the reduced one.
  Finiteness transfers along this surjection
  (\cref{lem:module_finite_of_surjective_mathlib}), giving finiteness of \(P/P'\)
  over \(\kappa[t_0, \dots, t_{r-1}]\).

  \emph{Why the free polynomial ring.} The smaller ring must be the \emph{free}
  polynomial ring \(\kappa[t_0, \dots, t_{r-1}]\), not the subalgebra
  \(\mathrm{adjoin}_\kappa\{t_0, \dots, t_{r-1}\} \subseteq
  \operatorname{End}_\kappa(M)\): relations among the endomorphisms inside
  \(\operatorname{End}_\kappa(M)\) would obstruct the polynomial-ring surjection
  \(\sigma\) that the scalar transfer requires
  (\cref{lem:fg_restrictScalars_of_surjective_mathlib}). Keeping free polynomial
  rings at every step keeps \(\sigma\) -- and hence the transfer -- available
  throughout. Everything is finite generation of ambient submodules of the fixed
  \(M\) and their quotients.
\end{proof}

\subsubsection*{Subquotient constructors}

The kernel/cokernel closure (\cref{lem:graded_subquotient_ker_coker}) and the
abstract finiteness transfer (\cref{lem:graded_subquotient_finite_transfer}) are
bundled into two constructors that take a length-\((r+1)\) subquotient datum to the
length-\(r\) kernel and cokernel data consumed by the induction. The two stability
lemmas below first record that the new pairs are stable under the \emph{whole}
endomorphism family (not merely the reduced one), which is what lets the over-\(\kappa[t_0,\dots,t_r]\)
finiteness inputs be stated before the variable is dropped.

\begin{lemma}

\leanok
[The kernel pair's lower module is stable under the full family]
  \label{lem:graded_ker_stable_full}
  \lean{AlgebraicGeometry.GradedModule.ker_stable_full}
  \uses{lem:graded_comap_map_le_of_commute}
  Let \((N, N')\) be a length-\((r+1)\) subquotient datum
  (\cref{def:graded_subquotientHilb}) with endomorphisms \(t_0, \dots, t_r\) and
  \(x = t_r\). Then the kernel pair's lower module \(N \cap x^{-1}(N')\) is stable
  under \emph{every} \(t_i\) (\(0 \le i \le r\)), not only the reduced family: each
  \(t_i\) preserves \(N\) and preserves the preimage \(x^{-1}(N')\) because it
  commutes with \(x\) and preserves \(N'\)
  (\cref{lem:graded_comap_map_le_of_commute}). This full stability is what allows the
  kernel subquotient's finiteness to be measured over \(\kappa[t_0, \dots, t_r]\)
  before the transfer drops the last variable.
\end{lemma}

\begin{lemma}

\leanok
[The cokernel pair's upper module is stable under the full family]
  \label{lem:graded_coker_stable_full}
  \lean{AlgebraicGeometry.GradedModule.coker_stable_full}
  \uses{lem:graded_map_map_le_of_commute}
  Under the same hypotheses, the cokernel pair's upper module \(N' + x \cdot N\) is
  stable under every \(t_i\) (\(0 \le i \le r\)): each \(t_i\) preserves \(N'\) and
  preserves the image \(x \cdot N\) because it commutes with \(x\) and preserves
  \(N\) (\cref{lem:graded_map_map_le_of_commute}), and a map preserves a sum of
  preserved submodules. This full stability supports the over-\(\kappa[t_0, \dots,
  t_r]\) finiteness input of the cokernel constructor.
\end{lemma}

\begin{definition}

\leanok
[Kernel constructor of a subquotient datum]
  \label{def:graded_subquotientDatum_ker}
  \lean{AlgebraicGeometry.GradedModule.SubquotientDatum.ker}
  \uses{def:graded_subquotientHilb, lem:graded_subquotient_ker_coker,
    lem:graded_ker_isHomogeneous, lem:graded_ker_le, lem:graded_ker_annihilate,
    lem:graded_comap_map_le_of_commute, lem:graded_ker_stable_full,
    lem:graded_subquotient_finite_transfer,
    lem:graded_polyQuot_finite_of_le_numerator}
  From a length-\((r+1)\) subquotient datum \((N, N')\) with endomorphisms
  \(t_0, \dots, t_r\) and \(x = t_r\) (\cref{def:graded_subquotientHilb}), the
  \emph{kernel constructor} produces the length-\(r\) datum
  \[
    \ker D \;=\; \bigl(\,N \cap x^{-1}(N'),\ N'\,\bigr)
  \]
  on the reduced family \((t_0, \dots, t_{r-1})\). Its
  non-finiteness fields are supplied by the ambient kernel calculus: homogeneity of
  \(N \cap x^{-1}(N')\) (\cref{lem:graded_ker_isHomogeneous}), the nesting
  \(N' \le N \cap x^{-1}(N')\) (\cref{lem:graded_ker_le}), and stability under the
  reduced family (\cref{lem:graded_ker_stable_full}). The finiteness witness
  is obtained by first shrinking the numerator from \(N\) to
  \(N \cap x^{-1}(N')\) (\cref{lem:graded_polyQuot_finite_of_le_numerator}) to get
  \(\kappa[t_0, \dots, t_r]\)-finiteness of the kernel subquotient, then transferring
  down one variable (\cref{lem:graded_subquotient_finite_transfer}), valid because
  \(x\) annihilates the kernel subquotient (\cref{lem:graded_ker_annihilate}). This
  realises the kernel half of \cref{lem:graded_subquotient_ker_coker} as a genuine
  length-\(r\) datum.
\end{definition}

\begin{definition}

\leanok
[Cokernel constructor of a subquotient datum]
  \label{def:graded_subquotientDatum_coker}
  \lean{AlgebraicGeometry.GradedModule.SubquotientDatum.coker}
  \uses{def:graded_subquotientHilb, lem:graded_subquotient_ker_coker,
    lem:graded_coker_isHomogeneous, lem:graded_coker_le, lem:graded_coker_annihilate,
    lem:graded_map_map_le_of_commute, lem:graded_coker_stable_full,
    lem:graded_subquotient_finite_transfer,
    lem:graded_polyQuot_finite_of_le_denominator}
  From the same length-\((r+1)\) datum, the \emph{cokernel constructor} produces the
  length-\(r\) datum
  \[
    \operatorname{coker} D \;=\; \bigl(\,N,\ N' + x \cdot N\,\bigr)
  \]
  on the reduced family. Its non-finiteness fields come from the ambient cokernel
  calculus: homogeneity of \(N' + x \cdot N\) (\cref{lem:graded_coker_isHomogeneous}),
  the nesting \(N' + x \cdot N \le N\) (\cref{lem:graded_coker_le}), and stability
  under the reduced family (\cref{lem:graded_coker_stable_full}). The finiteness field
  is obtained by enlarging the denominator from \(N'\) to \(N' + x \cdot N\)
  (\cref{lem:graded_polyQuot_finite_of_le_denominator}) to get
  \(\kappa[t_0, \dots, t_r]\)-finiteness of the cokernel subquotient, then
  transferring down one variable (\cref{lem:graded_subquotient_finite_transfer}),
  valid because \(x\) annihilates the cokernel subquotient
  (\cref{lem:graded_coker_annihilate}). This realises the cokernel half of
  \cref{lem:graded_subquotient_ker_coker} as a genuine length-\(r\) datum.
\end{definition}

\subsubsection*{The degreewise difference identity}

\begin{lemma}\leanok
[$D_5$ -- degreewise rank--nullity difference]
  \label{lem:graded_degreewise_finrank_diff}
  \lean{AlgebraicGeometry.GradedModule.degreewise_finrank_diff}
  \uses{lem:finrank_range_add_finrank_ker_mathlib,
    lem:finrank_ses_additive_mathlib}
  Let \(\varphi : V \to W\) be a \(\kappa\)-linear map between finite-dimensional
  \(\kappa\)-vector spaces. Then
  \[
    \dim_\kappa W - \dim_\kappa V
      \;=\; \dim_\kappa(W/\operatorname{range}\varphi) - \dim_\kappa(\ker\varphi).
  \]
  Applied degreewise with \(\varphi = (x : \mathcal{M}_n \to \mathcal{M}_{n+1})\),
  \(V = \mathcal{M}_n\), \(W = \mathcal{M}_{n+1}\) -- so that \(\ker\varphi\) is the
  degree-\(n\) piece \(K_n\) of the kernel subquotient and
  \(W/\operatorname{range}\varphi = \mathcal{M}_{n+1}/x \mathcal{M}_n\) the
  degree-\((n+1)\) piece \(C_{n+1}\) of the cokernel subquotient
  (\cref{lem:graded_subquotient_degreewise_diff}, via the split \(G_1\)
  \cref{lem:graded_homogeneousSubmodule_iSup_eq} and the degree shift
  \cref{lem:isHomogeneousElem_graded_smul_mathlib}) -- this gives, for every \(n\),
  \[
    \dim_\kappa \mathcal{M}_{n+1} - \dim_\kappa \mathcal{M}_n
      \;=\; \dim_\kappa C_{n+1} - \dim_\kappa K_n,
  \]
  which is the difference identity \(\mathrm{hdiff}\) consumed by
  \cref{lem:ratHilb_ofDiffEq}. No graded structure is used in proving the
  underlying vector-space identity -- it is pure rank--nullity.
\end{lemma}

\begin{proof}
  \uses{lem:finrank_range_add_finrank_ker_mathlib,
    lem:finrank_ses_additive_mathlib}
    \leanok
  By rank--nullity (\cref{lem:finrank_range_add_finrank_ker_mathlib}),
  \(\dim_\kappa V = \dim_\kappa(\ker\varphi) + \dim_\kappa(\operatorname{range}
  \varphi)\), so \(\dim_\kappa(\operatorname{range}\varphi) = \dim_\kappa V -
  \dim_\kappa(\ker\varphi)\). By additivity on the short exact sequence
  \(0 \to \operatorname{range}\varphi \to W \to W/\operatorname{range}\varphi \to 0\)
  (\cref{lem:finrank_ses_additive_mathlib}),
  \(\dim_\kappa(W/\operatorname{range}\varphi) = \dim_\kappa W -
  \dim_\kappa(\operatorname{range}\varphi) = \dim_\kappa W - \dim_\kappa V +
  \dim_\kappa(\ker\varphi)\). Rearranging gives
  \(\dim_\kappa W - \dim_\kappa V = \dim_\kappa(W/\operatorname{range}\varphi) -
  \dim_\kappa(\ker\varphi)\). The degreewise application identifies
  \(\ker(x : \mathcal{M}_n \to \mathcal{M}_{n+1})\) and
  \(\mathcal{M}_{n+1}/\operatorname{range}(x : \mathcal{M}_n \to \mathcal{M}_{n+1})\)
  with the degree-\(n\) piece of the kernel subquotient and the degree-\((n+1)\)
  piece of the cokernel subquotient (\cref{lem:graded_subquotient_degreewise_diff}),
  the latter because the degree-\((n+1)\) graded piece of the homogeneous submodule
  \(x\mathcal{M}\) is exactly \(x \mathcal{M}_n\)
  (\cref{lem:isHomogeneousElem_graded_smul_mathlib}).
\end{proof}

\subsubsection*{The base case ($r = 0$)}

\begin{lemma}\leanok
[Base-case finiteness over the empty polynomial ring]
  \label{lem:graded_finiteDimensional_of_mvPolynomial_isEmpty_finite}
  \lean{AlgebraicGeometry.GradedModule.finiteDimensional_of_mvPolynomial_isEmpty_finite}
  Let \(Q\) be a \(\kappa\)-vector space that is also a module over
  \(\operatorname{MvPolynomial}(\sigma, \kappa)\) for an \emph{empty} index set
  \(\sigma\) -- in particular \(\sigma = \{0, \dots, r-1\}\) with \(r = 0\) -- in a way
  compatible with the \(\kappa\)-action (scalar tower). If \(Q\) is module-finite over
  \(\operatorname{MvPolynomial}(\sigma, \kappa)\), then \(Q\) is finite-dimensional over
  \(\kappa\). Indeed for empty \(\sigma\) the polynomial ring is
  \(\operatorname{MvPolynomial}(\sigma, \kappa) \cong \kappa\), a finite \(\kappa\)-module,
  so finiteness of \(Q\) over it transfers to finiteness over \(\kappa\) along the scalar
  tower. This is the length-zero input to the ambient subquotient induction
  (\cref{lem:graded_subquotient_isRatHilb}): when \(r = 0\) the subquotient \(N/N'\) is a
  finite-dimensional \(\kappa\)-vector space, hence \(\mathrm{hilb}\) is eventually zero.
\end{lemma}

\begin{lemma}

\leanok
[independence of images modulo the kernel]
  \label{lem:graded_iSupIndep_map_of_mem_ker_sup}
  \lean{AlgebraicGeometry.GradedModule.iSupIndep_map_of_mem_ker_sup}
  Let \(\pi : A \to Q\) be a \(\kappa\)-linear map and \((g_i)_{i \in \iota}\) a family of
  subspaces of \(A\) that is \emph{independent modulo \(\ker\pi\)}: for each \(i\), every
  \(a \in g_i\) that also lies in \(\ker\pi + \sum_{j \ne i} g_j\) already lies in
  \(\ker\pi\). Then the image family \((\pi(g_i))_{i \in \iota}\) is independent
  (\(\mathrm{iSupIndep}\)) in \(Q\). Ring-agnostic lattice core of the base-case
  independence step below; applied to the quotient map \(N \twoheadrightarrow N/N'\).
\end{lemma}
\begin{proof}

\leanok
  By the disjointness criterion for independence it suffices to show, for each \(i\),
  that \(\pi(g_i) \cap \sum_{j \ne i} \pi(g_j) = 0\). Take \(q\) in that intersection;
  write \(q = \pi(a)\) with \(a \in g_i\), and (pushing the sum through \(\pi\), since
  \(\pi(\sum_{j\ne i} g_j) = \sum_{j\ne i}\pi(g_j)\)) \(q = \pi(b)\) with
  \(b \in \sum_{j \ne i} g_j\). Then \(a - b \in \ker\pi\), so
  \(a = (a-b) + b \in \ker\pi + \sum_{j \ne i} g_j\). By the hypothesis applied at \(i\),
  \(a \in \ker\pi\), whence \(q = \pi(a) = 0\).
\end{proof}

\begin{lemma}

\leanok
[Base case: the length-zero subquotient Hilbert function vanishes eventually]
  \label{lem:graded_subquotient_base_eventuallyZero}
  \lean{AlgebraicGeometry.GradedModule.subquotient_base_eventuallyZero}
  \uses{def:graded_subquotientHilb,
    lem:graded_finiteDimensional_of_mvPolynomial_isEmpty_finite,
    lem:graded_homogeneousSubmodule_iSupIndep, lem:graded_iSupIndep_map_of_mem_ker_sup,
    lem:finite_ne_bot_of_iSupIndep_mathlib, lem:ratHilb_ofEventuallyZero}
  % SOURCE: [Stacks Project], "Commutative Algebra", \S "Noetherian graded
  % rings", Proposition \texttt{proposition-graded-hilbert-polynomial} (Tag
  % 00K1), proof, base case (read from references/hilbert-serre-algebra.tex,
  % L13908-L13910).
  % SOURCE QUOTE:
  % "We prove this by induction on the minimal number of
  %  generators of $S_1$. If this number is $0$, then
  %  $M_n = 0$ for all $n \gg 0$ and the result holds."
  \textit{Source: [Stacks Project], Tag 00K1 (base case of the induction).}
  Let \((N, N')\) be a subquotient datum of length \(0\)
  (\cref{def:graded_subquotientHilb}): a homogeneous pair \(N' \le N\) of the fixed
  ambient graded module \(M = \bigoplus_n \mathcal{M}_n\) with \emph{no} endomorphisms,
  whose finiteness witness says the subquotient \(Q_0 := N/N'\) is module-finite over
  \(\operatorname{MvPolynomial}(\varnothing, \kappa)\). Then the ambient Hilbert
  function vanishes eventually: there is \(K \in \mathbb{N}\) with
  \(\mathrm{hilb}(n) = 0\) for all \(n > K\). Consequently
  \(\mathrm{IsRatHilb}(\mathrm{hilb}, 0)\) (\cref{lem:ratHilb_ofEventuallyZero}),
  which is the base case of the ambient subquotient induction
  (\cref{lem:graded_subquotient_isRatHilb}).
\end{lemma}

\begin{proof}
  \uses{lem:graded_finiteDimensional_of_mvPolynomial_isEmpty_finite,
    lem:graded_homogeneousSubmodule_iSupIndep,
    lem:finite_ne_bot_of_iSupIndep_mathlib, lem:ratHilb_ofEventuallyZero}
    \leanok
  Since the index set of variables is empty,
  \(\operatorname{MvPolynomial}(\varnothing, \kappa) \cong \kappa\), so the finiteness
  witness makes the subquotient \(Q_0 = N/N'\) a \emph{finite-dimensional}
  \(\kappa\)-vector space
  (\cref{lem:graded_finiteDimensional_of_mvPolynomial_isEmpty_finite}). For each
  degree \(n\), composing the inclusion of the homogeneous piece
  \(N \cap \mathcal{M}_n \hookrightarrow N\) with the quotient projection
  \(N \twoheadrightarrow Q_0\) gives a \(\kappa\)-linear map
  \(\psi_n : N \cap \mathcal{M}_n \to Q_0\), \(m \mapsto [m]\), whose image is the
  class of the degree-\(n\) homogeneous piece inside the subquotient.

  \emph{Step 1: independence of the images.} The family
  \(\bigl(\operatorname{range}(\psi_n)\bigr)_{n}\) is \(\mathrm{iSupIndep}\) in
  \(Q_0\). The reason is the ambient direct-sum grading of \(M\): the homogeneous
  pieces \(\mathcal{M}_n\) form an independent family
  (\cref{lem:graded_homogeneousSubmodule_iSupIndep}). An element of
  \(\bigsqcup_{j \ne n} \operatorname{range}(\psi_j)\) is a finite sum of classes of
  homogeneous elements of degrees \(j \ne n\); since \(N'\) is homogeneous, the
  degree-\(n\) homogeneous component of any representative of such a class lies in
  \(N'\), so the class has zero degree-\(n\) component in \(Q_0\). A nonzero class in
  \(\operatorname{range}(\psi_n)\), by contrast, is represented by a homogeneous
  element of degree \(n\) not in \(N'\), so has nonzero degree-\(n\) component.
  Comparing degree-\(n\) components forces the intersection
  \(\operatorname{range}(\psi_n) \cap \bigsqcup_{j \ne n} \operatorname{range}(\psi_j)
  = 0\); independence follows.
  % NOTE: ROUTE (b). Prove the iSupIndep by direct membership analysis in
  %   ⨆_{j ≠ n} range(ψ j), reading off the degree-n homogeneous component
  %   directly inside Q₀'s fixed κ-module structure.
  %   No outgoing linear map out of Q₀ is constructed, so no scalar ring other than κ
  %   is ever introduced.
  % NOTE: ROUTE (a). Building a κ-linear detector
  %   Φ : Q₀ → M/N', [m] ↦ [π_n m] (degree-n projection) via descent through the
  %   subquotient fails: the descended map carries the scalar action of Q₀'s base
  %   ring, but the target M/N' carries only a κ-module structure -- a scalar-ring
  %   incompatibility. Route (b) avoids this by working entirely within Q₀.

  \emph{Step 2: finiteness of the support.} Because \(Q_0\) is finite-dimensional and
  the ranges \(\operatorname{range}(\psi_n)\) are independent, only finitely many of
  them are nonzero: the set \(\{\,n : \operatorname{range}(\psi_n) \ne 0\,\}\) is
  finite (\cref{lem:finite_ne_bot_of_iSupIndep_mathlib}). In particular it is bounded
  above, by some \(K \in \mathbb{N}\).

  \emph{Step 3: eventual vanishing.} For \(n > K\) we have
  \(\operatorname{range}(\psi_n) = 0\), i.e.\ every homogeneous element of
  \(N \cap \mathcal{M}_n\) has class \(0\) in \(Q_0 = N/N'\); equivalently
  \(N \cap \mathcal{M}_n \le N'\), hence \(N \cap \mathcal{M}_n = N' \cap \mathcal{M}_n\)
  and
  \[
    \mathrm{hilb}(n) \;=\; \dim_\kappa(N \cap \mathcal{M}_n)
      - \dim_\kappa(N' \cap \mathcal{M}_n) \;=\; 0 .
  \]
  Thus \(\mathrm{hilb}\) is eventually zero, and
  \cref{lem:ratHilb_ofEventuallyZero} gives \(\mathrm{IsRatHilb}(\mathrm{hilb}, 0)\).
\end{proof}

\subsubsection*{The ambient subquotient induction}

\begin{lemma}\leanok
[Rationality of the ambient subquotient Hilbert function]
  \label{lem:graded_subquotient_isRatHilb}
  \lean{AlgebraicGeometry.GradedModule.subquotient_hilbertSeries_rational}
  \uses{def:graded_subquotientHilb, lem:graded_subquotient_ker_coker,
    def:graded_subquotientDatum_ker, def:graded_subquotientDatum_coker,
    lem:graded_subquotient_degreewise_diff,
    lem:graded_subquotient_finite_transfer, lem:ratHilb_ofEventuallyZero,
    lem:ratHilb_ofDiffEq, lem:ratHilb_bump,
    lem:graded_subquotient_base_eventuallyZero,
    lem:graded_finiteDimensional_of_mvPolynomial_isEmpty_finite}
  % SOURCE: [Stacks Project], "Commutative Algebra", \S "Noetherian graded
  % rings", Proposition \texttt{proposition-graded-hilbert-polynomial} (Tag
  % 00K1), proof (read from references/hilbert-serre-algebra.tex, L13908-L13914).
  \textit{Source: [Stacks Project], Tag 00K1 (induction on the number of
  degree-one generators).}
  Let \((N, N')\) be a subquotient datum of length \(r\)
  (\cref{def:graded_subquotientHilb}) with ambient Hilbert function
  \(\mathrm{hilb}\). Then \(\mathrm{IsRatHilb}(\mathrm{hilb}, r)\)
  (\cref{def:ratHilb}): the difference \(n \mapsto \dim_\kappa(N \cap \mathcal{M}_n)
  - \dim_\kappa(N' \cap \mathcal{M}_n)\) is, for \(n \gg 0\), the coefficient
  sequence of a rational series \(p \cdot (1 - X)^{-r}\) with \(p \in
  \mathbb{Q}[X]\). The induction is on the length \(r\), and every object that
  appears is an ambient pair of homogeneous submodules of the fixed graded module
  \(M\).
\end{lemma}

% SOURCE QUOTE PROOF: [Stacks Project], Tag 00K1, proof, induction set-up (read
% from references/hilbert-serre-algebra.tex, L13908-L13914).
% "We prove this by induction on the minimal number of
%  generators of $S_1$. If this number is $0$, then
%  $M_n = 0$ for all $n \gg 0$ and the result holds.
%  To prove the induction step, let $x\in S_1$
%  be one of a minimal set of generators, such that
%  the induction hypothesis applies to the
%  graded ring $S/(x)$."
\begin{proof}
  \uses{lem:graded_subquotient_ker_coker,
    def:graded_subquotientDatum_ker, def:graded_subquotientDatum_coker,
    lem:graded_subquotient_degreewise_diff,
    lem:graded_subquotient_finite_transfer, lem:ratHilb_ofEventuallyZero,
    lem:ratHilb_ofDiffEq, lem:ratHilb_bump,
    lem:graded_subquotient_base_eventuallyZero,
    lem:graded_finiteDimensional_of_mvPolynomial_isEmpty_finite}
    \leanok
  We induct on the length \(r\).

  \emph{Base case \(r = 0\).} A length-zero datum has a finiteness witness making the
  subquotient \(N/N'\) a finite-dimensional \(\kappa\)-vector space
  (\cref{lem:graded_finiteDimensional_of_mvPolynomial_isEmpty_finite}); its degreewise
  images form an independent family inside that finite-dimensional space, so only
  finitely many are nonzero and \(\mathrm{hilb}(n) = 0\) for \(n \gg 0\)
  (\cref{lem:graded_subquotient_base_eventuallyZero}). By
  \cref{lem:ratHilb_ofEventuallyZero} an eventually-zero function is rational of order
  \(0\), giving \(\mathrm{IsRatHilb}(\mathrm{hilb}, 0)\).

  \emph{Inductive step \(r \ge 1\).} Set \(x = t_{r-1}\). By the kernel/cokernel
  closure (\cref{lem:graded_subquotient_ker_coker}) the kernel subquotient
  \(K = (N \cap x^{-1}(N'),\, N')\) and the cokernel subquotient
  \(C = (N,\, N' + x \cdot N)\) are subquotient data of length \(r-1\), assembled by
  the kernel and cokernel constructors (\cref{def:graded_subquotientDatum_ker},
  \cref{def:graded_subquotientDatum_coker}) whose finiteness witnesses come from the
  abstract finiteness transfer (\cref{lem:graded_subquotient_finite_transfer}). The
  induction hypothesis applied
  to \(K\) and \(C\) gives \(\mathrm{IsRatHilb}(h_K, r-1)\) and
  \(\mathrm{IsRatHilb}(h_C, r-1)\) (\cref{def:ratHilb}). By the degreewise
  difference identity (\cref{lem:graded_subquotient_degreewise_diff}),
  \[
    \mathrm{hilb}(n+1) - \mathrm{hilb}(n) \;=\; h_C(n+1) - h_K(n)
    \qquad \text{for all } n .
  \]
  Feeding \(\mathrm{IsRatHilb}(h_C, r-1)\), \(\mathrm{IsRatHilb}(h_K, r-1)\) and this
  difference identity into the inductive-step engine
  (\cref{lem:ratHilb_ofDiffEq}) yields \(\mathrm{IsRatHilb}(\mathrm{hilb}, r)\). (The
  common order \(r-1\) for \(h_C\) and \(h_K\) is arranged, if necessary, by raising
  the lower one with \cref{lem:ratHilb_bump}.) This completes the induction.
\end{proof}

\section{Support and freeness predicates}
\label{sec:quot_predicates}

The Quot functor (\cref{def:quot_functor}) and the Grassmannian
(\cref{def:grassmannian_scheme}) classify quotient sheaves subject to two
geometric conditions: a coherent quotient must have \emph{schematic support
proper over the base}, and the Grassmannian's quotients must be \emph{locally
free of a fixed rank}. This section introduces both predicates.

\begin{definition}

\leanok
[Annihilator ideal sheaf of a sheaf of modules]
  \label{def:modules_annihilator}
  \lean{AlgebraicGeometry.Scheme.Modules.annihilator}
  % NOTE: The definition carries no \uses: the Lean construction is
  % `IdealSheafData.ofIdeals (fun U => Module.annihilator Γ(X,U) Γ(F,U))`, which is well-defined
  % unconditionally and does NOT depend on the localization bridges at definition time. The two
  % facts lem:annihilator_localization_eq_map and lem:qcoh_section_localization_basicOpen are needed
  % only for the downstream CHARACTERIZATION (the reverse inclusion / equality on affine opens, the
  % "annihilator forward direction"), which is not yet a blueprint block.
  % NOTE: The definition itself is formalized and axiom-clean; the inclusion direction
  % \(\mathrm{Ann}(\mathcal{F})(U) \subseteq \mathrm{Ann}_{\mathcal{O}_X(U)}(\mathcal{F}(U))\)
  % (\cref{lem:modules_annihilator_ideal_le}) is closed. The reverse inclusion, which identifies
  % the annihilator sheaf with the module-annihilator on affine opens, remains bridge-gated on
  % \cref{lem:qcoh_section_localization_basicOpen}.
  % SOURCE: [Nitsure], \S 1, sub-section "Stratification by Hilbert
  % Polynomials" (read from
  % references/nitsure-hilbert-quot-src/nitsure-hilbert-quot.tex, L468-L471).
  % SOURCE QUOTE:
  % "Let $X\to S$ be a finite type morphism of noetherian schemes,
  %  and let $L$ be a line bundle on $X$. Let $\F$
  %  be any coherent sheaf on
  %  $X$ whose schematic support is proper over $S$."
  \textit{Source: [Nitsure], \S 1 (FGA Explained Ch.~5).}
  Let \(X\) be a scheme and let \(\mathcal{F}\) be a sheaf of
  \(\mathcal{O}_X\)-modules that is \emph{quasi-coherent and of finite type}
  (over a locally noetherian base, so that on each affine open \(U\) the section
  module \(\mathcal{F}(U)\) is a finitely generated \(\mathcal{O}_X(U)\)-module).
  The \emph{annihilator ideal sheaf} of \(\mathcal{F}\) is the ideal sheaf
  \(\mathrm{Ann}(\mathcal{F})\) defined on affine opens: for each affine open
  \(U \subseteq X\),
  \[
    \mathrm{Ann}(\mathcal{F})(U)
    \;=\; \mathrm{Ann}_{\mathcal{O}_X(U)}\bigl(\mathcal{F}(U)\bigr)
    \;\subseteq\; \mathcal{O}_X(U),
  \]
  the annihilator of the \(\mathcal{O}_X(U)\)-module \(\mathcal{F}(U)\). These
  local data patch to a coherent ideal sheaf on \(X\) whose vanishing locus is
  the set-theoretic support of \(\mathcal{F}\).

  \emph{Well-definedness on basic opens.} The coherence condition for an ideal
  sheaf requires that, for an affine open \(U\) and \(f \in \mathcal{O}_X(U)\),
  the section \(\mathrm{Ann}(\mathcal{F})(U)\) restricts correctly to the basic
  open \(D(f) \subseteq U\), i.e. that
  \(\mathrm{Ann}_{\mathcal{O}_X(D(f))}(\mathcal{F}(D(f)))\) is the extension of
  \(\mathrm{Ann}_{\mathcal{O}_X(U)}(\mathcal{F}(U))\) along the localization map
  \(\mathcal{O}_X(U) \to \mathcal{O}_X(D(f))\). This is exactly where the two
  hypotheses enter. The bridge
  \cref{lem:qcoh_section_localization_basicOpen} supplies, for a quasi-coherent
  finite-type \(\mathcal{F}\), that \(\mathcal{O}_X(D(f)) = (\mathcal{O}_X(U))_f\)
  and that the restriction \(\mathcal{F}(U) \to \mathcal{F}(D(f))\) exhibits
  \(\mathcal{F}(D(f))\) as the localized module \((\mathrm{powers}\,f)^{-1}
  \mathcal{F}(U)\); the algebra engine
  \cref{lem:annihilator_localization_eq_map} then yields, using that
  \(\mathcal{F}(U)\) is finitely generated, the required identity
  \(\mathrm{Ann}_{(\mathcal{O}_X(U))_f}\bigl((\mathrm{powers}\,f)^{-1}
  \mathcal{F}(U)\bigr) = \mathrm{Ann}_{\mathcal{O}_X(U)}(\mathcal{F}(U)) \cdot
  (\mathcal{O}_X(U))_f\). Thus the local annihilators are compatible with basic-open
  restriction, which is the patching datum that makes \(\mathrm{Ann}(\mathcal{F})\)
  a well-defined ideal sheaf. The finite-type hypothesis is genuinely required:
  without finite generation of \(\mathcal{F}(U)\) the annihilator need not commute
  with localization, and the local data fail to patch.

  \emph{Note.} Mathlib does not carry an annihilator for sheaves of modules on a scheme.
\end{definition}

\begin{lemma}\leanok
[Annihilator ideal bounded by the module annihilator on affine opens]
  \label{lem:modules_annihilator_ideal_le}
  \lean{AlgebraicGeometry.Scheme.Modules.annihilator_ideal_le}
  \uses{def:modules_annihilator}
  Let \(X\) be a scheme, \(\mathcal{F}\) a sheaf of \(\mathcal{O}_X\)-modules and
  \(U \subseteq X\) an affine open. Then the section over \(U\) of the annihilator
  ideal sheaf \(\mathrm{Ann}(\mathcal{F})\) (\cref{def:modules_annihilator}) is
  contained in the module-annihilator of the sections of \(\mathcal{F}\) over
  \(U\):
  \[
    \mathrm{Ann}(\mathcal{F})(U)
    \;\subseteq\; \mathrm{Ann}_{\mathcal{O}_X(U)}\bigl(\mathcal{F}(U)\bigr).
  \]
  This is the inclusion direction available directly from the definition of
  \(\mathrm{Ann}(\mathcal{F})\); the reverse inclusion is the basic-open coherence
  statement gated on \cref{lem:qcoh_section_localization_basicOpen}. Proved directly
  in Lean. It supports the well-definedness of the schematic support
  (\cref{def:schematic_support}).
\end{lemma}

\begin{lemma}\leanok
[Annihilator family is basic-open coherent on affine opens]
  \label{lem:annihilator_map_basicOpen}
  \lean{AlgebraicGeometry.Scheme.Modules.annihilator_map_basicOpen}
  \uses{def:modules_annihilator, lem:annihilator_localization_eq_map,
    lem:qcoh_section_localization_basicOpen}
  Let \(X\) be a scheme, \(\mathcal{F}\) a quasi-coherent sheaf of \(\mathcal{O}_X\)-modules,
  \(V \subseteq X\) an affine open, and \(f \in \mathcal{O}_X(V)\), and suppose the section
  module \(\mathcal{F}(V)\) is \emph{finitely generated} over \(\mathcal{O}_X(V)\). Then the
  module-annihilator family is coherent across the basic open \(D(f) \subseteq V\): the
  extension of \(\mathrm{Ann}_{\mathcal{O}_X(V)}(\mathcal{F}(V))\) along the restriction map
  \(\mathcal{O}_X(V) \to \mathcal{O}_X(D(f))\) equals the module-annihilator over \(D(f)\),
  \[
    \bigl(\mathrm{Ann}_{\mathcal{O}_X(V)}(\mathcal{F}(V))\bigr)\cdot \mathcal{O}_X(D(f))
    \;=\; \mathrm{Ann}_{\mathcal{O}_X(D(f))}\bigl(\mathcal{F}(D(f))\bigr).
  \]
  This is exactly the \(\mathrm{map\_ideal\_basicOpen}\) coherence field needed to assemble the
  annihilator local data into an \(\mathrm{IdealSheafData}\) (\cref{lem:modules_annihilator_ideal}).
\end{lemma}
\begin{proof}
  \uses{lem:annihilator_localization_eq_map, lem:qcoh_section_localization_basicOpen}
  Since \(V\) is affine, \(\mathcal{O}_X(D(f)) = (\mathcal{O}_X(V))_f\) is the localization at the
  powers of \(f\), and by \cref{lem:qcoh_section_localization_basicOpen} the restriction
  \(\mathcal{F}(V) \to \mathcal{F}(D(f))\) exhibits \(\mathcal{F}(D(f))\) as the localized module
  \((\mathrm{powers}\,f)^{-1}\mathcal{F}(V)\). As \(\mathcal{F}(V)\) is finitely generated,
  \cref{lem:annihilator_localization_eq_map} identifies the annihilator of the localized module with
  the extension of \(\mathrm{Ann}_{\mathcal{O}_X(V)}(\mathcal{F}(V))\) along the structure map. This is
  the asserted equality.
\end{proof}

\begin{lemma}\leanok
[Annihilator ideal sheaf agrees with the module annihilator on affine opens]
  \label{lem:modules_annihilator_ideal}
  \lean{AlgebraicGeometry.Scheme.Modules.annihilator_ideal}
  \uses{def:modules_annihilator, lem:modules_annihilator_ideal_le,
    lem:annihilator_map_basicOpen}
  Let \(X\) be a scheme and \(\mathcal{F}\) a quasi-coherent sheaf of \(\mathcal{O}_X\)-modules,
  and suppose the section module \(\mathcal{F}(V)\) is \emph{finitely generated} over
  \(\mathcal{O}_X(V)\) at \emph{every} affine open \(V \subseteq X\). Then at every affine open
  \(U \subseteq X\) the section over \(U\) of the annihilator ideal sheaf
  \(\mathrm{Ann}(\mathcal{F})\) (\cref{def:modules_annihilator}) \emph{equals} the
  module-annihilator of the sections:
  \[
    \mathrm{Ann}(\mathcal{F})(U)
    \;=\; \mathrm{Ann}_{\mathcal{O}_X(U)}\bigl(\mathcal{F}(U)\bigr).
  \]
  The global finiteness hypothesis is necessary: \(\mathrm{Ann}(\mathcal{F})\) is the
  \(\mathrm{IdealSheafData.ofIdeals}\) of the local annihilator data, i.e.\ the \emph{largest
  coherent} ideal sheaf dominated by that family (not the naive section-wise infimum), so reading
  off its value at one \(U\) requires the family to cohere at every affine open — exactly the
  global condition, in parallel with Mathlib's \(\mathrm{Hom.ker\_apply}\) (global
  \(\mathrm{QuasiCompact}\)). For a \emph{finite-type} quasi-coherent \(\mathcal{F}\) the
  hypothesis is supplied at every affine \(V\) by the G1 affine-section finiteness lemma
  (\cref{lem:gf_qcoh_fintype_finite_sections}, finite type being affine-local), so the finite-type
  form is the specialisation that discharges the hypothesis once G1 lands.
\end{lemma}
\begin{proof}
  \uses{lem:modules_annihilator_ideal_le, lem:annihilator_map_basicOpen}
  The forward inclusion \(\subseteq\) at each \(U\) is \cref{lem:modules_annihilator_ideal_le}. For
  the equality we assemble the local annihilator data into a genuine ideal sheaf and read off its
  value. The family \(V \mapsto \mathrm{Ann}_{\mathcal{O}_X(V)}(\mathcal{F}(V))\) satisfies the
  basic-open coherence condition \(\mathrm{map\_ideal\_basicOpen}\) at every affine open, by
  \cref{lem:annihilator_map_basicOpen} (using the finiteness hypothesis at each \(V\)); hence it
  defines an \(X.\mathrm{IdealSheafData}\) \(I\) with \(I(V) = \mathrm{Ann}_{\mathcal{O}_X(V)}(\mathcal{F}(V))\).
  By construction \(\mathrm{Ann}(\mathcal{F}) = \mathrm{ofIdeals}(I)\), and
  \(\mathrm{ofIdeals\_ideal}\) gives \(\mathrm{ofIdeals}(I).\mathrm{ideal} = I.\mathrm{ideal}\).
  Evaluating at the affine open \(U\) yields
  \(\mathrm{Ann}(\mathcal{F})(U) = I(U) = \mathrm{Ann}_{\mathcal{O}_X(U)}(\mathcal{F}(U))\),
  simultaneously at every affine \(U\).
\end{proof}

\begin{lemma}\leanok
[Annihilator of a finitely generated module commutes with
  localization]
  \label{lem:annihilator_localization_eq_map}
  \lean{Module.annihilator_isLocalizedModule_eq_map}
  Let \(R\) be a commutative ring, let \(S \subseteq R\) be a multiplicative
  submonoid, and let \(R_S = S^{-1}R\) be the localization of \(R\) at \(S\). Let
  \(M\) be a \emph{finitely generated} \(R\)-module and let \(f : M \to M_S\) be
  an \(R\)-linear map exhibiting \(M_S\) as the localized module \(S^{-1}M\)
  (i.e.\ \(f\) is a localization of \(M\) at \(S\)). Then the annihilator of
  \(M_S\) over \(R_S\) is the extension of the annihilator of \(M\) along the
  structure map \(R \to R_S\):
  \[
    \mathrm{Ann}_{R_S}(M_S)
    \;=\; (\mathrm{Ann}_R M)\cdot R_S
    \;=\; \mathrm{map}\bigl(R \to R_S\bigr)\bigl(\mathrm{Ann}_R M\bigr).
  \]
  Equivalently, for finitely generated \(M\) the annihilator commutes with
  localization: \(\mathrm{Ann}(S^{-1}M) = S^{-1}\,\mathrm{Ann}(M)\).

  \emph{The finite-generation hypothesis is essential.} For a module that is not
  finitely generated the inclusion \(S^{-1}\,\mathrm{Ann}(M) \subseteq
  \mathrm{Ann}(S^{-1}M)\) can be strict, so the identity fails. This is the
  standard commutative-algebra fact that the annihilator of a finite module
  localizes; it is recorded here as the load-bearing algebra engine for the
  basic-open coherence of \cref{def:modules_annihilator}.
\end{lemma}

\begin{proof}
  \(\supseteq\): if \(r \in \mathrm{Ann}_R M\), then \(r\) kills every element of
  \(M\), hence its image \(r/1 \in R_S\) kills every element of \(M_S\) (each such
  element is \(m/s = (1/s)\cdot f(m)\) for some \(m \in M\), \(s \in S\), and
  \((r/1)\cdot(m/s) = (1/s)\cdot f(r\,m) = 0\)). Thus
  \((\mathrm{Ann}_R M)\cdot R_S \subseteq \mathrm{Ann}_{R_S}(M_S)\).

  \(\subseteq\): write \(M = \langle m_1, \dots, m_n\rangle\) using the
  finite-generation hypothesis, and let \(y \in \mathrm{Ann}_{R_S}(M_S)\). Write
  \(y = a/s\) with \(a \in R\), \(s \in S\). For each generator \(m_i\) the element
  \(y\cdot f(m_i) = (a/s)\cdot f(m_i) = 0\) in \(M_S\); by the definition of the
  localized module this means there is \(u_i \in S\) with \(u_i\,a\,m_i = 0\) in
  \(M\). Put \(u = \prod_{i=1}^n u_i \in S\). Then \(u\,a\) kills every generator
  \(m_i\), hence kills all of \(M\), so \(u\,a \in \mathrm{Ann}_R M\). Finally
  \(y = a/s = (u\,a)/(u\,s)\) lies in \(S^{-1}\,\mathrm{Ann}_R M =
  (\mathrm{Ann}_R M)\cdot R_S\). Therefore \(\mathrm{Ann}_{R_S}(M_S) \subseteq
  (\mathrm{Ann}_R M)\cdot R_S\), and combined with the reverse inclusion the two
  ideals coincide. The use of a single common multiplier \(u = \prod_i u_i\) is
  exactly the step that requires finitely many generators.
\end{proof}

\begin{lemma}[Sections on a basic open localize the affine ring]
  \label{lem:isLocalization_basicOpen_mathlib}
  \lean{AlgebraicGeometry.IsAffineOpen.isLocalization_basicOpen}
  \mathlibok
  \textit{Provided by Mathlib
  (\texttt{Mathlib.AlgebraicGeometry.AffineScheme}).}
  Let \(X\) be a scheme, \(U \subseteq X\) an affine open, and
  \(f \in \Gamma(X, U)\). Then the restriction map
  \(\Gamma(X, U) \to \Gamma\bigl(X,\, D(f)\bigr)\) on the basic open
  \(D(f) = X.\mathrm{basicOpen}\,f\) exhibits \(\Gamma(X, D(f))\) as the
  localization of \(\Gamma(X, U)\) away from \(f\); that is, it is
  \(\mathrm{IsLocalization.Away}\,f\), so
  \(\Gamma(X, D(f)) = (\Gamma(X, U))_f\) as a \(\Gamma(X,U)\)-algebra.
\end{lemma}

\begin{lemma}\leanok
[Quasi-coherent sections on a basic open are the localized module]
  \label{lem:qcoh_section_localization_basicOpen}
  \lean{AlgebraicGeometry.Scheme.Modules.isLocalizedModule_basicOpen}
  \uses{lem:isLocalization_basicOpen_mathlib, lem:qcoh_affine_section_localization,
    lem:qcoh_affine_isIso_fromTildeΓ,
    lem:isLocalizedModule_restrict_of_isIso_fromTildeΓ,
    lem:qcoh_pullback_fromSpec, lem:modules_restrict_linear,
    lem:modules_restrict_basicOpen_linear, lem:fromSpec_image_top_section_coherence,
    lem:section_localization_hfr_aux_general, lem:isLocalizedModule_basicOpen_of_hP1}
  % SOURCE: [Stacks Project], "Schemes" (chapter tag 020J), \S 7
  % "Quasi-coherent sheaves on affines", Lemma \texttt{lemma-spec-sheaves},
  % item (4) (read from references/stacks-schemes.tex, L691-L702).
  % SOURCE QUOTE:
  % "Let $R$ be a ring. Let $M$ be an $R$-module. Let $\widetilde M$
  %  be the sheaf of $\mathcal{O}_{\Spec(R)}$-modules
  %  associated to $M$.
  %  [...]
  %  \item For every $f\in R$ we have $\Gamma(D(f), \widetilde M) = M_f$
  %  as an $R_f$-module."
  \textit{Source: [Stacks Project], "Schemes", \S 7, Lemma on the sheaves
  \(\widetilde M\) on \(\Spec(R)\), item (4).}
  Let \(X\) be a scheme and let \(\mathcal{M}\) be a \emph{quasi-coherent} sheaf
  of modules on \(X\). Let \(U \subseteq X\) be an
  affine open and \(f \in \Gamma(X, U)\). Then:
  \begin{enumerate}
    \item \(\Gamma\bigl(X, D(f)\bigr)\) is the localization \((\Gamma(X,U))_f\)
      of \(\Gamma(X,U)\) away from \(f\)
      (\cref{lem:isLocalization_basicOpen_mathlib}); and
    \item the section-restriction map
      \(\mathcal{M}(U) \to \mathcal{M}\bigl(D(f)\bigr)\) exhibits
      \(\mathcal{M}(D(f))\) as the localized module
      \((\mathrm{powers}\,f)^{-1}\,\mathcal{M}(U)\); that is, it is a localization
      of the \(\Gamma(X,U)\)-module \(\mathcal{M}(U)\) at the submonoid
      \(\mathrm{powers}\,f = \{1, f, f^2, \dots\}\).
  \end{enumerate}
\end{lemma}

\begin{proof}
  \uses{lem:isLocalization_basicOpen_mathlib, lem:qcoh_affine_section_localization,
    lem:qcoh_pullback_fromSpec, lem:qcoh_affine_isIso_fromTildeΓ,
    lem:section_localization_hfr_aux_general, lem:fromSpec_image_top_section_coherence,
    lem:modules_restrict_basicOpen_linear, def:gamma_image_ring_equiv,
    lem:isLocalizedModule_basicOpen_of_hP1}
  Part (1) is \cref{lem:isLocalization_basicOpen_mathlib}, the Mathlib statement
  that on an affine open \(U\) the sections over the basic open \(D(f)\) form the
  localization of \(\Gamma(X,U)\) away from \(f\).

  Part (2) is the \emph{gap-2 transport}. Because \(U\) is affine, Mathlib's
  canonical affine immersion \(\mathrm{fromSpec} : \Spec \Gamma(X, U) \to X\) has
  image \(U\) and identifies \(U\) with the spectrum of its ring of sections. We
  obtain (2) by instantiating the general-ambient-scheme transport core
  \cref{lem:section_localization_hfr_aux_general} at the affine open immersion
  \(j := \mathrm{fromSpec}\), then bridging its conclusion to the basic-open
  restriction.

  \emph{Supplying the P1 datum.} The core requires
  \(\mathrm{IsIso}\bigl(((\mathrm{pullback}\,\mathrm{fromSpec}).\mathrm{obj}\,
  \mathcal{M}).\mathrm{fromTildeΓ}\bigr)\). This is produced in two steps: first the
  pullback \((\mathrm{pullback}\,\mathrm{fromSpec}).\mathrm{obj}\,\mathcal{M}\) is
  again quasi-coherent by \cref{lem:qcoh_pullback_fromSpec} (quasi-coherence is
  preserved under pullback along \(\mathrm{fromSpec}\)); then, since this pullback
  lives on the affine scheme \(\Spec \Gamma(X, U)\),
  \cref{lem:qcoh_affine_isIso_fromTildeΓ} turns that quasi-coherence into the
  required isomorphism of the \(\widetilde{(-)}\)-counit.

  \emph{Applying the core.} Feeding this P1 datum, the element \(f\), and the slice
  element \(f' := \sigma^{-1}(f)\) --- where
  \(\sigma = (\Gamma\Spec\text{-iso})^{-1} \circ
  \mathrm{gammaImageRingEquiv}\,\mathrm{fromSpec}\,\top\)
  (\cref{def:gamma_image_ring_equiv}) is the \(\top\)-level section ring iso, so the
  transport hypothesis \(\sigma(f') = f\) holds by construction --- to
  \cref{lem:section_localization_hfr_aux_general} yields that the section
  restriction
  \(\Gamma(\mathcal{M}, \mathrm{fromSpec}\,''^{U}\,\top) \to
  \Gamma(\mathcal{M}, \mathrm{fromSpec}\,''^{U}\,D(f'))\) is
  \(\mathrm{IsLocalizedModule}(\mathrm{powers}\,f)\) over \(\Gamma(X, U)\).

  \emph{Bridging to the basic-open restriction.} It remains to identify this
  statement, posed on the image opens \(\mathrm{fromSpec}\,''^{U}\,\top\) and
  \(\mathrm{fromSpec}\,''^{U}\,D(f')\), with
  \(\mathrm{IsLocalizedModule}(\mathrm{powers}\,f)\) for the consumer-facing
  basic-open restriction \(\mathrm{restrictBasicOpenₗ}\)
  (\cref{lem:modules_restrict_basicOpen_linear}). The image opens are
  \(\mathrm{fromSpec}\,''^{U}\,\top = U\) and
  \(\mathrm{fromSpec}\,''^{U}\,D(f') = D(f)\). The one nontrivial coherence is that,
  under the identification \(\mathrm{fromSpec}\,''^{U}\,\top = U\), the section ring
  iso \(\sigma\) restricts to the identity --- precisely \(\rho(\sigma\,f') = f\),
  where \(\rho\) is the presheaf transport along that open identification --- which
  is exactly the section-coherence crux
  \cref{lem:fromSpec_image_top_section_coherence}. With it the localized-module
  property transports across the open identification isomorphisms to the
  basic-open restriction \(\mathcal{M}(U) \to \mathcal{M}(D(f))\), establishing (2).
\end{proof}

\begin{lemma}

\leanok
[Basic-open localization given the pullback-$\mathrm{fromTildeΓ}$ isomorphism datum]
  \label{lem:isLocalizedModule_basicOpen_of_hP1}
  \lean{AlgebraicGeometry.Scheme.Modules.isLocalizedModule_basicOpen_of_hP1}
  \uses{lem:section_localization_hfr_aux_general, lem:fromSpec_image_top_section_coherence,
    lem:modules_restrict_basicOpen_linear, lem:isLocalizedModule_ringEquiv_semilinear,
    def:gamma_image_ring_equiv}
  \textit{The $\mathrm{hP1}$-explicit form of \cref{lem:qcoh_section_localization_basicOpen} part (2):
  the basic-open restriction localizes when the quasi-coherence of the pullback is given as a hypothesis,
  and the proof reduces to a purely mechanical open/ring-iso bridge.}
  Let \(\mathcal{M}\) be a sheaf of modules on a scheme \(X\), \(U \subseteq X\) an affine open with
  canonical affine immersion \(\mathrm{fromSpec} : \Spec \Gamma(X,U) \to X\), and \(f \in \Gamma(X,U)\).
  Suppose
  \(\mathrm{hP1} : \mathrm{IsIso}\bigl(((\mathrm{pullback}\,\mathrm{fromSpec}).\mathrm{obj}\,
  \mathcal{M}).\mathrm{fromTildeΓ}\bigr)\). Then the basic-open restriction
  \(\mathrm{restrictBasicOpenₗ}\,\mathcal{M}\,f : \Gamma(\mathcal{M}, U) \to
  \Gamma(\mathcal{M}, X.\mathrm{basicOpen}\,f)\) is \(\mathrm{IsLocalizedModule}(\mathrm{powers}\,f)\) over
  \(\Gamma(X,U)\).
\end{lemma}
\begin{proof}
  \uses{lem:section_localization_hfr_aux_general, lem:fromSpec_image_top_section_coherence,
    lem:modules_restrict_basicOpen_linear, lem:isLocalizedModule_ringEquiv_semilinear,
    def:gamma_image_ring_equiv}
  Instantiate the general-ambient transport core \cref{lem:section_localization_hfr_aux_general} at the
  affine open immersion \(j := \mathrm{fromSpec}\) (its base ring is then \(\Gamma(X,U)\)), with slice
  element \(f' := \sigma^{-1}(f)\) where \(\sigma = \mathrm{gammaImageRingEquiv}\,j\,\top\)
  (\cref{def:gamma_image_ring_equiv}); the transport hypothesis \(\sigma(f') = f\) holds by construction.
  Fed the datum \(\mathrm{hP1}\), the core yields \(\mathrm{IsLocalizedModule}(\mathrm{powers}\,f')\) for
  the restriction \(\mathrm{restrictₗ}\,\mathcal{M} : \Gamma(\mathcal{M}, j\,''^{U}\top) \to
  \Gamma(\mathcal{M}, j\,''^{U}D(f'))\) over \(A := \Gamma(X, j\,''^{U}\top)\). Transport this to the
  consumer-facing basic-open restriction across: (i) the open identifications \(j\,''^{U}\top = U\) and
  \(j\,''^{U}D(f') = X.\mathrm{basicOpen}\,f\) (the latter Mathlib's
  \(\mathrm{IsAffineOpen.fromSpec\_image\_basicOpen}\)); (ii) the section ring iso \(\rho : A \cong
  \Gamma(X,U)\) induced by \(j\,''^{U}\top = U\), for which \(\rho(\sigma f') = f\) is the section-coherence
  crux \cref{lem:fromSpec_image_top_section_coherence}; and (iii) the semilinear localization-transfer
  lemma \cref{lem:isLocalizedModule_ringEquiv_semilinear} along \(\rho\), with
  \((\mathrm{powers}\,f').\mathrm{map}\,\rho = \mathrm{powers}\,f\) by \(\mathrm{Submonoid.map\_powers}\).
  This yields \(\mathrm{IsLocalizedModule}(\mathrm{powers}\,f)\) for
  \(\mathrm{restrictBasicOpenₗ}\,\mathcal{M}\,f\) (\cref{lem:modules_restrict_basicOpen_linear}).
\end{proof}

% --- gap-2 supporting infrastructure (the named pieces of the gap-2 transport above) ---

\begin{lemma}\leanok
[$\Gamma(X,U)$-linear section restriction]
  \label{lem:modules_restrict_linear}
  \lean{AlgebraicGeometry.Scheme.Modules.restrictₗ}
  Let \(\mathcal{M}\) be a sheaf of modules on a scheme \(X\) and let \(V \subseteq U\) be an inclusion
  of opens. The presheaf restriction \(\Gamma(\mathcal{M}, U) \to \Gamma(\mathcal{M}, V)\) is
  \(\Gamma(X, U)\)-linear, where the codomain \(\Gamma(\mathcal{M}, V)\) carries the
  \(\Gamma(X, U)\)-module structure obtained by restricting scalars along the structure-sheaf
  restriction \(\Gamma(X, U) \to \Gamma(X, V)\).
\end{lemma}
\begin{proof}
  Additivity is automatic. Compatibility with scalar multiplication is the semilinearity of presheaf
  restriction for a sheaf of modules: restriction commutes with the \(\Gamma(X, U)\)-action, which over
  the restrict-scalars codomain is exactly \(\Gamma(X, U)\)-linearity.
\end{proof}

\begin{lemma}\leanok
[$\Gamma(X,U)$-linear restriction to a basic open]
  \label{lem:modules_restrict_basicOpen_linear}
  \lean{AlgebraicGeometry.Scheme.Modules.restrictBasicOpenₗ}
  \uses{lem:modules_restrict_linear}
  Let \(\mathcal{M}\) be a sheaf of modules on a scheme \(X\), \(U \subseteq X\) an open, and
  \(f \in \Gamma(X, U)\). The presheaf restriction
  \(\Gamma(\mathcal{M}, U) \to \Gamma(\mathcal{M}, X.\mathrm{basicOpen}\,f)\) is \(\Gamma(X, U)\)-linear,
  where the codomain carries its \(\Gamma(X, U)\)-module structure through the scalar tower
  \(\Gamma(X, U) \to \Gamma(X, X.\mathrm{basicOpen}\,f)\) acting on
  \(\Gamma(\mathcal{M}, X.\mathrm{basicOpen}\,f)\).
\end{lemma}
\begin{proof}
  \uses{lem:modules_restrict_linear}
  Same content as \cref{lem:modules_restrict_linear}, in the scalar-tower presentation: linearity is the
  semilinearity of presheaf restriction together with the scalar-tower compatibility relating the
  \(\Gamma(X, U)\)- and \(\Gamma(X, X.\mathrm{basicOpen}\,f)\)-actions, using that the algebra map
  \(\Gamma(X, U) \to \Gamma(X, X.\mathrm{basicOpen}\,f)\) is the structure-sheaf restriction along
  \(X.\mathrm{basicOpen}\,f \le U\).
\end{proof}

\begin{lemma}\leanok
[Section coherence of the affine immersion $\mathrm{fromSpec}$]
  \label{lem:fromSpec_image_top_section_coherence}
  \lean{AlgebraicGeometry.Scheme.Modules.fromSpec_image_top_section_coherence}
  \uses{def:gamma_image_ring_equiv}
  Let \(X\) be a scheme and \(U \subseteq X\) an affine open, with canonical affine immersion
  \(\mathrm{fromSpec} : \Spec \Gamma(X, U) \to X\) of image \(U\), and write
  \(e_{\top} : \mathrm{fromSpec}\,''^{U}\,\top = U\) for the canonical identification of its image
  with \(U\). Let
  \(\sigma = (\Gamma\Spec\text{-iso})^{-1} \circ \mathrm{gammaImageRingEquiv}\,\mathrm{fromSpec}\,\top\)
  (\cref{def:gamma_image_ring_equiv}) be the \(\top\)-level section ring isomorphism, and let
  \(\rho : \mathcal{O}_X(U) \xrightarrow{\sim} \mathcal{O}_X(\mathrm{fromSpec}''^{U}\,\top)\) be the
  structure-sheaf isomorphism induced by \(e_{\top}\). Then \(\rho \circ \sigma = \mathrm{id}\), i.e.,
  \(\rho(\sigma\,f) = f\) for every \(f \in \Gamma(X, U)\).
\end{lemma}
\begin{proof}
  The composite \(\rho \circ \sigma\) is a ring endomorphism of \(\Gamma(X, U)\). Expanding the
  definition of \(\sigma\) via \cref{def:gamma_image_ring_equiv} and applying the self-application
  identity of \(\mathrm{fromSpec}\) over \(U\)---which identifies the image of the total open of
  \(\Spec\Gamma(X,U)\) under \(\mathrm{fromSpec}\) with \(U\)---together with presheaf naturality,
  the composite \(\rho \circ \sigma\) reduces to a \(\Gamma(X,U)\)-algebra endomorphism of
  \(\Gamma(X,U)\). Since the only such endomorphism is the identity, \(\rho(\sigma\,f) = f\).
\end{proof}

\begin{lemma}\leanok
[Basic-open Hfr along an abstract affine open immersion, general ambient scheme]
  \label{lem:section_localization_hfr_aux_general}
  \lean{AlgebraicGeometry.Scheme.Modules.section_localization_hfr_aux_general}
  \uses{lem:isLocalizedModule_restrict_of_isIso_fromTildeΓ, lem:gamma_pullback_image_iso,
    lem:gamma_pullback_image_iso_hom_semilinear, lem:gamma_pullback_image_iso_hom_naturality,
    def:gamma_image_ring_equiv, lem:isLocalizedModule_ringEquiv_semilinear,
    lem:modules_restrict_linear}
  This is the general-ambient-scheme analogue of \cref{lem:section_localization_hfr_aux}: the ambient
  scheme \(X\) need not be affine, and the localization ring is the section ring
  \(\Gamma(X, j\,''^{U}\,\top)\) of the image open.
  Let \(\mathcal{M}\) be a sheaf of modules on a scheme \(X\) and \(j : \Spec S \to X\) an affine open
  immersion with
  \(\mathrm{IsIso}\bigl(((\mathrm{pullback}\,j).\mathrm{obj}\,\mathcal{M}).\mathrm{fromTildeΓ}\bigr)\)
  (the P1 datum). Fix a slice element \(f' \in S\) and a section \(f \in \Gamma(X, j\,''^{U}\,\top)\) with
  \(\mathrm{gammaImageRingEquiv}\,j\,\top\,((\Gamma\Spec\text{-iso}_S)^{-1} f') = f\). Then the section
  restriction \(\Gamma(\mathcal{M}, j\,''^{U}\,\top) \to \Gamma(\mathcal{M}, j\,''^{U}\,D(f'))\) is
  \(\mathrm{IsLocalizedModule}(\mathrm{powers}\,f)\) over \(A := \Gamma(X, j\,''^{U}\,\top)\), the
  conclusion map being \(\mathrm{restrictₗ}\) of \cref{lem:modules_restrict_linear}.
\end{lemma}
\begin{proof}
  \uses{lem:isLocalizedModule_restrict_of_isIso_fromTildeΓ, lem:gamma_pullback_image_iso,
    lem:gamma_pullback_image_iso_hom_semilinear, lem:gamma_pullback_image_iso_hom_naturality,
    def:gamma_image_ring_equiv, lem:isLocalizedModule_ringEquiv_semilinear}
  Set \(M' = (\mathrm{pullback}\,j).\mathrm{obj}\,\mathcal{M}\); the P1 hypothesis gives
  \(\mathrm{IsIso}\,M'.\mathrm{fromTildeΓ}\), so \cref{lem:isLocalizedModule_restrict_of_isIso_fromTildeΓ}
  localizes the slice restriction \(\Gamma(M', \top) \to \Gamma(M', D(f'))\) at \(\mathrm{powers}\,f'\)
  over the slice ring \(S\). The \(\sigma\)-semilinear additive isomorphisms
  \(e_1 = \mathrm{gammaPullbackImageIso}\,j\,\mathcal{M}\,\top\) and
  \(e_2 = \mathrm{gammaPullbackImageIso}\,j\,\mathcal{M}\,D(f')\) (\cref{lem:gamma_pullback_image_iso}),
  over the section ring iso
  \(\sigma = (\Gamma\Spec\text{-iso}_S)^{-1} \circ \mathrm{gammaImageRingEquiv}\,j\,\top\)
  (\cref{def:gamma_image_ring_equiv}), are semilinear by
  \cref{lem:gamma_pullback_image_iso_hom_semilinear} and intertwine the slice localization map with the
  section restriction by \cref{lem:gamma_pullback_image_iso_hom_naturality}. Bridge (I)
  \cref{lem:isLocalizedModule_ringEquiv_semilinear} transports the localization across \(\sigma\),
  landing \(\mathrm{IsLocalizedModule}\) at \((\mathrm{powers}\,f').\mathrm{map}\,\sigma\); since
  \(\sigma(f') = f\) and \(\sigma\) carries powers to powers (ring maps preserve the submonoid of powers),
  this submonoid is \(\mathrm{powers}\,f\). Because here the base and
  target rings coincide (\(A = \Gamma(X, j\,''^{U}\,\top)\) is itself the localization ring), no further
  restriction-of-scalars bridge is needed, and the conclusion is
  \(\mathrm{IsLocalizedModule}(\mathrm{powers}\,f)\) over \(A\) for the restriction \(\mathrm{restrictₗ}\).
\end{proof}

% --- Pullback of a quasi-coherent sheaf along an open immersion is quasi-coherent ---

\begin{definition}\leanok
[Inverse of the slice equivalence sends the unit module to the unit module]
  \label{def:over_restrict_unit_iso_inv}
  \lean{AlgebraicGeometry.Scheme.Modules.overRestrictUnitIsoInv}
  \uses{def:over_restrict_equiv, lem:isIso_unitToPushforwardObjUnit_of_isIso}
  \textit{The inverse of the slice equivalence functor preserves the structure-sheaf (unit)
  module; dual to \cref{def:over_restrict_unit_iso}.}
  Let \(V \subseteq X\) be an open subscheme. The \emph{inverse} functor of the slice
  equivalence (\cref{def:over_restrict_equiv}) carries the unit (structure-sheaf) module \(\mathbf{1}\)
  of \(V.\mathrm{toScheme}.\mathrm{ringCatSheaf}\) to the unit module of the over-site
  \(X.\mathrm{ringCatSheaf}.\mathrm{over}\,V\):
  \[
    (\mathrm{overRestrictEquiv}\,V).\mathrm{inverse}.\mathrm{obj}\,
      \bigl(\mathbf{1}\,(V.\mathrm{toScheme}.\mathrm{ringCatSheaf})\bigr)
      \;\cong\; \mathbf{1}\,(X.\mathrm{ringCatSheaf}.\mathrm{over}\,V) .
  \]
\end{definition}
\begin{proof}
  \uses{def:over_restrict_equiv, lem:isIso_unitToPushforwardObjUnit_of_isIso}
  As in the forward case (\cref{def:over_restrict_unit_iso}), the inverse functor of the equivalence is
  the pushforward \(\mathrm{pushforward}\,\varphi\) of sheaves of modules along an over-site/open-site
  ring comparison \(\varphi\) built from the (op of the) unit isomorphism of
  \((\mathrm{Opens.overEquivalence}\,V)\); since that unit is an isomorphism, \(\varphi\) is an
  isomorphism. The canonical unit comparison \(\mathrm{unitToPushforwardObjUnit}\,\varphi\) is then an
  isomorphism by the \(\mathrm{IsIso}\,\varphi\)-driven criterion
  (\cref{lem:isIso_unitToPushforwardObjUnit_of_isIso}), and its inverse is the asserted isomorphism.
\end{proof}

\begin{definition}\leanok
[Geometric presentation back-transported to a slice presentation]
  \label{def:over_restrict_presentation_inv}
  \lean{AlgebraicGeometry.Scheme.Modules.overRestrictPresentationInv}
  \uses{def:over_restrict_unit_iso_inv, lem:over_restrict_pullback_iso, lem:presentation_map_mathlib}
  \textit{Dual to \cref{def:over_restrict_presentation}: a presentation of the geometric pullback
  yields a presentation of the abstract Grothendieck slice.}
  Let \(V \subseteq X\) be open and \(M\) a sheaf of modules on \(X\). A presentation of the geometric
  pullback \((V.\iota^{*})\,M\) yields a presentation of the abstract Grothendieck-slice restriction
  \(M.\mathrm{over}\,V\):
  \[
    \bigl((V.\iota^{*})\,M\bigr).\mathrm{Presentation}
      \;\longrightarrow\;
    (M.\mathrm{over}\,V).\mathrm{Presentation} .
  \]
\end{definition}
\begin{proof}
  \uses{def:over_restrict_unit_iso_inv, lem:over_restrict_pullback_iso, lem:presentation_map_mathlib}
  First transport the given presentation of \((V.\iota^{*})\,M\) across the inverse of the pullback
  packaging isomorphism \((\mathrm{overRestrictEquiv}\,V).\mathrm{functor}.\mathrm{obj}\,
  (M.\mathrm{over}\,V) \cong (V.\iota^{*})\,M\) (\cref{lem:over_restrict_pullback_iso}) via
  \(\mathrm{Presentation.ofIsIso}\), obtaining a presentation of the transported slice module
  \((\mathrm{overRestrictEquiv}\,V).\mathrm{functor}.\mathrm{obj}\,(M.\mathrm{over}\,V)\). Then apply
  \(\mathrm{Presentation.map}\) (\cref{lem:presentation_map_mathlib}) along the inverse functor
  \((\mathrm{overRestrictEquiv}\,V).\mathrm{inverse}\), using \cref{def:over_restrict_unit_iso_inv} to
  match the image of the structure-sheaf (unit) module with the target unit module; this lands a
  presentation of \((\mathrm{overRestrictEquiv}\,V).\mathrm{inverse}.\mathrm{obj}\,
  \bigl((\mathrm{overRestrictEquiv}\,V).\mathrm{functor}.\mathrm{obj}\,(M.\mathrm{over}\,V)\bigr)\).
  Finally collapse the round trip across the counit isomorphism of the equivalence
  (\(\mathrm{inverse} \circ \mathrm{functor} \cong \mathrm{id}\)) via \(\mathrm{Presentation.ofIsIso}\),
  giving a presentation of \(M.\mathrm{over}\,V\).
\end{proof}

\begin{definition}\leanok
[Geometric pullback-square iso over a preimage cover member]
  \label{def:pullback_preimage_iota_iso}
  \lean{AlgebraicGeometry.Scheme.Modules.pullbackPreimageιIso}
  \uses{lem:modules_pullback_mathlib}
  \textit{Pseudofunctoriality iso realising the commuting preimage square.}
  Let \(g : Y \to X\) be an open immersion, \(M\) a sheaf of modules on \(X\), and \(U \subseteq X\) an
  open with preimage \(g^{-1}(U)\). Writing \(k : g^{-1}(U).\mathrm{toScheme} \to U.\mathrm{toScheme}\) for
  the open immersion induced by \(g\) (so \(k \circ U.\iota = (g^{-1}U).\iota \circ g\)), the
  pseudofunctoriality of pullback (\(\mathrm{pullbackComp}\)) gives a natural isomorphism
  \[
    \bigl((g^{-1}U).\iota^{*}\bigr)\,(g^{*}M) \;\cong\; k^{*}\bigl(U.\iota^{*}\,M\bigr) .
  \]
\end{definition}
\begin{proof}
  \uses{lem:modules_pullback_mathlib}
  Compose the pullback-composition coherence isos along the two ways of reading
  \((g^{-1}U).\iota \circ g = k \circ U.\iota\): \(\mathrm{pullbackComp}\,k\,U.\iota\), the congruence
  \(\mathrm{pullbackCongr}\) of the square identity, and \((\mathrm{pullbackComp}\,(g^{-1}U).\iota\,g)\)
  inverted. Each is an instance of the \(\mathrm{pullbackComp}\) pseudofunctoriality of the pullback of
  sheaves of modules (\cref{lem:modules_pullback_mathlib}).
\end{proof}

\begin{definition}\leanok
[Presentation of the pullback's geometric restriction to a preimage cover member]
  \label{def:presentation_pullback_iota_preimage}
  \lean{AlgebraicGeometry.Scheme.Modules.presentationPullbackιPreimage}
  \uses{def:presentation_pullback_iota_of_quasicoherentData, lem:presentation_map_mathlib,
    lem:modules_pullback_mathlib, def:pullback_preimage_iota_iso, def:pullback_open_immersion_unit_iso}
  \textit{For an open immersion $g$, the pulled-back sheaf has a presentation on the geometric
  restriction to each preimage cover member.}
  Let \(g : Y \to X\) be an open immersion, \(M\) a sheaf of modules on \(X\) with quasi-coherence data
  \(q\), and \(i\) a cover index. Write \(N := (\mathrm{pullback}\,g).\mathrm{obj}\,M\) and
  \(W_i := g^{-1}(q.X_i)\). Then the geometric restriction
  \((W_i.\iota^{*})\,N\) admits a \(\mathrm{SheafOfModules.Presentation}\) on \(W_i.\mathrm{toScheme}\).
\end{definition}
\begin{proof}
  \uses{def:presentation_pullback_iota_of_quasicoherentData, lem:presentation_map_mathlib,
    lem:modules_pullback_mathlib}
  The open immersion \(g\) restricts to an open immersion
  \(k : W_i.\mathrm{toScheme} \to (q.X_i).\mathrm{toScheme}\) on preimages, characterized by the
  commuting square \(k \circ (q.X_i).\iota = W_i.\iota \circ g\) (the morphism induced by \(g\) between
  the inclusion of \(W_i = g^{-1}(q.X_i)\) and that of \(q.X_i\)). By pseudofunctoriality of pullback
  (\(\mathrm{pullbackComp}\), named in prose) this square gives a natural isomorphism
  \[
    (W_i.\iota^{*})\,N \;=\; (W_i.\iota^{*})\,(g^{*}M)
      \;\cong\; (W_i.\iota \circ g)^{*}M
      \;=\; (k \circ (q.X_i).\iota)^{*}M
      \;\cong\; k^{*}\bigl((q.X_i).\iota^{*}\,M\bigr) .
  \]
  Now \((q.X_i).\iota^{*}\,M\) carries the geometric presentation
  \cref{def:presentation_pullback_iota_of_quasicoherentData}. Apply \(\mathrm{Presentation.map}\)
  (\cref{lem:presentation_map_mathlib}) along the pullback functor \(k^{*}\)
  (\cref{lem:modules_pullback_mathlib}), which preserves colimits and carries the unit sheaf to the
  unit sheaf, to obtain a presentation of \(k^{*}\bigl((q.X_i).\iota^{*}\,M\bigr)\); transport it across
  the displayed isomorphism via \(\mathrm{Presentation.ofIsIso}\) to land a presentation of
  \((W_i.\iota^{*})\,N\).
\end{proof}

\begin{lemma}\leanok
[The pulled-back sheaf is quasi-coherent on each preimage cover member]
  \label{lem:isQuasicoherent_over_preimage}
  \lean{AlgebraicGeometry.Scheme.Modules.isQuasicoherent_over_preimage}
  \uses{def:over_restrict_presentation_inv, def:presentation_pullback_iota_preimage,
    lem:presentation_isQuasicoherent_mathlib}
  \textit{The abstract Grothendieck slice of the pullback over a preimage cover member is
  quasi-coherent.}
  With \(g : Y \to X\) an open immersion, \(M\) quasi-coherent with datum \(q\),
  \(N := (\mathrm{pullback}\,g).\mathrm{obj}\,M\) and \(W_i := g^{-1}(q.X_i)\), the Grothendieck-slice
  restriction \(N.\mathrm{over}\,W_i\) is quasi-coherent.
\end{lemma}
\begin{proof}
  \uses{def:over_restrict_presentation_inv, def:presentation_pullback_iota_preimage,
    lem:presentation_isQuasicoherent_mathlib}
  Feed the geometric presentation of \((W_i.\iota^{*})\,N\)
  (\cref{def:presentation_pullback_iota_preimage}) into the geometric-to-slice back-transport
  \cref{def:over_restrict_presentation_inv} at \(V = W_i\), producing a presentation of the slice
  \(N.\mathrm{over}\,W_i\). A sheaf with a presentation is quasi-coherent
  (\cref{lem:presentation_isQuasicoherent_mathlib}).
\end{proof}

\begin{lemma}\leanok
[The preimage family of a quasi-coherence cover covers the source]
  \label{lem:coversTop_preimage}
  \lean{AlgebraicGeometry.Scheme.Modules.coversTop_preimage}
  \uses{lem:isQuasicoherent_quasicoherentData_mathlib}
  \textit{Preimages of a top cover under a morphism again cover the top.}
  Let \(g : Y \to X\) be a morphism of schemes and \(q\) quasi-coherence data for a sheaf \(M\) on
  \(X\), with cover \(\{q.X_i\}\) of \(X\) (\cref{lem:isQuasicoherent_quasicoherentData_mathlib}). Then
  the preimage family \(\{\,W_i := g^{-1}(q.X_i)\,\}\) is a \(\mathrm{CoversTop}\) family on \(Y\).
\end{lemma}
\begin{proof}
  \uses{lem:isQuasicoherent_quasicoherentData_mathlib}
  The covering property \(q.\mathrm{coversTop}\) says \(\{q.X_i\}\) generates the maximal sieve on
  \(\top_X\); pulling back along the continuous map underlying \(g\) commutes with the sieve operations
  and sends \(\top_X\) to \(\top_Y\) (\(g^{-1}\top = \top\) and \(g^{-1}\) preserves suprema), so the
  preimage family \(\{g^{-1}(q.X_i)\}\) generates the maximal sieve on \(\top_Y\), i.e.\ it covers the
  top.
\end{proof}

\begin{lemma}\leanok
[Pullback of a quasi-coherent sheaf along an open immersion is quasi-coherent]
  \label{lem:isQuasicoherent_pullback_of_isOpenImmersion}
  \lean{AlgebraicGeometry.Scheme.Modules.isQuasicoherent_pullback_of_isOpenImmersion}
  \uses{lem:isQuasicoherent_quasicoherentData_mathlib, lem:isQuasicoherent_over_preimage,
    lem:coversTop_preimage, lem:isQuasicoherent_of_coversTop_mathlib}
  \textit{The general statement of which \cref{lem:qcoh_pullback_fromSpec} is the
  \(g = \mathrm{fromSpec}\) instance.}
  Let \(g : Y \to X\) be an open immersion and \(M\) a quasi-coherent sheaf of modules on \(X\). Then
  the pullback \((\mathrm{pullback}\,g).\mathrm{obj}\,M\) is quasi-coherent.
\end{lemma}
\begin{proof}
  \uses{lem:isQuasicoherent_quasicoherentData_mathlib, lem:isQuasicoherent_over_preimage,
    lem:coversTop_preimage, lem:isQuasicoherent_of_coversTop_mathlib}
  Choose quasi-coherence data \(q\) for \(M\) (\cref{lem:isQuasicoherent_quasicoherentData_mathlib})
  and set \(N := (\mathrm{pullback}\,g).\mathrm{obj}\,M\). The preimage family
  \(\{W_i := g^{-1}(q.X_i)\}\) covers \(Y\) (\cref{lem:coversTop_preimage}), and on each member the
  slice \(N.\mathrm{over}\,W_i\) is quasi-coherent (\cref{lem:isQuasicoherent_over_preimage}). By the
  local-to-global criterion \(\mathrm{IsQuasicoherent.of\_coversTop}\)
  (\cref{lem:isQuasicoherent_of_coversTop_mathlib}) applied to this cover, \(N\) is quasi-coherent.
\end{proof}

\begin{lemma}\leanok
[Quasi-coherence is preserved under pullback along $\mathrm{fromSpec}$]
  \label{lem:qcoh_pullback_fromSpec}
  \lean{AlgebraicGeometry.Scheme.Modules.isQuasicoherent_pullback_fromSpec}
  \uses{lem:isQuasicoherent_pullback_of_isOpenImmersion}
  Let \(\mathcal{M}\) be a quasi-coherent sheaf of modules on a scheme \(X\) and let \(U \subseteq X\) be
  an affine open with canonical affine immersion \(\mathrm{fromSpec} : \Spec \Gamma(X, U) \to X\). Then
  the pullback \((\mathrm{pullback}\,\mathrm{fromSpec}).\mathrm{obj}\,\mathcal{M}\) is quasi-coherent.
\end{lemma}
\begin{proof}
  \uses{lem:isQuasicoherent_pullback_of_isOpenImmersion}
  The affine immersion \(\mathrm{fromSpec} : \Spec \Gamma(X, U) \to X\) of an affine open is an open
  immersion (\(\mathrm{IsAffineOpen.fromSpec}\) is \(\mathrm{IsOpenImmersion}\), named in prose). Hence
  this is the \(g := \mathrm{fromSpec}\) instance of the general statement
  \cref{lem:isQuasicoherent_pullback_of_isOpenImmersion}: the pullback of the quasi-coherent sheaf
  \(\mathcal{M}\) along the open immersion \(\mathrm{fromSpec}\) is quasi-coherent.
\end{proof}

\begin{lemma}\leanok
[Basic-open restriction of a $\widetilde{N}$ sheaf is a module localization]
  \label{lem:isLocalizedModule_tilde_restrict}
  \lean{AlgebraicGeometry.isLocalizedModule_tilde_restrict}
  % SOURCE: [Stacks Project], "Schemes" (tag 020J), \S 7, Lemma on the sheaves
  % \widetilde M on Spec(R), item (4) (read from references/stacks-schemes.tex, L691-L702).
  % SOURCE QUOTE:
  % "For every $f\in R$ we have $\Gamma(D(f), \widetilde M) = M_f$ as an $R_f$-module."
  \textit{Source: [Stacks Project], "Schemes", \S 7, item (4); Spec-local core of
  \cref{lem:qcoh_section_localization_basicOpen}, part (2).}
  Let \(N\) be an \(R\)-module and \(f \in R\). Then the presheaf restriction map of the
  associated sheaf \(\widetilde{N}\) on \(\Spec R\) from the global sections
  \(\Gamma(\widetilde{N}, \top)\) to the sections on the basic open \(\Gamma(\widetilde{N}, D(f))\)
  exhibits \(\Gamma(\widetilde{N}, D(f))\) as the localized module
  \((\mathrm{powers}\,f)^{-1} N\); that is, it is \(\mathrm{IsLocalizedModule}(\mathrm{powers}\,f)\)
  over \(R\).
\end{lemma}
\begin{proof}
  Mathlib supplies the localization-out-of-\(N\) instance
  \(\mathrm{IsLocalizedModule}(\mathrm{powers}\,f)\) for the canonical map
  \(N \to \Gamma(\widetilde N, D(f))\). Since the global-sections map
  \(N \to \Gamma(\widetilde N, \top)\) is an isomorphism and the canonical map to \(D(f)\) factors
  through it followed by the restriction, precomposing the localization map with the inverse
  isomorphism transfers the \(\mathrm{IsLocalizedModule}\) property to the restriction map
  (Mathlib's transfer lemma for localization along a linear equivalence).
\end{proof}

\begin{lemma}\leanok
[Basic-open restriction localizes for a sheaf in the essential image of $\widetilde{(-)}$]
  \label{lem:isLocalizedModule_restrict_of_isIso_fromTildeΓ}
  \lean{AlgebraicGeometry.isLocalizedModule_restrict_of_isIso_fromTildeΓ}
  \uses{lem:isLocalizedModule_tilde_restrict}
  \textit{Affine engine of \cref{lem:qcoh_section_localization_basicOpen}: the basic-open
  restriction localizes for any module sheaf that is \(\widetilde{N}\) up to isomorphism.}
  Let \(M\) be a sheaf of modules on \(\Spec R\) such that the counit
  \(M.\mathrm{fromTildeΓ}\) is an isomorphism (equivalently, \(M\) lies in the essential image of
  \(\widetilde{(-)}\); Mathlib's \(\mathrm{isIso\_fromTildeΓ\_iff}\)), and let \(f \in R\). Then the
  presheaf restriction of \(M\) from \(\Gamma(M, \top)\) to \(\Gamma(M, D(f))\) is
  \(\mathrm{IsLocalizedModule}(\mathrm{powers}\,f)\) over \(R\).
\end{lemma}
\begin{proof}
  The image of \(M.\mathrm{fromTildeΓ}\) under \(\mathrm{modulesSpecToSheaf}\) is an isomorphism;
  forgetting to the presheaf gives a natural isomorphism \(\psi\) whose components are isomorphisms
  of \(R\)-modules. Its naturality square for the inclusion \(D(f) \hookrightarrow \top\)
  intertwines the restriction maps of \(\widetilde{N}\) (for \(N = \Gamma(M, \top)\)) and of \(M\).
  Transporting \cref{lem:isLocalizedModule_tilde_restrict} across the two component isomorphisms
  (pre- and post-composing by the linear equivalences induced by the components of \(\psi\))
  identifies the restriction of \(M\) with a localization map, hence it is
  \(\mathrm{IsLocalizedModule}(\mathrm{powers}\,f)\).
\end{proof}

\begin{lemma}[A section restricted to its non-vanishing locus is a unit (endomorphism form)]
  \label{lem:isUnit_algebraMap_end_basicOpen_mathlib}
  \lean{AlgebraicGeometry.tilde.isUnit_algebraMap_end_basicOpen}
  \mathlibok
  \textit{Provided by Mathlib (\texttt{Mathlib.AlgebraicGeometry.Modules.Tilde}).}
  Let \(M\) be a sheaf of modules on \(\Spec R\) and \(f \in R\). Then \(f\) acts as a unit
  endomorphism of the sections \(\Gamma(M, D(f))\) over the basic open
  \(D(f) = \mathrm{basicOpen}\,f\); that is,
  \[
    \mathrm{IsUnit}\bigl(\mathrm{algebraMap}\,R\,(\mathrm{Module.End}_R\,\Gamma(M, D(f)))\,f\bigr).
  \]
  Equivalently, the image of \(f\) in the endomorphism ring of \(\Gamma(M, D(f))\) is invertible.
  No quasi-coherence of \(M\) is required: it follows from the structure ring
  \(\Gamma(\mathcal{O}_{\Spec R}, D(f)) = R_f\) being the away-localization
  (\(\mathrm{IsLocalization.Away}\,f\)), through which the \(R\)-action on \(\Gamma(M, D(f))\)
  factors.
\end{lemma}

\begin{lemma}\leanok
[Basic-open restriction localizes, for a globally presented module]
  \label{lem:isLocalizedModule_basicOpen_of_presentation}
  \lean{AlgebraicGeometry.isLocalizedModule_basicOpen_of_presentation}
  \uses{lem:isIso_fromTildeΓ_of_presentation_mathlib,
    lem:isLocalizedModule_restrict_of_isIso_fromTildeΓ}
  \textit{The Route-F endpoint for the globally-presented case.}
  Let \(M\) be a sheaf of modules on \(\Spec R\) admitting a global presentation (a
  \(\mathrm{SheafOfModules.Presentation}\), i.e.\ a cokernel-of-free-modules description), and let
  \(f \in R\). Then the section restriction \(\Gamma(M, \top) \to \Gamma(M, D(f))\) exhibits the
  target as the localized module \((\mathrm{powers}\,f)^{-1}\,\Gamma(M, \top)\); that is, it is
  \(\mathrm{IsLocalizedModule}(\mathrm{powers}\,f)\) over \(R\).
\end{lemma}
\begin{proof}
  \uses{lem:isIso_fromTildeΓ_of_presentation_mathlib,
    lem:isLocalizedModule_restrict_of_isIso_fromTildeΓ}
  A global presentation forces the counit \(M.\mathrm{fromTildeΓ}\) to be an isomorphism
  (\cref{lem:isIso_fromTildeΓ_of_presentation_mathlib}). With \(\mathrm{IsIso}\,M.\mathrm{fromTildeΓ}\)
  in hand, the affine engine \cref{lem:isLocalizedModule_restrict_of_isIso_fromTildeΓ} delivers, for
  the given \(f\), that the section restriction is \(\mathrm{IsLocalizedModule}(\mathrm{powers}\,f)\).
\end{proof}

\begin{lemma}\leanok
[The \emph{map\_units} field holds for any sheaf of modules]
  \label{lem:map_units_restrict_basicOpen}
  \lean{AlgebraicGeometry.map_units_restrict_basicOpen}
  \uses{lem:isUnit_algebraMap_end_basicOpen_mathlib}
  \textit{The \(\mathrm{map\_units}\) field of G1-core, holding unconditionally.}
  Let \(M\) be a sheaf of modules on \(\Spec R\) and \(f \in R\). Then every element of the submonoid
  \(\mathrm{powers}\,f = \{1, f, f^2, \dots\}\) acts invertibly on the sections \(\Gamma(M, D(f))\):
  for each \(x \in \mathrm{powers}\,f\),
  \(\mathrm{IsUnit}(\mathrm{algebraMap}\,R\,(\mathrm{Module.End}_R\,\Gamma(M, D(f)))\,x)\). This is
  precisely the \emph{map\_units} field of the three-field \(\mathrm{IsLocalizedModule}\) predicate,
  in the shape the constructor consumes.
\end{lemma}
\begin{proof}
  \uses{lem:isUnit_algebraMap_end_basicOpen_mathlib}
  By \cref{lem:isUnit_algebraMap_end_basicOpen_mathlib} the element \(f\) acts as a unit
  endomorphism of \(\Gamma(M, D(f))\); a power of a unit is a unit, so each \(f^n\) does too. As this
  needs no quasi-coherence of \(M\), the \emph{map\_units} field is free and contributes nothing to
  the genuine content of G1-core (which is \emph{surj} / \emph{exists\_of\_eq}, supplied through the
  \(\widetilde{(-)}\)-counit being an isomorphism).
\end{proof}

\begin{lemma}[A presentation makes the $\widetilde{(-)}$-counit an isomorphism]
  \label{lem:isIso_fromTildeΓ_of_presentation_mathlib}
  \lean{AlgebraicGeometry.isIso_fromTildeΓ_of_presentation}
  \mathlibok
  \textit{Provided by Mathlib (\texttt{Mathlib.AlgebraicGeometry.Modules.Tilde}).}
  Let \(M\) be a sheaf of modules on \(\Spec R\) equipped with a presentation (a
  cokernel-of-free-modules description). Then the counit
  \(M.\mathrm{fromTildeΓ} : \widetilde{\Gamma(M,\top)} \to M\) is an isomorphism. Applied on an
  affine \(\Spec R_g\) this identifies \(M\) with \(\widetilde{N_g}\); together with the affine
  tilde localization instance (wrapped by \cref{lem:isLocalizedModule_tilde_restrict}), it supplies
  the base case of the descent.
\end{lemma}

\begin{lemma}[Sheaf Mayer--Vietoris: the pullback cone over a two-open cover is a limit]
  \label{lem:isLimitPullbackCone_mathlib}
  \lean{TopCat.Sheaf.isLimitPullbackCone}
  \mathlibok
  \textit{Provided by Mathlib
  (\texttt{Mathlib.Topology.Sheaves.SheafCondition.PairwiseIntersections}).}
  For a sheaf \(F\) on a topological space and two opens \(U, V\), the square relating
  \(F(U \cup V)\), \(F(U)\), \(F(V)\), \(F(U \cap V)\) is a pullback (limit). This is the generic
  module Mayer--Vietoris used in the inductive step of the descent.
\end{lemma}

\begin{lemma}[Unique gluing of compatible sections]
  \label{lem:existsUnique_gluing_mathlib}
  \lean{TopCat.Sheaf.existsUnique_gluing}
  \mathlibok
  \textit{Provided by Mathlib
  (\texttt{Mathlib.Topology.Sheaves.SheafCondition.UniqueGluing}); the companion
  \(\mathrm{eq\_of\_locally\_eq'}\) lives in the same file.}
  For a sheaf \(F\) and a family of opens \(\{U_i\}\) with union \(V\), any family of sections
  \(s_i \in F(U_i)\) agreeing on overlaps glues to a unique \(s \in F(V)\); and a section of \(V\)
  is determined by its restrictions to a cover (\(\mathrm{eq\_of\_locally\_eq'}\)). These discharge
  the existence and uniqueness obligations in the inductive step. The primed companion
  \(\mathrm{existsUnique\_gluing'}\), which presents the same statement for a family of opens whose
  supremum is a fixed open \(V\), is the form consumed by the cover-form descent
  (\cref{lem:section_localization_descent_of_cover}).
\end{lemma}

\begin{lemma}[Separatedness: a section is determined by its restrictions to a cover]
  \label{lem:eq_of_locally_eq_mathlib}
  \lean{TopCat.Sheaf.eq_of_locally_eq'}
  \mathlibok
  \textit{Provided by Mathlib
  (\texttt{Mathlib.Topology.Sheaves.SheafCondition.UniqueGluing}).}
  Let \(F\) be a sheaf on a topological space, \(V\) an open, and \(\{U_i\}\) a family of opens with
  \(\bigsqcup_i U_i = V\). If two sections \(s, t \in F(V)\) have equal restrictions
  \(s|_{U_i} = t|_{U_i}\) for every \(i\), then \(s = t\). This separatedness over a cover is the
  uniqueness half of the sheaf condition; it discharges the \(\mathrm{exists\_of\_eq}\) obligation
  (a global section vanishing on \(D(f)\) is killed by a power of \(f\)) and the final
  ``conclude on \(D(f)\)'' step of the cover-form descent
  (\cref{lem:section_localization_descent_of_cover}).
\end{lemma}

\begin{lemma}\leanok
[Section-localization for a quasi-coherent sheaf on $\Spec R$ (G1-core)]
  \label{lem:qcoh_affine_section_localization}
  \lean{AlgebraicGeometry.Scheme.Modules.isLocalizedModule_basicOpen_of_isQuasicoherent}
  \uses{lem:qcoh_affine_isIso_fromTildeΓ, lem:isLocalizedModule_restrict_of_isIso_fromTildeΓ,
    lem:isLocalization_basicOpen_mathlib}
  % NOTE: the Lean decl does NOT yet exist. G1-core is now a DOWNSTREAM COROLLARY of gap1
  % (\cref{lem:qcoh_affine_isIso_fromTildeΓ}): it is no longer proved by a direct 3-field
  % compact-open induction. The single irreducible content is gap1 (IsIso M.fromTildeΓ); once that
  % holds, all three IsLocalizedModule fields are delivered at once by
  % \cref{lem:isLocalizedModule_restrict_of_isIso_fromTildeΓ}. G1-core and gap1 are in fact
  % EQUIVALENT via \cref{lem:isIso_fromTildeΓ_iff_isLocalizedModule_restrict} (both engines are
  % axiom-clean in-file). It is the affine instance (X = Spec R, U = ⊤) of
  % \cref{lem:qcoh_section_localization_basicOpen}. Route in
  % analogies/quot-qcoh-affine-globalization.md, Route C / Stacks 01HA.
  % SOURCE: [Stacks Project], "Properties of Schemes", tag 01HA (\S Quasi-coherent sheaves,
  % "$\mathcal{F}(D(f)) = \mathcal{F}(X)_f$ for $\mathcal{F}$ quasi-coherent on affine $X$")
  % (read from references/stacks-properties.tex, L2152--2170). Cf. [Hartshorne], II.5.2.
  % SOURCE QUOTE:
  % "Let $X$ be a scheme. Let $f \in \Gamma(X, \mathcal{O}_X)$.
  %  Denote $X_f \subset X$ the open where $f$ is invertible. [...]
  %  if $\mathcal{F}$ is a quasi-coherent sheaf of $\mathcal{O}_X$-modules the map
  %  $$\Gamma(X, \mathcal{F})_f \longrightarrow \Gamma(X_f, \mathcal{F})$$
  %  is an isomorphism."
  % (Stacks, lemma-invert-f-sections; affine case $X = \Spec R$, $X_f = D(f)$.)
  \textit{Source: [Stacks Project], tag 01HA; cf.\ [Hartshorne], II.5.2.}
  Let \(M\) be a quasi-coherent sheaf of modules on \(\Spec R\) and \(f \in R\). Then the section
  restriction map \(\Gamma(M, \top) \to \Gamma(M, D(f))\) exhibits \(\Gamma(M, D(f))\) as the
  localized module \((\mathrm{powers}\,f)^{-1}\,\Gamma(M, \top)\); i.e.\ it is
  \(\mathrm{IsLocalizedModule}(\mathrm{powers}\,f)\) over \(R\).
\end{lemma}
% NOTE: G1-core is now derived from gap1, not by a direct 3-field compact-open induction. The proof
% is the forward leg of the equivalence \cref{lem:isIso_fromTildeΓ_iff_isLocalizedModule_restrict}.
% The LEMMA STATEMENT is cited to Stacks 01HA above; the proof is a project-internal reduction (no
% transcribed Stacks proof), so there is no `% SOURCE QUOTE PROOF`.
\begin{proof}
  \uses{lem:qcoh_affine_isIso_fromTildeΓ, lem:isLocalizedModule_restrict_of_isIso_fromTildeΓ,
    lem:isIso_fromTildeΓ_iff_isLocalizedModule_restrict}
  Since \(M\) is quasi-coherent, gap1 (\cref{lem:qcoh_affine_isIso_fromTildeΓ}) gives that the
  \(\widetilde{(-)}\)-counit \(M.\mathrm{fromTildeΓ}\) is an isomorphism. With
  \(\mathrm{IsIso}\,M.\mathrm{fromTildeΓ}\) in hand, the affine engine
  \cref{lem:isLocalizedModule_restrict_of_isIso_fromTildeΓ} delivers, for the given \(f\), that the
  section restriction \(\Gamma(M, \top) \to \Gamma(M, D(f))\) is
  \(\mathrm{IsLocalizedModule}(\mathrm{powers}\,f)\) over \(R\) -- this single application supplies
  all three predicate fields (\emph{map\_units}, \emph{surj}, \emph{exists\_of\_eq}) at once. The two
  statements are in fact equivalent: \cref{lem:isIso_fromTildeΓ_iff_isLocalizedModule_restrict}
  records that \(\mathrm{IsIso}\,M.\mathrm{fromTildeΓ}\) holds iff the section restriction localizes
  on every basic open, so G1-core and gap1 are interderivable and the genuine content sits entirely
  in gap1.
\end{proof}

% --- Quasi-coherence datum and the gap1 sub-build ---

\begin{lemma}[Quasi-coherence = a cover with local presentations]
  \label{lem:isQuasicoherent_quasicoherentData_mathlib}
  \lean{SheafOfModules.QuasicoherentData}
  \mathlibok
  \textit{Provided by Mathlib
  (\texttt{Mathlib.Algebra.Category.ModuleCat.Sheaf.Quasicoherent}).}
  A sheaf of modules \(M\) on a scheme is \emph{quasi-coherent}
  (\(\mathrm{IsQuasicoherent}\,M\)) exactly when the type of \emph{quasi-coherence data}
  \(\mathrm{QuasicoherentData}\,M\) is nonempty. A datum
  \(q : \mathrm{QuasicoherentData}\,M\) consists of an index family \(\{X_i\}\) of opens covering the
  whole scheme together with, for each \(i\), a presentation
  \((M.\mathrm{over}\,X_i).\mathrm{Presentation}\) of the restriction of \(M\) to \(X_i\)
  (a cokernel-of-free-modules description on the site over \(X_i\)). Thus quasi-coherence packages a
  cover on which \(M\) is locally presented; on each member of the cover
  \cref{lem:isIso_fromTildeΓ_of_presentation_mathlib} applies.
\end{lemma}

% --- Mathlib dependency anchors for the gap1 slice-transport route ---

\begin{lemma}[Pullback of sheaves of modules along a morphism of schemes]
  \label{lem:modules_pullback_mathlib}
  \lean{AlgebraicGeometry.Scheme.Modules.pullback}
  \mathlibok
  \textit{Provided by Mathlib (\texttt{Mathlib.AlgebraicGeometry.Modules.Sheaf}).}
  For a morphism of schemes \(f : X \to Y\), Mathlib supplies the pullback functor
  \(f^{*} : Y.\mathrm{Modules} \to X.\mathrm{Modules}\) on sheaves of modules, left adjoint to
  pushforward, hence colimit-preserving, and carrying the unit sheaf to the unit sheaf. The project
  already uses this functor in the definition of local freeness. It is the geometric restriction
  through which a presentation is transported to an open subscheme.
\end{lemma}

\begin{lemma}[Restriction of sheaves of modules along an open immersion]
  \label{lem:modules_restrictFunctor_mathlib}
  \lean{AlgebraicGeometry.Scheme.Modules.restrictFunctor}
  \mathlibok
  \textit{Provided by Mathlib (\texttt{Mathlib.AlgebraicGeometry.Modules.Sheaf}).}
  For an open immersion \(f : X \to Y\), Mathlib supplies the restriction functor
  \(f^{\mathrm{res}} : Y.\mathrm{Modules} \to X.\mathrm{Modules}\), with
  \(\Gamma(f^{\mathrm{res}} M, U) = \Gamma(M, f(U))\) definitionally; it is a left adjoint (colimit
  preserving) and carries the unit sheaf to the unit sheaf. This is the geometric
  ``restrict to open \(U\)'' functor on \((\Spec R).\mathrm{Modules}\).
\end{lemma}

\begin{lemma}[Restriction is isomorphic to pullback]
  \label{lem:modules_restrictFunctorIsoPullback_mathlib}
  \lean{AlgebraicGeometry.Scheme.Modules.restrictFunctorIsoPullback}
  \mathlibok
  \textit{Provided by Mathlib (\texttt{Mathlib.AlgebraicGeometry.Modules.Sheaf}).}
  For an open immersion \(f\), the restriction functor (\cref{lem:modules_restrictFunctor_mathlib})
  is naturally isomorphic to the pullback functor (\cref{lem:modules_pullback_mathlib}):
  \(f^{\mathrm{res}} \cong f^{*}\). This identifies the restriction functor with the
  universal-property pullback, so a presentation transported along one transports along the other.
\end{lemma}

\begin{lemma}[Transport of a presentation along a colimit-preserving functor]
  \label{lem:presentation_map_mathlib}
  \lean{SheafOfModules.Presentation.map}
  \mathlibok
  \textit{Provided by Mathlib (\texttt{Mathlib.Algebra.Category.ModuleCat.Sheaf.Quasicoherent}).}
  Let \(F\) be a functor between categories of sheaves of modules that preserves colimits and is
  equipped with an isomorphism \(\eta : F(\mathrm{unit}_R) \cong \mathrm{unit}_S\) sending the unit
  sheaf to the unit sheaf. Then any presentation \(P : \mathrm{Presentation}\,M\) (a
  cokernel-of-free-modules description) induces a presentation
  \(P.\mathrm{map}\,F\,\eta : \mathrm{Presentation}\,(F(M))\) of the image. Applied to the pullback
  and restriction functors (\cref{lem:modules_pullback_mathlib},
  \cref{lem:modules_restrictFunctor_mathlib}), it is how a presentation of \(M\) over one site is
  carried to a presentation over another.
\end{lemma}

\begin{lemma}[Slice-presentation transport for quasi-coherence data]
  \label{lem:quasicoherentData_bind_mathlib}
  \lean{SheafOfModules.QuasicoherentData.bind}
  \mathlibok
  \textit{Provided by Mathlib (\texttt{Mathlib.Algebra.Category.ModuleCat.Sheaf.Quasicoherent}).}
  Mathlib's \(\mathrm{QuasicoherentData.bind}\) transports a presentation across an iterated
  Grothendieck slice: from quasi-coherence data on a cover \(\{X_i\}\) and, for each \(i\), data on a
  cover of \(M.\mathrm{over}\,X_i\), it assembles quasi-coherence data on the composite cover, by
  applying \(\mathrm{Presentation.map}\) (\cref{lem:presentation_map_mathlib}) along the equivalence
  \(\mathrm{pushforwardPushforwardEquivalence}(\mathrm{Over.iteratedSliceEquiv}\,\cdots)\). It is the
  worked template for the per-element presentation transport used in \cref{lem:over_restrict_iso} and
  \cref{lem:isIso_fromTildeΓ_basicOpen_of_quasicoherent}, handling the slice-site adjointness and
  sheafification typeclass synthesis.
\end{lemma}

\begin{lemma}[A presentation makes a sheaf of modules quasi-coherent]
  \label{lem:presentation_isQuasicoherent_mathlib}
  \lean{SheafOfModules.Presentation.isQuasicoherent}
  \mathlibok
  \textit{Provided by Mathlib
  (\texttt{Mathlib.Algebra.Category.ModuleCat.Sheaf.Quasicoherent}).}
  If a sheaf of \(R\)-modules \(M\) admits a presentation
  \(P : \mathrm{Presentation}\,M\), then \(M\) is quasi-coherent (\(\mathrm{IsQuasicoherent}\,M\)).
\end{lemma}

\begin{lemma}[Quasi-coherence is local: a cover on which $M$ is quasi-coherent makes $M$ quasi-coherent]
  \label{lem:isQuasicoherent_of_coversTop_mathlib}
  \lean{SheafOfModules.IsQuasicoherent.of_coversTop}
  \mathlibok
  \textit{Provided by Mathlib
  (\texttt{Mathlib.Algebra.Category.ModuleCat.Sheaf.Quasicoherent}).}
  Let \(M\) be a sheaf of \(R\)-modules and \(\{X_i\}\) an indexed family of objects of the site with
  \(\mathrm{CoversTop}\,\{X_i\}\). If the Grothendieck-slice restriction \(M.\mathrm{over}\,X_i\) is
  quasi-coherent for every \(i\), then \(M\) is quasi-coherent.
\end{lemma}

\begin{lemma}[Localization is flat]
  \label{lem:isLocalization_flat_mathlib}
  \lean{IsLocalization.flat}
  \mathlibok
  \textit{Provided by Mathlib (\texttt{Mathlib.RingTheory.Flat.Localization}).}
  If \(S\) is a localization of a commutative ring \(R\) at a multiplicative submonoid, then \(S\) is
  a flat \(R\)-module. Consequently localization is exact and commutes with finite limits; this is
  what lets the section-localization functor pass through the finite sheaf equalizer in the descent
  (\cref{lem:section_localization_descent}).
\end{lemma}

% --- The finite basic-open cover (topological precursor) ---

\begin{lemma}

\leanok
[Finite basic-open cover refining a quasi-coherence cover]
  \label{lem:exists_finite_basicOpen_cover_le_quasicoherentData}
  \lean{AlgebraicGeometry.exists_finite_basicOpen_cover_le_quasicoherentData}
  \uses{lem:isQuasicoherent_quasicoherentData_mathlib, lem:isBasis_basic_opens_mathlib}
  \textit{The topological finite-cover front of the gap1 transport: any quasi-coherence cover refines
  to a finite basic-open cover.}
  Let \(M\) be a sheaf of modules on \(\Spec R\) and let \(q\) be quasi-coherence data for \(M\)
  (\cref{lem:isQuasicoherent_quasicoherentData_mathlib}): a — possibly infinite — open cover
  \(\{q.X_i\}\) of \(\top\) together with a presentation of \(M.\mathrm{over}\,(q.X_i)\) on each
  member. Then there is a \emph{finite} subset \(t \subseteq R\) whose basic opens cover \(\Spec R\)
  (equivalently \(\mathrm{span}\,t = (1)\)) such that each \(D(r)\), \(r \in t\), is contained in some
  member \(q.X_i\) of the presentation cover.
\end{lemma}
\begin{proof}
  \uses{lem:isBasis_basic_opens_mathlib}
  The set \(G = \{\,r \in R : D(r) \le q.X_i \text{ for some } i\,\}\) spans the unit ideal. Indeed
  every prime of \(R\) lies in some member \(q.X_i\) by the covering property \(q.\mathrm{coversTop}\)
  (read off the \(\mathrm{Opens}\) Grothendieck-topology sieve), and the basic opens form a basis of
  the Zariski topology (\cref{lem:isBasis_basic_opens_mathlib}), so a basic open \(D(r) \subseteq
  q.X_i\) is a neighbourhood of that prime; hence \(\bigsqcup_{r \in G} D(r) = \top\), which is to say
  \(\mathrm{span}\,G = (1)\). A unit-ideal membership uses only finitely many of the generators,
  yielding a finite \(t \subseteq G\) with \(\mathrm{span}\,t = (1)\) and each \(D(r) \le q.X_i\) for
  \(r \in t\).
\end{proof}

% --- The over-site/open-subspace sheaf equivalence (bridge C, topological layer) ---

% The blocks in this sub-section record the project-local infrastructure that fills the explicit
% TODO of Mathlib's `Topology/Sheaves/Over.lean`: the equivalence of sheaf categories induced by
% the site equivalence `Opens.overEquivalence U`. Together with the three Mathlib dependency
% anchors below they form the topos-theoretic (purely topological) layer of bridge C
% (\cref{lem:over_restrict_iso}); the geometric identification of the structure sheaves is the
% layer above them.

\begin{lemma}[Slice of opens equivalent to opens of the open subspace]
  \label{lem:opens_overEquivalence_mathlib}
  \lean{TopologicalSpace.Opens.overEquivalence}
  \mathlibok
  \textit{Provided by Mathlib (\texttt{Mathlib.Topology.Sheaves.Over}).}
  For a topological space \(X\) and an open \(U \subseteq X\), the slice category of opens of
  \(X\) over \(U\) (the opens of \(X\) contained in \(U\), with inclusions) is equivalent to the
  opens of the subspace \(U\):
  \((\mathrm{Opens}\,X)_{/U} \simeq \mathrm{Opens}(U)\). Mathlib supplies this underlying
  equivalence of sites; it leaves the lift to sheaf categories as an explicit TODO.
\end{lemma}

\begin{lemma}[Sheaf categories of equivalent sites are equivalent]
  \label{lem:equivalence_sheafCongr_mathlib}
  \lean{CategoryTheory.Equivalence.sheafCongr}
  \mathlibok
  \textit{Provided by Mathlib (\texttt{Mathlib.CategoryTheory.Sites.Equivalence}).}
  Let \(e : C \simeq D\) be an equivalence of categories carrying Grothendieck topologies \(J\)
  on \(C\) and \(K\) on \(D\), and suppose \(e.\mathrm{inverse}\) is a dense subsite for
  \((K, J)\). Then for any target category \(A\) the sheaf categories are equivalent:
  \(\mathrm{Sheaf}(J, A) \simeq \mathrm{Sheaf}(K, A)\).
\end{lemma}

\begin{lemma}[Module sheaves transport across an equivalence of ringed sites]
  \label{lem:pushforwardPushforwardEquivalence_mathlib}
  \lean{SheafOfModules.pushforwardPushforwardEquivalence}
  \mathlibok
  \textit{Provided by Mathlib
  (\texttt{Mathlib.Algebra.Category.ModuleCat.Sheaf.PushforwardContinuous}).}
  Let \(e : C \simeq D\) be an equivalence of sites with both \(e.\mathrm{functor}\) and
  \(e.\mathrm{inverse}\) continuous, and let \(S\) (a sheaf of rings on \(J\)) and \(R\) (a sheaf
  of rings on \(K\)) be related by mutually inverse comparison maps across the continuous
  pushforwards. Then the categories of sheaves of modules are equivalent,
  \(\mathrm{SheafOfModules}\,R \simeq \mathrm{SheafOfModules}\,S\), compatibly with the
  underlying ring-sheaf identification. This is the module-level lift used in step~3 of
  \cref{lem:over_restrict_iso}.
\end{lemma}

\begin{lemma}\leanok
[The over-equivalence functor is cocontinuous]
  \label{lem:overEquivalence_functor_isCocontinuous}
  \lean{AlgebraicGeometry.overEquivalence_functor_isCocontinuous}
  \uses{lem:opens_overEquivalence_mathlib}
  \textit{Topological layer of bridge C; fills the \texttt{Topology/Sheaves/Over.lean} TODO.}
  The functor of \(\mathrm{Opens.overEquivalence}\,U\) (\cref{lem:opens_overEquivalence_mathlib})
  is cocontinuous (cover-lifting) from the \(U\)-slice \(J_X|_U\) of the opens Grothendieck
  topology of \(X\) to the opens Grothendieck topology \(J_U\) of the subspace \(U\).
\end{lemma}
\begin{proof}
  \uses{lem:opens_overEquivalence_mathlib}
  A sieve covers in the slice topology \(J_X|_U\) iff its image sieve covers in \(J_X\), and
  covering for the opens topology is pointwise: a sieve covers iff its members' opens cover the
  base set-theoretically. Given a point \(x\) of a slice object and a cover member in \(J_U\)
  containing the image of \(x\), transport it across the open embedding \(U \hookrightarrow X\),
  which is injective: the set-theoretic image of an open \(V \subseteq U\) is an open of \(X\)
  contained in the base, a neighbourhood of \(x\) lying in the image sieve. This witnesses the
  cover-lifting condition.
\end{proof}

\begin{lemma}\leanok
[The over-equivalence inverse is cocontinuous]
  \label{lem:overEquivalence_inverse_isCocontinuous}
  \lean{AlgebraicGeometry.overEquivalence_inverse_isCocontinuous}
  \uses{lem:opens_overEquivalence_mathlib}
  \textit{Topological layer of bridge C; fills the \texttt{Topology/Sheaves/Over.lean} TODO.}
  Symmetrically, the inverse of \(\mathrm{Opens.overEquivalence}\,U\) is cocontinuous from the
  opens topology \(J_U\) of the subspace \(U\) to the \(U\)-slice \(J_X|_U\).
\end{lemma}
\begin{proof}
  \uses{lem:opens_overEquivalence_mathlib}
  As in \cref{lem:overEquivalence_functor_isCocontinuous}, reduce to the pointwise covering
  criterion and the over-sieve image criterion; given a point and a slice-cover member, pull the
  member back along the embedding \(U \hookrightarrow X\) (preimage of an open of \(X\) is an open
  of the subspace) to obtain a cover member of \(J_U\) witnessing the lift.
\end{proof}

\begin{lemma}\leanok
[The over-equivalence inverse is a dense subsite]
  \label{lem:overEquivalence_inverse_isDenseSubsite}
  \lean{AlgebraicGeometry.overEquivalence_inverse_isDenseSubsite}
  \uses{lem:overEquivalence_functor_isCocontinuous, lem:overEquivalence_inverse_isCocontinuous}
  \textit{Topological layer of bridge C; the dense-subsite witness assembled from the two
  cocontinuities.}
  The inverse of \(\mathrm{Opens.overEquivalence}\,U\) is a dense subsite for \((J_U, J_X|_U)\).
  This is formal: an equivalence both of whose legs are cocontinuous
  (\cref{lem:overEquivalence_functor_isCocontinuous},
  \cref{lem:overEquivalence_inverse_isCocontinuous}) yields a dense subsite on its inverse.
\end{lemma}

\begin{lemma}\leanok
[The over-equivalence functor is continuous]
  \label{lem:overEquivalence_functor_isContinuous}
  \lean{AlgebraicGeometry.overEquivalence_functor_isContinuous}
  \uses{lem:overEquivalence_inverse_isCocontinuous}
  \textit{Topological layer of bridge C; needed for the module-level lift of step~3.}
  The functor of \(\mathrm{Opens.overEquivalence}\,U\) is continuous (cover-preserving) for
  \((J_X|_U, J_U)\). It is obtained from the cocontinuity of the inverse
  (\cref{lem:overEquivalence_inverse_isCocontinuous}) through the equivalence adjunction: the
  left adjoint of a cocontinuous functor is continuous.
\end{lemma}

\begin{lemma}\leanok
[The over-equivalence inverse is continuous]
  \label{lem:overEquivalence_inverse_isContinuous}
  \lean{AlgebraicGeometry.overEquivalence_inverse_isContinuous}
  \uses{lem:overEquivalence_functor_isCocontinuous}
  \textit{Topological layer of bridge C; needed for the module-level lift of step~3.}
  Symmetrically, the inverse of \(\mathrm{Opens.overEquivalence}\,U\) is continuous for
  \((J_U, J_X|_U)\), obtained from the cocontinuity of the functor
  (\cref{lem:overEquivalence_functor_isCocontinuous}) through the equivalence adjunction.
\end{lemma}

\begin{definition}\leanok
[The over-site/open-subspace sheaf equivalence]
  \label{def:overEquivalence_sheafCongr}
  \lean{AlgebraicGeometry.overEquivalence_sheafCongr}
  \uses{lem:opens_overEquivalence_mathlib, lem:equivalence_sheafCongr_mathlib,
    lem:overEquivalence_inverse_isDenseSubsite}
  \textit{Keystone of the topological layer of bridge C.}
  For any category \(A\), the site equivalence \(\mathrm{Opens.overEquivalence}\,U\)
  (\cref{lem:opens_overEquivalence_mathlib}) lifts to an equivalence of sheaf categories
  \[
    \mathrm{Sheaf}(J_X|_U,\, A) \;\simeq\; \mathrm{Sheaf}(J_U,\, A),
  \]
  obtained by feeding the dense-subsite witness
  (\cref{lem:overEquivalence_inverse_isDenseSubsite}) to the general sheaf-congruence
  construction (\cref{lem:equivalence_sheafCongr_mathlib}). This is exactly the equivalence left
  as a TODO in Mathlib's \texttt{Topology/Sheaves/Over.lean}, and it is the object across which
  \cref{lem:over_restrict_iso} transports the structure sheaf and the module \(M\).
\end{definition}

% --- (C, step 3) The module-level equivalence and its functor identification ---

% The two blocks below are the module-category layer assembled on top of the topological
% sheaf-congruence def:overEquivalence_sheafCongr: the step-3 equivalence of categories of sheaves of
% modules (def:over_restrict_equiv) and the step-4 functor identification (lem:over_restrict_functor_iso)
% whose object component is the slice-touching bridge lem:over_restrict_iso below. They are
% project-bespoke infrastructure (no external source).

\begin{definition}

\leanok
[The slice-to-geometric module-category equivalence (gap1, C, step 3)]
  \label{def:over_restrict_equiv}
  \lean{AlgebraicGeometry.Scheme.Modules.overRestrictEquiv}
  \uses{def:overEquivalence_sheafCongr, lem:opens_overEquivalence_mathlib,
    lem:pushforwardPushforwardEquivalence_mathlib, lem:overEquivalence_functor_isContinuous,
    lem:overEquivalence_inverse_isContinuous}
  \textit{Module-level lift of the topological bridge C; the step-3 equivalence on which the
  slice-touching isomorphism is built.}
  Let \(X\) be a scheme and \(U \subseteq X\) an open. The category of sheaves of modules over the
  sliced structure sheaf \(\mathcal{O}_X.\mathrm{over}\,U\) on the slice site \(J.\mathrm{over}\,U\) is
  equivalent to the category \(U.\mathrm{toScheme}.\mathrm{Modules}\) of sheaves of modules on the open
  subscheme \(U.\mathrm{toScheme}\):
  \[
    \mathrm{SheafOfModules}\bigl(\mathcal{O}_X.\mathrm{over}\,U\bigr)
      \;\simeq\; U.\mathrm{toScheme}.\mathrm{Modules}.
  \]
  This is the module-category lift of the topological sheaf equivalence
  \cref{def:overEquivalence_sheafCongr}.
\end{definition}
\begin{proof}
  \uses{def:overEquivalence_sheafCongr, lem:opens_overEquivalence_mathlib,
    lem:pushforwardPushforwardEquivalence_mathlib, lem:overEquivalence_functor_isContinuous,
    lem:overEquivalence_inverse_isContinuous}
  Apply the module-sheaf transport \(\mathrm{pushforwardPushforwardEquivalence}\)
  (\cref{lem:pushforwardPushforwardEquivalence_mathlib}) to the site equivalence
  \(\mathrm{Opens.overEquivalence}\,U\) (\cref{lem:opens_overEquivalence_mathlib}), whose two legs are
  both continuous (\cref{lem:overEquivalence_functor_isContinuous},
  \cref{lem:overEquivalence_inverse_isContinuous}). The transport consumes a pair of mutually inverse
  ring-sheaf comparison morphisms between the sliced structure sheaf
  \(\mathcal{O}_X.\mathrm{over}\,U\) and the geometric structure sheaf
  \(\mathcal{O}_{U.\mathrm{toScheme}}\); for these one takes the whiskered (co)unit of the site
  equivalence \(\mathrm{Opens.overEquivalence}\,U\) into the structure presheaf, paired with the
  identity. The identity is legal precisely because the geometric structure sheaf is
  the transport of the sliced one across the ring-valued instance of
  \cref{def:overEquivalence_sheafCongr} (this is step~2 of the bridge, which holds by
  construction). Orienting the resulting equivalence so that its functor runs from the slice
  side to the geometric side gives the stated equivalence.
\end{proof}

\begin{lemma}

\leanok
[Functor identification of the slice-to-geometric equivalence (gap1, C, step 4)]
  \label{lem:over_restrict_functor_iso}
  \lean{AlgebraicGeometry.Scheme.Modules.overRestrictFunctorIso}
  \uses{def:over_restrict_equiv, lem:modules_restrictFunctor_mathlib,
    lem:opens_overEquivalence_mathlib}
  \textit{The composite of the slice functor with the step-3 equivalence is the geometric restriction
  functor.}
  With \(X\) and \(U \subseteq X\) as above, the composite of the abstract Grothendieck-slice functor
  \(M \mapsto M.\mathrm{over}\,U\) (pushforward along the identity of the sliced ring sheaf) with the
  functor of the step-3 equivalence (\cref{def:over_restrict_equiv}) is naturally isomorphic to the
  geometric restriction functor along the open immersion \(U.\iota\):
  \[
    \bigl(\mathrm{pushforward}\,\mathbf{1}\bigr) \;\circ\; (\mathrm{overRestrictEquiv}\,U).\mathrm{functor}
      \;\cong\; \mathrm{restrictFunctor}\,U.\iota
  \]
  (\cref{lem:modules_restrictFunctor_mathlib}).
\end{lemma}
\begin{proof}
  \uses{def:over_restrict_equiv, lem:modules_restrictFunctor_mathlib,
    lem:opens_overEquivalence_mathlib}
  Both sides are pushforwards of sheaves of modules along the \emph{same} opens functor. Indeed the
  opens functor of the open immersion factors as
  \[
    U.\iota.\mathrm{opensFunctor}
      \;=\; (\mathrm{Opens.overEquivalence}\,U).\mathrm{inverse}\;\circ\;\mathrm{Over.forget}\,U
  \]
  (\cref{lem:opens_overEquivalence_mathlib}), an equality that holds by construction, and
  \(\mathrm{restrictFunctor}\,U.\iota\) (\cref{lem:modules_restrictFunctor_mathlib}) is precisely
  pushforward along \(U.\iota.\mathrm{opensFunctor}\). The left-hand composite is pushforward along the
  factored functor on the right. Identify the two via the composition law for module-sheaf
  pushforward (\(\mathrm{pushforwardComp}\)) followed by congruence of pushforward along the equal ring
  morphisms (\(\mathrm{pushforwardCongr}\)); the required ring-morphism equality is routine.
\end{proof}

% --- (C) The single slice-touching bridge ---

\begin{lemma}\leanok
[The Grothendieck-slice restriction equals the geometric restriction (gap1, C)]
  \label{lem:over_restrict_iso}
  \lean{AlgebraicGeometry.Scheme.Modules.overRestrictIso}
  \uses{lem:modules_restrictFunctor_mathlib, lem:modules_restrictFunctorIsoPullback_mathlib,
    lem:isLocalization_basicOpen_mathlib, def:overEquivalence_sheafCongr,
    lem:overEquivalence_functor_isContinuous, lem:overEquivalence_inverse_isContinuous,
    lem:opens_overEquivalence_mathlib, lem:pushforwardPushforwardEquivalence_mathlib}
  % NOTE: RESOLVED and axiom-clean. AlgebraicGeometry.Scheme.Modules.overRestrictIso
  % now exists (#print axioms = {propext, Classical.choice, Quot.sound}); gap1 bridge C
  % (steps 1->4) is fully closed. The step-2 geometric ring-sheaf identification collapsed to
  % `rfl` (U.toScheme.ringCatSheaf is definitionally the transport of the sliced sheaf, via
  % toScheme_presheaf_obj/map). The statement is routed through the step-3 equivalence functor:
  % overRestrictIso : (overRestrictEquiv U).functor.obj (M.over U) ≅ (restrictFunctor U.ι).obj M.
  \textit{The one lemma that must touch the Grothendieck slice site; everything downstream of it is
  geometric.}
  Let \(X\) be a scheme and \(U \subseteq X\) an open. The \emph{abstract Grothendieck-slice}
  restriction \(M.\mathrm{over}\,U\) — an object of the category of sheaves of modules over the
  sliced structure sheaf \(\mathcal{O}_X.\mathrm{over}\,U\) on the slice site \(J.\mathrm{over}\,U\) —
  is canonically isomorphic to the \emph{geometric} restriction
  \((\mathrm{restrictFunctor}\,U.\iota).\mathrm{obj}\,M\)
  (\cref{lem:modules_restrictFunctor_mathlib}), equivalently to the pullback \(U.\iota^{*} M\)
  (\cref{lem:modules_restrictFunctorIsoPullback_mathlib}), as a sheaf of modules on the small Zariski
  site of the open subscheme \(U.\mathrm{toScheme}\). The isomorphism is induced by the equivalence
  of sites \(J.\mathrm{over}\,U \simeq \mathrm{Opens}(U.\mathrm{toScheme})\) coming from the opens
  functor of the open immersion \(U.\iota\), under which the sliced sheaf of rings
  \(\mathcal{O}_X.\mathrm{over}\,U\) corresponds to the structure sheaf of \(U.\mathrm{toScheme}\).
  For \(X = \Spec R\) and \(U = D(r)\) a basic open, the target category is
  \((\Spec R_r).\mathrm{Modules}\) under the affine identification \(D(r) \cong \Spec R_r\)
  (\cref{lem:isLocalization_basicOpen_mathlib}).
\end{lemma}
\begin{proof}
  \uses{def:overEquivalence_sheafCongr, lem:overEquivalence_functor_isContinuous,
    lem:overEquivalence_inverse_isContinuous, lem:opens_overEquivalence_mathlib,
    lem:pushforwardPushforwardEquivalence_mathlib, lem:modules_restrictFunctor_mathlib,
    lem:modules_restrictFunctorIsoPullback_mathlib, lem:isLocalization_basicOpen_mathlib}
  The isomorphism is built in four steps, ascending from the purely topological layer to the
  module-level statement.

  \emph{Step 1 (topological layer --- done).} The opens functor of the open immersion
  \(U.\iota\) realizes the slice of the opens site of \(X\) over \(U\) as the opens site of the
  subspace \(U\); on sheaf categories this is the equivalence
  \(\mathrm{Sheaf}(J_X|_U,\,A) \simeq \mathrm{Sheaf}(J_U,\,A)\) of
  \cref{def:overEquivalence_sheafCongr}, for every target category \(A\). This is the
  topos-theoretic layer, assembled from the cocontinuity/continuity of both legs of
  \(\mathrm{Opens.overEquivalence}\,U\) (\cref{lem:opens_overEquivalence_mathlib}).

  \emph{Step 2 (geometric ring-sheaf identification --- the current obstacle).} Identify the
  sliced structure sheaf \(\mathcal{O}_X.\mathrm{over}\,U\) with the structure sheaf
  \(\mathcal{O}_{U}\) of the open subscheme \(U.\mathrm{toScheme}\), transported across the
  \(\mathrm{RingCat}\)-valued instance of \cref{def:overEquivalence_sheafCongr}; this is an
  isomorphism of sheaves of rings. The bridge is the factorization of the opens functor of the
  immersion,
  \[
    U.\iota.\mathrm{opensFunctor}
      \;=\; (\mathrm{Opens.overEquivalence}\,U).\mathrm{inverse}
        \;\circ\; \mathrm{Over.forget}\,U ,
  \]
  under which the structure-sheaf comparison reduces to the standard open-immersion
  structure-sheaf identifications; recall that \(\mathrm{restrictFunctor}\,U.\iota\)
  (\cref{lem:modules_restrictFunctor_mathlib}) is precisely pushforward along
  \(U.\iota.\mathrm{opensFunctor}\). This ring-sheaf identification is the remaining geometric
  obstacle.

  \emph{Step 3 (lift to modules).} Given the step-2 identification of the underlying ring sheaves
  and the continuity of both legs (\cref{lem:overEquivalence_functor_isContinuous},
  \cref{lem:overEquivalence_inverse_isContinuous}), lift the sheaf-of-rings comparison to an
  equivalence of categories of sheaves of modules via
  \(\mathrm{pushforwardPushforwardEquivalence}\)
  (\cref{lem:pushforwardPushforwardEquivalence_mathlib}). Under this equivalence the sliced
  module \(M.\mathrm{over}\,U\) corresponds to the geometric restriction
  \((\mathrm{restrictFunctor}\,U.\iota).\mathrm{obj}\,M\).

  \emph{Step 4 (compose).} Compose with
  \(\mathrm{restrictFunctorIsoPullback}\) (\cref{lem:modules_restrictFunctorIsoPullback_mathlib})
  to land the isomorphism \(M.\mathrm{over}\,U \cong M.\mathrm{restrict}\,U.\iota \cong
  U.\iota^{*}M\). The two sides live a priori in different module categories, so the literal Lean
  statement of \(\mathrm{overRestrictIso}\) is phrased through the step-3 equivalence functor; for
  \(X = \Spec R\) and \(U = D(r)\) the target is \((\Spec R_r).\mathrm{Modules}\) under
  \(D(r) \cong \Spec R_r\) (\cref{lem:isLocalization_basicOpen_mathlib}).
\end{proof}

% --- (C, step 4') The pullback form of the slice-to-geometric isomorphism ---

\begin{lemma}

\leanok
[The slice-to-geometric isomorphism in pullback form (gap1, C, step 4')]
  \label{lem:over_restrict_pullback_iso}
  \lean{AlgebraicGeometry.Scheme.Modules.overRestrictPullbackIso}
  \uses{lem:over_restrict_iso, lem:modules_restrictFunctorIsoPullback_mathlib,
    lem:modules_pullback_mathlib}
  \textit{The pullback packaging of the slice-touching bridge; the form consumed by the P1 transport.}
  With \(X\) and \(U \subseteq X\) as above and \(M\) a sheaf of modules on \(X\), the transport of the
  abstract Grothendieck-slice restriction \(M.\mathrm{over}\,U\) under the step-3 equivalence functor
  (\cref{def:over_restrict_equiv}) is canonically isomorphic to the inverse-image (pullback) of \(M\)
  along the open immersion \(U.\iota\):
  \[
    (\mathrm{overRestrictEquiv}\,U).\mathrm{functor}.\mathrm{obj}\,(M.\mathrm{over}\,U)
      \;\cong\; (U.\iota^{*})\,M
  \]
  (\cref{lem:modules_pullback_mathlib}). This is the pullback packaging of the slice-touching bridge
  \cref{lem:over_restrict_iso}; it is exactly the form in which a presentation of
  \(M.\mathrm{over}\,U\) is transported into a presentation of the geometric pullback \(U.\iota^{*}M\),
  and it is the ingredient consumed by the per-element local-tilde transport
  \cref{lem:isIso_fromTildeΓ_basicOpen_of_quasicoherent}.
\end{lemma}
\begin{proof}
  \uses{lem:over_restrict_iso, lem:modules_restrictFunctorIsoPullback_mathlib}
  Compose the slice-to-geometric isomorphism \cref{lem:over_restrict_iso}, which identifies the
  transported slice restriction with the geometric restriction
  \((\mathrm{restrictFunctor}\,U.\iota).\mathrm{obj}\,M\), with the component at \(M\) of the natural
  isomorphism \(\mathrm{restrictFunctorIsoPullback}\,U.\iota\)
  (\cref{lem:modules_restrictFunctorIsoPullback_mathlib}) identifying the geometric restriction with
  the pullback \(U.\iota^{*}M\).
\end{proof}

% --- (P1 prep) Slice→geometric Presentation-transport machinery ---
% These three project-bespoke helpers transport a presentation across the step-3
% slice equivalence into a presentation of the geometric pullback; they feed the
% P1 keystone (lem:isIso_fromTildeΓ_basicOpen_of_quasicoherent). No external source.

\begin{lemma}\leanok
[Unit-comparison iso from an iso ring comparison]
  \label{lem:isIso_unitToPushforwardObjUnit_of_isIso}
  \lean{AlgebraicGeometry.isIso_unitToPushforwardObjUnit_of_isIso'}
  \textit{Project-bespoke auxiliary (general continuous-functor fact).}
  Let \(Fu : C \to D\) be a continuous functor of sites and \(\psi : Fu^{*}S \to R\) (equivalently a
  comparison of sheaves of rings) with \(\psi\) an isomorphism. Then the canonical unit comparison
  \(\mathrm{unitToPushforwardObjUnit}\,\psi : \mathbf{1}_S \to (\mathrm{pushforward}\,\psi).\mathrm{obj}\,
  \mathbf{1}_R\) is an isomorphism.
\end{lemma}
\begin{proof}
  Reflecting through the forgetful passage to presheaves, the component of the unit comparison over each
  object \(V\) is the functor-image of \(\psi.\mathrm{app}\,V\); since \(\psi\) is an isomorphism each
  such component is, hence so is the unit comparison. (This is the \(\mathrm{IsIso}\,\psi\)-driven
  criterion, in place of Mathlib's finality-hypothesis form.)
\end{proof}

\begin{definition}\leanok
[Step-3 equivalence sends the unit module to the unit module]
  \label{def:over_restrict_unit_iso}
  \lean{AlgebraicGeometry.Scheme.Modules.overRestrictUnitIso}
  \uses{def:over_restrict_equiv, lem:isIso_unitToPushforwardObjUnit_of_isIso}
  \textit{The step-3 slice equivalence functor preserves the structure-sheaf (unit) module.}
  Let \(U \subseteq X\) be an open subscheme. The functor of the step-3 slice
  equivalence (\cref{def:over_restrict_equiv}) carries the unit
  (structure-sheaf) module \(\mathbf{1}\) of the over-site
  \(X.\mathrm{ringCatSheaf}.\mathrm{over}\,U\) to the unit module of
  \(U.\mathrm{toScheme}.\mathrm{ringCatSheaf}\):
  \[
    (\mathrm{overRestrictEquiv}\,U).\mathrm{functor}.\mathrm{obj}\,
      \bigl(\mathbf{1}\,(X.\mathrm{ringCatSheaf}.\mathrm{over}\,U)\bigr)
      \;\cong\; \mathbf{1}\,(U.\mathrm{toScheme}.\mathrm{ringCatSheaf}) .
  \]
\end{definition}
\begin{proof}
  \uses{def:over_restrict_equiv}
  By construction the functor of the equivalence is the pushforward
  \(\mathrm{pushforward}\,\psi_0\) of sheaves of modules along the
  over-site/open-site comparison, where \(\psi_0 = \mathrm{homMk}(\mathbf{1})\)
  is the identity ring comparison. The canonical unit comparison
  \(\mathrm{unitToPushforwardObjUnit}\,\psi_0 : \mathbf{1}_S \to
  (\mathrm{pushforward}\,\psi_0).\mathrm{obj}\,\mathbf{1}_R\) is an isomorphism
  whenever the underlying ring comparison \(\psi_0\) is one: reflecting through the
  forgetful passage to presheaves (\(\mathrm{toSheaf}/\mathrm{sheafToPresheaf}\)),
  its component over each object \(V\) is the functor-image of
  \(\psi_0.\mathrm{app}\,V\), hence an isomorphism (this \(\mathrm{IsIso}\,\psi\)
  driven criterion is supplied by the auxiliary
  \(\mathrm{isIso\_unitToPushforwardObjUnit\_of\_isIso}'\)). Here
  \(\psi_0 = \mathrm{homMk}(\mathbf{1})\) is the identity, so the unit comparison
  is an isomorphism, and its inverse is the asserted isomorphism. (Mathlib's
  isomorphism statement for \(\mathrm{unitToPushforwardObjUnit}\) instead requires
  a finality hypothesis on the comparison functor; the \(\mathrm{IsIso}\,\psi\)
  driven form used here is the one the slice equivalence needs.)
\end{proof}

\begin{definition}\leanok
[Transport of a slice presentation to a geometric presentation]
  \label{def:over_restrict_presentation}
  \lean{AlgebraicGeometry.Scheme.Modules.overRestrictPresentation}
  \uses{def:over_restrict_unit_iso, lem:over_restrict_pullback_iso, lem:presentation_map_mathlib}
  \textit{A slice presentation yields a presentation of the geometric pullback.}
  Let \(U \subseteq X\) be open and \(M\) a sheaf of modules on \(X\). A
  presentation of the abstract Grothendieck-slice restriction \(M.\mathrm{over}\,U\)
  yields a presentation of the geometric pullback \((U.\iota^{*})\,M\):
  \[
    (M.\mathrm{over}\,U).\mathrm{Presentation}
      \;\longrightarrow\;
    \bigl((U.\iota^{*})\,M\bigr).\mathrm{Presentation} .
  \]
\end{definition}
\begin{proof}
  \uses{def:over_restrict_unit_iso, lem:over_restrict_pullback_iso, lem:presentation_map_mathlib}
  Apply the functorial transport of presentations \(\mathrm{Presentation.map}\)
  (\cref{lem:presentation_map_mathlib}) along the step-3 slice equivalence functor
  \((\mathrm{overRestrictEquiv}\,U).\mathrm{functor}\), using
  \cref{def:over_restrict_unit_iso} to identify the image of the structure-sheaf
  (unit) module with the target unit module; this lands a presentation of the
  transported slice module
  \((\mathrm{overRestrictEquiv}\,U).\mathrm{functor}.\mathrm{obj}\,(M.\mathrm{over}\,U)\).
  Finally transport across the pullback packaging isomorphism
  \((\mathrm{overRestrictEquiv}\,U).\mathrm{functor}.\mathrm{obj}\,(M.\mathrm{over}\,U)
  \cong (U.\iota^{*})\,M\) (\cref{lem:over_restrict_pullback_iso}) via
  \(\mathrm{Presentation.ofIsIso}\), giving a presentation of the geometric
  pullback \((U.\iota^{*})\,M\).
\end{proof}

\begin{definition}\leanok
[Per-cover-member geometric presentation from quasi-coherence data]
  \label{def:presentation_pullback_iota_of_quasicoherentData}
  \lean{AlgebraicGeometry.Scheme.Modules.presentationPullbackιOfQuasicoherentData}
  \uses{def:over_restrict_presentation, lem:quasicoherentData_bind_mathlib}
  \textit{Quasi-coherence data gives a presentation on each cover member's pullback.}
  Let \(M\) be a sheaf of modules on \(X\), let \(q\) be quasi-coherence data for
  \(M\), and let \(i\) be a cover index. There is a presentation of the geometric
  pullback of \(M\) to the open subscheme \(q.X_i\):
  \[
    \bigl((q.X_i).\iota^{*}\,M\bigr).\mathrm{Presentation} .
  \]
\end{definition}
\begin{proof}
  \uses{def:over_restrict_presentation, lem:quasicoherentData_bind_mathlib}
  The datum \(q\) packages, for each index \(i\), a slice presentation
  \(q.\mathrm{presentation}\,i : (M.\mathrm{over}\,(q.X_i)).\mathrm{Presentation}\)
  over the cover member \(q.X_i\) (the per-member presentations whose binding is
  recorded by \(\mathrm{QuasicoherentData}\), \cref{lem:quasicoherentData_bind_mathlib}).
  Feeding this slice presentation to \cref{def:over_restrict_presentation} at
  \(U = q.X_i\) produces the desired presentation of the geometric pullback
  \((q.X_i).\iota^{*}\,M\).
\end{proof}

% --- (P1 prep) Affine-open presentation transport: slice → geometrize → affine-identify ---
% Six project-bespoke Scheme.Modules helpers (no external source — Mathlib presentation/pullback
% bookkeeping): they carry the global presentation of a cover member down to an arbitrary open W,
% then transport it across the affine identification W ≅ Spec Γ(Z,W). Together they assemble the
% general keystone lem:isIso_fromTildeΓ_presentationPullback, of which the basic-open P1 keystone
% lem:isIso_fromTildeΓ_basicOpen_of_quasicoherent is the affine instance W = (q.X i).ι ⁻¹ᵁ D(r).

\begin{definition}

\leanok
[Presentation of the geometric restriction to an open of a cover member]
  \label{def:presentation_pullback_iota_restrict}
  \lean{AlgebraicGeometry.Scheme.Modules.presentationPullbackιRestrict}
  \uses{def:presentation_pullback_iota_of_quasicoherentData, def:over_restrict_presentation,
    lem:presentation_map_mathlib}
  \textit{A global presentation on a cover member restricts to a presentation on any open of it.}
  Let \(M\) be a sheaf of modules on \(X\) with quasi-coherence data \(q\), let \(i\) be a cover
  index, and let \(W \subseteq Z\) be \emph{any} open of the cover-member subscheme
  \(Z := (q.X_i).\mathrm{toScheme}\). Then the iterated geometric restriction
  \((W.\iota^{*})\,\bigl((q.X_i).\iota^{*}\,M\bigr)\) of \(M\) admits a
  \(\mathrm{SheafOfModules.Presentation}\) on the open subscheme \(W.\mathrm{toScheme}\).
\end{definition}
\begin{proof}
  \uses{def:presentation_pullback_iota_of_quasicoherentData, def:over_restrict_presentation,
    lem:presentation_map_mathlib}
  Write \(N := (q.X_i).\iota^{*}\,M\), which carries the global presentation
  \cref{def:presentation_pullback_iota_of_quasicoherentData} on \(Z\). Slice this global presentation
  down to the Grothendieck-slice \(N.\mathrm{over}\,W\) by a single application of
  \(\mathrm{Presentation.map}\) (\cref{lem:presentation_map_mathlib}) along the over-functor
  \(\mathrm{pushforward}\,(\mathbf{1})\) — no iterated-slice equivalence is needed since \(N\) already
  lives on the genuine scheme \(Z\). Then geometrize this slice presentation via
  \cref{def:over_restrict_presentation} at \(U = W\), yielding a presentation of the geometric
  restriction \((W.\iota^{*})\,N\) on \(W.\mathrm{toScheme}\).
\end{proof}

\begin{definition}

\leanok
[The opens functor of an iso of schemes is an equivalence of opens sites]
  \label{def:opens_map_equiv_of_iso}
  \lean{AlgebraicGeometry.Scheme.Modules.opensMapEquivOfIso}
  % NOTE: the Mathlib pseudofunctoriality isos Opens.mapComp / Opens.mapId / Opens.mapIso and
  % Scheme.Hom.comp_base carry no blueprint anchor; they are named in prose rather than \uses'd.
  \textit{The inverse-image opens functor of a scheme isomorphism is an equivalence.}
  For an isomorphism of schemes \(\varphi : Y \cong Z\), the inverse-image functor
  \(\mathrm{Opens.map}\,\varphi_{\mathrm{inv}}^{\mathrm{base}} :
  \mathrm{Opens}\,Y \to \mathrm{Opens}\,Z\) on opens sites is an equivalence of categories, with
  inverse \(\mathrm{Opens.map}\,\varphi_{\mathrm{hom}}^{\mathrm{base}}\).
\end{definition}
\begin{proof}
  The continuous map \(\varphi_{\mathrm{inv}}^{\mathrm{base}}\) underlying \(\varphi.\mathrm{inv}\) is
  a homeomorphism, so its opens map is an equivalence. Concretely, the unit and counit isomorphisms
  are assembled from the pseudofunctoriality isos \(\mathrm{Opens.mapComp}\) / \(\mathrm{Opens.mapId}\)
  / \(\mathrm{Opens.mapIso}\), with the composite identities
  \(\varphi_{\mathrm{hom}}^{\mathrm{base}} \circ \varphi_{\mathrm{inv}}^{\mathrm{base}} = \mathrm{id}\)
  and \(\varphi_{\mathrm{inv}}^{\mathrm{base}} \circ \varphi_{\mathrm{hom}}^{\mathrm{base}} =
  \mathrm{id}\) coming from \(\varphi.\mathrm{hom\_inv\_id}\) / \(\varphi.\mathrm{inv\_hom\_id}\) via
  \(\mathrm{Scheme.Hom.comp\_base}\).
\end{proof}

\begin{lemma}

\leanok
[The opens functor of an iso of schemes is final]
  \label{lem:opens_map_final_of_scheme_iso}
  \lean{AlgebraicGeometry.Scheme.Modules.opensMap_final_of_schemeIso}
  \uses{def:opens_map_equiv_of_iso}
  \textit{The inverse-image opens functor of a scheme isomorphism is final.}
  For an isomorphism of schemes \(\varphi : Y \cong Z\), the inverse-image opens functor
  \(\mathrm{Opens.map}\,\varphi_{\mathrm{inv}}^{\mathrm{base}}\) is \(\mathrm{Final}\).
\end{lemma}
\begin{proof}
  \uses{def:opens_map_equiv_of_iso}
  An equivalence of categories is final; the functor is an equivalence by
  \cref{def:opens_map_equiv_of_iso}.
\end{proof}

\begin{definition}

\leanok
[Pullback along an iso of schemes sends the unit module to the unit module]
  \label{def:pullback_scheme_iso_unit_iso}
  \lean{AlgebraicGeometry.Scheme.Modules.pullbackSchemeIsoUnitIso}
  \uses{def:opens_map_equiv_of_iso, lem:opens_map_final_of_scheme_iso}
  % NOTE: the Mathlib canonical comparison SheafOfModules.pullbackObjUnitToUnit carries no blueprint
  % anchor; it is named in prose rather than \uses'd.
  \textit{Pullback along a scheme isomorphism preserves the structure-sheaf module.}
  For an isomorphism of schemes \(\varphi : Y \cong Z\), the pullback functor along
  \(\varphi.\mathrm{inv} : Z \to Y\) carries the structure-sheaf (unit) module of \(Y\) to the unit
  module of \(Z\):
  \[
    \bigl(\mathrm{pullback}\,\varphi.\mathrm{inv}\bigr).\mathrm{obj}\,
      \bigl(\mathbf{1}\,Y\bigr) \;\cong\; \mathbf{1}\,Z .
  \]
\end{definition}
\begin{proof}
  \uses{def:opens_map_equiv_of_iso, lem:opens_map_final_of_scheme_iso}
  The canonical comparison \(\mathrm{pullbackObjUnitToUnit}\) is an isomorphism whenever the site
  functor \(\mathrm{Opens.map}\,\varphi_{\mathrm{inv}}^{\mathrm{base}}\) is \(\mathrm{Final}\), which
  holds by \cref{lem:opens_map_final_of_scheme_iso} (itself from the equivalence
  \cref{def:opens_map_equiv_of_iso}). Its inverse is the asserted unit isomorphism.
\end{proof}

\begin{definition}

\leanok
[Pullback along an open immersion sends the unit module to the unit module]
  \label{def:pullback_open_immersion_unit_iso}
  \lean{AlgebraicGeometry.Scheme.Modules.pullbackOpenImmersionUnitIso}
  \uses{def:pullback_scheme_iso_unit_iso}
  % NOTE: the Mathlib canonical comparison SheafOfModules.pullbackObjUnitToUnit carries no blueprint
  % anchor; it is named in prose rather than \uses'd.
  \textit{Pullback along an open immersion preserves the structure-sheaf module.}
  Generalising \cref{def:pullback_scheme_iso_unit_iso} from isomorphisms to open immersions: for an
  open immersion \(k : Y \to X\), the pullback functor along \(k\) carries the structure-sheaf (unit)
  module of \(X\) to the unit module of \(Y\):
  \[
    \bigl(\mathrm{pullback}\,k\bigr).\mathrm{obj}\,\bigl(\mathbf{1}\,X\bigr) \;\cong\; \mathbf{1}\,Y .
  \]
\end{definition}
\begin{proof}
  \uses{def:pullback_scheme_iso_unit_iso}
  The canonical comparison \(\mathrm{pullbackObjUnitToUnit}\) is an isomorphism whenever the site functor
  \(\mathrm{Opens.map}\,k_{\mathrm{base}}\) is \(\mathrm{Final}\). Since \(k\) is an open immersion,
  \(k_{\mathrm{base}}\) is an open map, so \(\mathrm{Opens.map}\,k_{\mathrm{base}}\) is a right adjoint
  (the open-map adjunction) and hence \(\mathrm{Final}\). Its inverse is the asserted unit isomorphism.
\end{proof}

\begin{definition}

\leanok
[A presentation transports across pullback by an iso of schemes]
  \label{def:presentation_pullback_of_scheme_iso}
  \lean{AlgebraicGeometry.Scheme.Modules.presentationPullbackOfSchemeIso}
  \uses{def:pullback_scheme_iso_unit_iso, lem:presentation_map_mathlib,
    lem:modules_pullback_mathlib}
  \textit{Pullback by a scheme isomorphism carries a presentation to a presentation.}
  Given an isomorphism of schemes \(\varphi : Y \cong Z\) and a presentation of a sheaf of modules
  \(N\) on \(Y\), the geometric pullback \((\mathrm{pullback}\,\varphi.\mathrm{inv}).\mathrm{obj}\,N\)
  of \(N\) to \(Z\) admits a presentation.
\end{definition}
\begin{proof}
  \uses{def:pullback_scheme_iso_unit_iso, lem:presentation_map_mathlib,
    lem:modules_pullback_mathlib}
  Apply \(\mathrm{Presentation.map}\) (\cref{lem:presentation_map_mathlib}) along the pullback functor
  \(\mathrm{pullback}\,\varphi.\mathrm{inv}\) (\cref{lem:modules_pullback_mathlib}), which preserves
  colimits as the left adjoint of the pullback--pushforward adjunction
  (\(\mathrm{leftAdjoint\_preservesColimits}\)), feeding it the unit identification
  \cref{def:pullback_scheme_iso_unit_iso}. Applied with \(\varphi\) the affine identification
  \(\mathrm{IsAffineOpen.isoSpec}\), this is the step that moves a presentation onto a genuine
  \(\Spec\).
\end{proof}

\begin{lemma}

\leanok
[Quasi-coherent restricts to a tilde on every affine open of a cover member (gap1, P1, general keystone)]
  \label{lem:isIso_fromTildeΓ_presentationPullback}
  \lean{AlgebraicGeometry.Scheme.Modules.isIso_fromTildeΓ_presentationPullback}
  \uses{def:presentation_pullback_iota_restrict, def:presentation_pullback_of_scheme_iso,
    lem:isIso_fromTildeΓ_of_presentation_mathlib}
  % NOTE: the Mathlib affine identification IsAffineOpen.isoSpec carries no blueprint anchor; it is
  % named in prose rather than \uses'd. This lemma is STRICTLY MORE GENERAL than the basic-open
  % keystone lem:isIso_fromTildeΓ_basicOpen_of_quasicoherent (which is its instance
  % W = (q.X i).ι ⁻¹ᵁ D(r)); it is the form the gap1 descent (D, lem:section_localization_descent)
  % consumes per affine cover member.
  \textit{On every affine open of a cover member, the restricted sheaf is a geometric tilde.}
  Let \(M\) be a sheaf of modules on \(X\) with quasi-coherence data \(q\), let \(i\) be a cover
  index, and let \(W \subseteq Z := (q.X_i).\mathrm{toScheme}\) be an \emph{affine} open. Then the
  geometric restriction of \(M\) to the affine \(W \cong \Spec \Gamma(Z, W)\) (pulled back to \(Z\),
  then to \(W\), then across \(\mathrm{IsAffineOpen.isoSpec}\)) has an isomorphism
  \(\mathrm{fromTildeΓ}\) counit:
  \[
    \mathrm{IsIso}\bigl((\mathrm{isoSpec}_W.\mathrm{inv}^{*}\,(W.\iota^{*}\,(q.X_i).\iota^{*}\,M))
      .\mathrm{fromTildeΓ}\bigr).
  \]
  This is strictly more general than the basic-open keystone
  \cref{lem:isIso_fromTildeΓ_basicOpen_of_quasicoherent}, of which it is the affine instance
  \(W = (q.X_i).\iota^{-1}D(r)\); it is the form consumed per affine cover member by the gap1 descent
  \cref{lem:section_localization_descent}.
\end{lemma}
\begin{proof}
  \uses{def:presentation_pullback_iota_restrict, def:presentation_pullback_of_scheme_iso,
    lem:isIso_fromTildeΓ_of_presentation_mathlib}
  The slice presentation supplied by the quasi-coherence datum geometrizes via
  \cref{def:presentation_pullback_iota_restrict} to a global presentation of
  \(W.\iota^{*}\,(q.X_i).\iota^{*}\,M\) on \(W.\mathrm{toScheme}\). Since \(W\) is affine,
  \(\mathrm{IsAffineOpen.isoSpec}\) identifies \(W.\mathrm{toScheme} \cong \Spec \Gamma(Z, W)\);
  transporting the presentation across this iso by
  \cref{def:presentation_pullback_of_scheme_iso} lands a global presentation on the genuine affine
  \(\Spec \Gamma(Z, W)\). A global presentation forces the \(\widetilde{(-)}\)-counit to be an
  isomorphism (\cref{lem:isIso_fromTildeΓ_of_presentation_mathlib}).
\end{proof}

% --- (P1) Per-element local-tilde transport ---

\begin{lemma}\leanok
[Quasi-coherent restricts to a tilde on each basic open of the cover (gap1, P1)]
  \label{lem:isIso_fromTildeΓ_basicOpen_of_quasicoherent}
  \lean{AlgebraicGeometry.Scheme.Modules.isIso_fromTildeΓ_restrict_basicOpen}
  \uses{lem:isIso_fromTildeΓ_presentationPullback, lem:over_restrict_iso,
    lem:over_restrict_pullback_iso, lem:presentation_map_mathlib,
    lem:quasicoherentData_bind_mathlib, lem:modules_pullback_mathlib,
    lem:isIso_fromTildeΓ_of_presentation_mathlib, lem:isLocalization_basicOpen_mathlib,
    lem:exists_finite_basicOpen_cover_le_quasicoherentData}
  % NOTE: RESOLVED and axiom-clean. `isIso_fromTildeΓ_restrict_basicOpen` is formalized
  % (no sorry; {propext, Classical.choice, Quot.sound}) via the 5-step affine descent: global
  % presentation of N = (q.X i)^* M on the scheme Z, slice to N.over W, geometrize via
  % `overRestrictPresentation`, transport across `IsAffineOpen.isoSpec` (W affine = preimage of D(r)),
  % then `isIso_fromTildeΓ_of_presentation`. The general form is `isIso_fromTildeΓ_presentationPullback`
  % (per affine open of a cover member). It replaces the former lem:exists_isIso_fromTildeΓ_basicOpen_cover,
  % whose \lean{} pin named a non-existent decl.
  \textit{Per-element step: on each basic open of the refined cover, the restricted sheaf is a
  geometric tilde.}
  Let \(M\) be a sheaf of modules on \(\Spec R\) with quasi-coherence data \(q\), and let \(r \in R\)
  be such that \(D(r) \le q.X_i\) for some member of the cover (as produced by
  \cref{lem:exists_finite_basicOpen_cover_le_quasicoherentData}). Then the restricted counit is an
  isomorphism,
  \[
    \mathrm{IsIso}\bigl((M|_{D(r)}).\mathrm{fromTildeΓ}\bigr),
  \]
  i.e. \(M|_{D(r)}\) is identified with the tilde \(\widetilde{N}\) of the \(R_r\)-module
  \(N = \Gamma(M, D(r))\) on \(\Spec R_r \cong D(r)\).
\end{lemma}
\begin{proof}
  \uses{lem:isIso_fromTildeΓ_presentationPullback, lem:over_restrict_iso,
    lem:over_restrict_pullback_iso, lem:presentation_map_mathlib,
    lem:quasicoherentData_bind_mathlib, lem:modules_pullback_mathlib,
    lem:isIso_fromTildeΓ_of_presentation_mathlib, lem:isLocalization_basicOpen_mathlib,
    lem:exists_finite_basicOpen_cover_le_quasicoherentData,
    def:over_restrict_presentation, def:presentation_pullback_iota_of_quasicoherentData}
  The quasi-coherence datum supplies a presentation
  \(q.\mathrm{presentation}\,i : (M.\mathrm{over}\,(q.X_i)).\mathrm{Presentation}\) on the slice over
  \(q.X_i\). Restrict it along the inclusion of opens \(D(r) \subseteq q.X_i\) by transporting across
  the iterated Grothendieck slice — the template is \(\mathrm{QuasicoherentData.bind}\)
  (\cref{lem:quasicoherentData_bind_mathlib}) — using \(\mathrm{Presentation.map}\)
  (\cref{lem:presentation_map_mathlib}) along the colimit-preserving slice functor; this yields a
  presentation of \(M.\mathrm{over}\,(D(r))\). Pushing through the slice-to-geometric bridge
  \cref{lem:over_restrict_iso} converts it into a presentation of the geometric restriction
  \((\mathrm{restrictFunctor}\,(D(r)).\iota).\mathrm{obj}\,M\); then, across the affine
  identification \(D(r) \cong \Spec R_r\) (\cref{lem:isLocalization_basicOpen_mathlib}, by
  \(\mathrm{Presentation.map}\) along the pullback functor \cref{lem:modules_pullback_mathlib}), into
  a genuine global presentation of \(M|_{D(r)}\) as a sheaf of modules on \(\Spec R_r\). A global
  presentation forces the \(\widetilde{(-)}\)-counit to be an isomorphism
  (\cref{lem:isIso_fromTildeΓ_of_presentation_mathlib}), giving
  \(\mathrm{IsIso}\bigl((M|_{D(r)}).\mathrm{fromTildeΓ}\bigr)\).
\end{proof}

% --- (D) The section-localization descent: the gap1 keystone ---

\begin{lemma}\leanok
[Section-localization descent from a cover hypothesis (gap1 keystone, D, cover form)]
  \label{lem:section_localization_descent_of_cover}
  \lean{AlgebraicGeometry.Scheme.Modules.isLocalizedModule_basicOpen_descent_of_cover}
  \uses{lem:map_units_restrict_basicOpen, lem:existsUnique_gluing_mathlib,
    lem:eq_of_locally_eq_mathlib, lem:isLimitPullbackCone_mathlib,
    lem:isLocalization_flat_mathlib}
  % NOTE: this is the LANDED, axiom-clean keystone (Lean decl
  % isLocalizedModule_basicOpen_descent_of_cover EXISTS). It takes the per-cover-element localization
  % data Hfr as a HYPOTHESIS and assembles all three IsLocalizedModule fields by pure finite-cover
  % sheaf theory; it does NOT route through the global QCoh ≃ Mod equivalence (= gap1 itself). The
  % NAMED form lem:section_localization_descent is this lemma instantiated at the quasi-coherence
  % cover once Hfr is produced (lem:pullback_gamma_top_iso).
  % SOURCE: [Stacks Project], "Properties of Schemes", \S "Sections over principal opens",
  % Lemma \texttt{lemma-invert-f-sections} (read from references/stacks-properties.tex, L2152-L2171).
  % Cf. [Hartshorne], II.5.3.
  % SOURCE QUOTE:
  % "Let $X$ be a scheme. Let $f \in \Gamma(X, \mathcal{O}_X)$.
  %  Denote $X_f \subset X$ the open where $f$ is invertible, see
  %  Schemes, Lemma \ref{schemes-lemma-f-open}.
  %  If $X$ is quasi-compact and quasi-separated, the canonical map
  %  $$
  %  \Gamma(X, \mathcal{O}_X)_f \longrightarrow \Gamma(X_f, \mathcal{O}_X)
  %  $$
  %  is an isomorphism. Moreover, if $\mathcal{F}$ is a quasi-coherent
  %  sheaf of $\mathcal{O}_X$-modules the map
  %  $$
  %  \Gamma(X, \mathcal{F})_f \longrightarrow \Gamma(X_f, \mathcal{F})
  %  $$
  %  is an isomorphism."
  \textit{Source: [Stacks Project], "Properties of Schemes", \S ``Sections over principal opens'',
  Lemma \texttt{lemma-invert-f-sections}; cf.\ [Hartshorne], II.5.3.}
  Let \(M\) be a sheaf of modules on \(\Spec R\) and \(f \in R\). Suppose given a finite basic-open
  cover of \(\Spec R\): a finite set \(t \subseteq R\) with \(\mathrm{span}(t) = (1)\), so that
  \(\{D(r)\}_{r \in t}\) covers \(\Spec R\). Suppose moreover the \emph{per-piece localization data}
  \[
    \mathrm{Hfr} : \quad \text{for every open } U \le D(r) \text{ for some } r \in t, \quad
    \bigl(\Gamma(M, U) \longrightarrow \Gamma(M,\, D(f) \sqcap U)\bigr)
    \text{ is } \mathrm{IsLocalizedModule}(\mathrm{powers}\,f),
  \]
  i.e.\ on every open contained in a cover member -- in particular each \(D(r)\) and each overlap
  \(D(r) \sqcap D(r')\) -- the restriction to the \(f\)-locus exhibits the target as the localization
  of the source at the powers of \(f\). Then the global section restriction
  \(\Gamma(M, \top) \to \Gamma(M, D(f))\) is itself \(\mathrm{IsLocalizedModule}(\mathrm{powers}\,f)\)
  over \(R\); that is, \(\Gamma(M, D(f)) \cong (\mathrm{powers}\,f)^{-1}\,\Gamma(M, \top)\).
\end{lemma}
% NOTE: the lemma STATEMENT is cited to Stacks lemma-invert-f-sections / Hartshorne II.5.3; the proof
% below is the project-internal finite-cover descent (gluing + separatedness + flatness of
% localization) that deliberately does NOT route through the global affine equivalence (which is gap1
% itself), so there is no transcribed `% SOURCE QUOTE PROOF`.
\begin{proof}
  \uses{lem:map_units_restrict_basicOpen, lem:existsUnique_gluing_mathlib,
    lem:eq_of_locally_eq_mathlib, lem:isLimitPullbackCone_mathlib,
    lem:isLocalization_flat_mathlib}
  We verify the three fields of \(\mathrm{IsLocalizedModule}(\mathrm{powers}\,f)\) for the restriction
  \(\theta : \Gamma(M, \top) \to \Gamma(M, D(f))\).

  \emph{\(\mathrm{map\_units}\).} For every \(n\) the element \(f^n\) acts invertibly on
  \(\Gamma(M, D(f))\): on the basic open \(D(f)\) the global section \(f\) is a unit of
  \(\Gamma(\mathcal{O}, D(f)) = R_f\), and scalar multiplication by a unit is bijective. This holds
  for an arbitrary sheaf of modules \(M\) and is \cref{lem:map_units_restrict_basicOpen}; it does not
  use \(\mathrm{Hfr}\).

  \emph{\(\mathrm{surj}\) (every section over \(D(f)\) is a fraction).} Let
  \(y \in \Gamma(M, D(f))\). On each cover member \(D(r)\), the restriction
  \(y|_{D(f) \sqcap D(r)}\) lies in the localization \(\Gamma(M, D(f) \sqcap D(r))\) of
  \(\Gamma(M, D(r))\) at \(\mathrm{powers}\,f\) (the case \(U = D(r)\) of \(\mathrm{Hfr}\)), so there
  are \(x_r \in \Gamma(M, D(r))\) and \(n_r\) with \(f^{n_r}\, y|_{D(f)\sqcap D(r)} = x_r|_{D(f)\sqcap
  D(r)}\). Taking a common exponent \(N \ge \max_r n_r\) (the cover is finite) and rescaling, we may
  assume \(f^{N} y\) is represented on each \(D(r)\) by a section \(\tilde x_r\). On an overlap
  \(D(r) \sqcap D(r')\) the two representatives \(\tilde x_r, \tilde x_{r'}\) restrict to sections of
  \(\Gamma(M, D(r) \sqcap D(r'))\) that agree after further multiplication by \(f\): both compute
  \(f^{N} y\) on \(D(f) \sqcap D(r) \sqcap D(r')\), and the overlap restriction
  \(\Gamma(M, D(r)\sqcap D(r')) \to \Gamma(M, D(f) \sqcap D(r) \sqcap D(r'))\) is
  \(\mathrm{IsLocalizedModule}(\mathrm{powers}\,f)\) (the case \(U = D(r) \sqcap D(r')\) of
  \(\mathrm{Hfr}\)), hence injective up to a power of \(f\); so there is a common further exponent
  \(P\) with \(f^{P}\tilde x_r\) and \(f^{P}\tilde x_{r'}\) equal on the overlap. The family
  \(\{f^{P}\tilde x_r\}_r\) is then compatible on overlaps and glues, by the sheaf condition over the
  finite cover (\cref{lem:existsUnique_gluing_mathlib}), to a global section
  \(x \in \Gamma(M, \top)\). Its restriction to \(D(f)\) equals \(f^{N+P} y\): the two agree on each
  \(D(f) \sqcap D(r)\), and \(\{D(f)\sqcap D(r)\}_r\) covers \(D(f)\), so they agree by separatedness
  (\cref{lem:eq_of_locally_eq_mathlib}). Thus \(\theta(x) = f^{N+P} y\), exhibiting \(y\) as the
  fraction \(x / f^{N+P}\).

  \emph{\(\mathrm{exists\_of\_eq}\) (a vanishing fraction is annihilated by a power of \(f\)).}
  Let \(x \in \Gamma(M, \top)\) with \(\theta(x) = x|_{D(f)} = 0\). On each cover member \(D(r)\), the
  restriction \(x|_{D(r)}\) maps to \(0\) in \(\Gamma(M, D(f) \sqcap D(r))\), so by the localization
  hypothesis (the case \(U = D(r)\) of \(\mathrm{Hfr}\), which gives the \(\mathrm{exists\_of\_eq}\)
  field on that piece) there is \(n_r\) with \(f^{n_r}\, x|_{D(r)} = 0\). A common exponent
  \(N \ge \max_r n_r\) over the finite cover gives \(f^{N} x|_{D(r)} = 0\) for every \(r\); since the
  \(D(r)\) cover \(\Spec R\), separatedness (\cref{lem:eq_of_locally_eq_mathlib}) forces
  \(f^{N} x = 0\) in \(\Gamma(M, \top)\).

  These three fields are exactly \(\mathrm{IsLocalizedModule}(\mathrm{powers}\,f)\) for \(\theta\).
  Equivalently the comparison \((\mathrm{powers}\,f)^{-1}\Gamma(M, \top) \to \Gamma(M, D(f))\) is an
  isomorphism -- the \(\Gamma(X, \mathcal{F})_f \cong \Gamma(X_f, \mathcal{F})\) of Stacks
  \texttt{lemma-invert-f-sections} on \(X = \Spec R\), \(X_f = D(f)\), obtained from local data over a
  finite cover by sheaf gluing and separatedness, with the finite limit interacting with localization
  through the flatness of localization (\cref{lem:isLocalization_flat_mathlib},
  \cref{lem:isLimitPullbackCone_mathlib}), never through the global affine equivalence.
\end{proof}

% --- (D, engines) The private sheaf-gluing sub-steps of the cover-form descent (coverage) ---

\begin{lemma}

\leanok
[Restriction maps of a sheaf compose (descent bookkeeping)]
  \label{lem:res_comp}
  \lean{AlgebraicGeometry.Scheme.Modules.res_comp}
  \textit{Project-local bookkeeping for the section-localization descent.}
  For a sheaf \(F\) on \(\Spec R\) and opens \(C \le B \le A\), the composite of the restriction maps
  \(\Gamma(F, A) \to \Gamma(F, B) \to \Gamma(F, C)\) equals the direct restriction
  \(\Gamma(F, A) \to \Gamma(F, C)\).
\end{lemma}
\begin{proof}
  The two poset inclusions \(C \le B \le A\) compose to \(C \le A\) in the poset of opens (uniqueness
  of poset-hom), and \(F.\mathrm{map}\) is functorial, so the two presheaf maps compose to the direct
  one.
\end{proof}

\begin{lemma}

\leanok
[A finite spanning family covers $\Spec R$ by basic opens (descent bookkeeping)]
  \label{lem:iSup_basicOpen_subtype_eq_top}
  \lean{AlgebraicGeometry.Scheme.Modules.iSup_basicOpen_subtype_eq_top}
  \uses{lem:isLocalization_basicOpen_mathlib}
  \textit{Project-local glue for the sheaf-gluing reduction of the descent.}
  If a finite family \(t \subseteq R\) spans the unit ideal, \(\mathrm{span}(t) = (1)\), then the
  supremum of the basic opens \(D(r)\) over \(r \in t\) is \(\top\); i.e.\ \(\{D(r)\}_{r \in t}\) is a
  finite basic-open cover of \(\Spec R\).
\end{lemma}
\begin{proof}
  Standard: \(\bigsqcup_{r \in t} D(r) = \top\) iff the \(r\) generate the unit ideal
  (\(\mathrm{PrimeSpectrum.iSup\_basicOpen\_eq\_top\_iff}\)), applied to \(\mathrm{span}(t) = (1)\).
\end{proof}

\begin{lemma}

\leanok
[Overlap representatives agree after a power of $f$ (descent surjectivity sub-step)]
  \label{lem:descent_overlap_agree}
  \lean{AlgebraicGeometry.Scheme.Modules.descent_overlap_agree}
  \uses{lem:res_comp}
  \textit{Project-local sub-step of the cover-form descent.}
  In the surjectivity argument of \cref{lem:section_localization_descent_of_cover}: given a section
  \(y\) over \(D(f)\) and, on a cover member \(D(r)\), a representative \(\tilde x_r\) with
  \(f^N y = \tilde x_r|_{D(f)\sqcap D(r)}\), the restrictions of two such representatives
  \(\tilde x_r, \tilde x_{r'}\) to the overlap \(D(r) \sqcap D(r')\) become equal after multiplication
  by a common power of \(f\).
\end{lemma}
\begin{proof}
  Both \(\tilde x_r\) and \(\tilde x_{r'}\) compute \(f^N y\) on \(D(f) \sqcap D(r) \sqcap D(r')\); the
  overlap restriction \(\Gamma(M, D(r)\sqcap D(r')) \to \Gamma(M, D(f) \sqcap D(r) \sqcap D(r'))\) is a
  localization at \(\mathrm{powers}\,f\) (the overlap case of \(\mathrm{Hfr}\)), hence injective up to a
  power of \(f\); compatibility of the restriction maps is \cref{lem:res_comp}.
\end{proof}

\begin{lemma}

\leanok
[Surjectivity field of the cover-form descent]
  \label{lem:descent_surj}
  \lean{AlgebraicGeometry.Scheme.Modules.descent_surj}
  \uses{lem:descent_overlap_agree, lem:res_comp, lem:iSup_basicOpen_subtype_eq_top,
    lem:existsUnique_gluing_mathlib, lem:eq_of_locally_eq_mathlib}
  \textit{Project-local: the \(\mathrm{surj}\) engine of the cover-form descent.}
  Given the per-cover-element localization data \(\mathrm{Hfr}\) on a finite basic-open cover, every
  section \(y \in \Gamma(M, D(f))\) is, up to a power of \(f\), the restriction of a global section: there
  are \(x \in \Gamma(M, \top)\) and \(n\) with \(x|_{D(f)} = f^n\, y\).
\end{lemma}
\begin{proof}
  Per cover member, \(\mathrm{Hfr}\) at \(U = D(r)\) gives a representative of \(f^{n_r} y\); a common
  exponent \(N\) (finite cover, \cref{lem:iSup_basicOpen_subtype_eq_top}) yields representatives
  \(\tilde x_r\), which agree on overlaps after a further power of \(f\)
  (\cref{lem:descent_overlap_agree}). The patched family glues
  (\cref{lem:existsUnique_gluing_mathlib}) to \(x \in \Gamma(M, \top)\), and \(x|_{D(f)} = f^{N+P} y\) by
  separatedness over the cover \(\{D(f)\sqcap D(r)\}\) (\cref{lem:eq_of_locally_eq_mathlib}), the
  restrictions composing by \cref{lem:res_comp}.
\end{proof}

\begin{lemma}

\leanok
[Separatedness/torsion field of the cover-form descent]
  \label{lem:descent_smul_eq_zero}
  \lean{AlgebraicGeometry.Scheme.Modules.descent_smul_eq_zero}
  \uses{lem:iSup_basicOpen_subtype_eq_top, lem:eq_of_locally_eq_mathlib, lem:res_comp}
  \textit{Project-local: the \(\mathrm{exists\_of\_eq}\) engine of the cover-form descent.}
  Given the per-cover-element localization data \(\mathrm{Hfr}\), any global section
  \(x \in \Gamma(M, \top)\) with \(x|_{D(f)} = 0\) is annihilated by a power of \(f\): there is \(n\) with
  \(f^n\, x = 0\).
\end{lemma}
\begin{proof}
  Per cover member, \(\mathrm{Hfr}\) at \(U = D(r)\) (its \(\mathrm{exists\_of\_eq}\) field) kills
  \(x|_{D(r)}\) by some \(f^{n_r}\); the finite sup of these powers
  (\cref{lem:iSup_basicOpen_subtype_eq_top}) kills every \(x|_{D(r)}\), and separatedness over the cover
  (\cref{lem:eq_of_locally_eq_mathlib}, restrictions composing by \cref{lem:res_comp}) lifts this to
  \(f^n\, x = 0\) globally.
\end{proof}

% --- (D, basic-open cover form) The form gap1 actually consumes ---

\begin{lemma}

\leanok
[Section-localization descent from a BASIC-open cover hypothesis (gap1 keystone, D, basic-open form)]
  \label{lem:section_localization_descent_of_basicOpen_cover}
  \lean{AlgebraicGeometry.Scheme.Modules.isLocalizedModule_basicOpen_descent_of_basicOpen_cover}
  \uses{lem:section_localization_descent_of_cover, lem:descent_surj, lem:descent_smul_eq_zero,
    lem:map_units_restrict_basicOpen, lem:existsUnique_gluing_mathlib,
    lem:eq_of_locally_eq_mathlib}
  % NOTE: this is the form gap1 USES. Same conclusion as the general-U cover form
  % lem:section_localization_descent_of_cover, but Hfr is demanded only at BASIC opens D(s) ≤ D(r).
  % The general-U form is axiom-clean, but its Hfr (a localization on EVERY open U ≤ D(r)) cannot be
  % discharged for an arbitrary quasi-coherent M — a general open of Spec R_r need not be quasi-compact.
  % The sheaf-gluing engines descent_surj/descent_smul_eq_zero only ever consult Hfr at basic opens
  % D(s) and overlaps D(r·r') = D(r)⊓D(r') (PrimeSpectrum.basicOpen_mul), all quasi-compact, so the
  % basic-open hypothesis suffices to rebuild all three IsLocalizedModule fields.
  \textit{Source: [Stacks Project], "Properties of Schemes", \S ``Sections over principal opens'',
  Lemma \texttt{lemma-invert-f-sections}; cf.\ [Hartshorne], II.5.3.}
  Let \(M\) be a sheaf of modules on \(\Spec R\) and \(f \in R\). Suppose given a finite basic-open
  cover of \(\Spec R\): a finite set \(t \subseteq R\) with \(\mathrm{span}(t) = (1)\). Suppose moreover
  the \emph{basic-open} per-cover-element localization data
  \[
    \mathrm{Hfr} : \quad \text{for every } s \in R \text{ with } D(s) \le D(r) \text{ for some } r \in t,
    \quad \bigl(\Gamma(M, D(s)) \longrightarrow \Gamma(M,\, D(f) \sqcap D(s))\bigr)
    \text{ is } \mathrm{IsLocalizedModule}(\mathrm{powers}\,f).
  \]
  Then the global section restriction \(\Gamma(M, \top) \to \Gamma(M, D(f))\) is itself
  \(\mathrm{IsLocalizedModule}(\mathrm{powers}\,f)\) over \(R\).
\end{lemma}
\begin{proof}
  \uses{lem:section_localization_descent_of_cover, lem:descent_surj, lem:descent_smul_eq_zero,
    lem:map_units_restrict_basicOpen}
  A thin wrapper over the cover-form descent. The three \(\mathrm{IsLocalizedModule}\) fields are rebuilt
  exactly as in \cref{lem:section_localization_descent_of_cover}, but every appeal to the per-piece data
  is at a basic open: \(\mathrm{map\_units}\) is \cref{lem:map_units_restrict_basicOpen} (independent of
  \(\mathrm{Hfr}\)); \(\mathrm{surj}\) is \cref{lem:descent_surj}, whose only \(\mathrm{Hfr}\) call is at
  the cover member \(D(r) = D(r)\) itself; and \(\mathrm{exists\_of\_eq}\) is
  \cref{lem:descent_smul_eq_zero}, again calling \(\mathrm{Hfr}\) only at each \(D(r)\). The overlaps it
  consults are \(D(r) \sqcap D(r') = D(r\cdot r')\) (a basic open, by
  \(\mathrm{PrimeSpectrum.basicOpen\_mul}\)), so the basic-open hypothesis \(\mathrm{Hfr}\) supplies every
  datum the engines need. (The general-\(U\) form
  \cref{lem:section_localization_descent_of_cover} would instead demand a localization on every open
  \(U \le D(r)\), which is not available for an arbitrary quasi-coherent \(M\); this basic-open form is
  the one gap1 uses.)
\end{proof}

% --- (D, section transport) The Hfr ingredient: pullback-along-open-immersion section comparison ---

\begin{lemma}\leanok
[Global sections of a pullback along an open immersion are sections over the image
  (gap1, Hfr section transport)]
  \label{lem:pullback_gamma_top_iso}
  \lean{AlgebraicGeometry.Scheme.Modules.gammaPullbackTopIso}
  \uses{lem:modules_pullback_mathlib, lem:modules_restrictFunctor_mathlib,
    lem:modules_restrictFunctorIsoPullback_mathlib, lem:over_restrict_pullback_iso,
    lem:isIso_fromTildeΓ_basicOpen_of_quasicoherent,
    lem:isLocalizedModule_restrict_of_isIso_fromTildeΓ}
  % NOTE: The SECTION analogue of P1's OBJECT-level transport
  % (\cref{lem:isIso_fromTildeΓ_basicOpen_of_quasicoherent}). TWO Mathlib-absent ingredients bridge
  % gammaPullbackTopIso and Hfr: (I) a ring-iso-semilinear IsLocalizedModule transport (Mathlib only
  % has same-ring of_linearEquiv), and (II) a base-change-of-localization R -> R_r bridge.
  % Project-bespoke infrastructure: no external verbatim source for the construction; it stands on
  % its sketch.
  \textit{Project-bespoke infrastructure: the section-level transport feeding the per-cover-element
  hypothesis \(\mathrm{Hfr}\) of \cref{lem:section_localization_descent_of_cover}.}
  Let \(j : U \hookrightarrow X\) be an open immersion of schemes and let \(M\) be a sheaf of modules
  on \(X\). The global sections of the pullback module along \(j\) are canonically isomorphic to the
  sections of \(M\) over the image of \(j\):
  \[
    \Gamma\bigl((\mathrm{pullback}\,j).\mathrm{obj}\,M,\ \top\bigr)
    \;\cong\; \Gamma\bigl(M,\ \mathrm{range}\,j\bigr),
  \]
  naturally in the open argument: for opens \(V \le \top\) of \(U\) the comparison
  \(\Gamma((\mathrm{pullback}\,j).\mathrm{obj}\,M, V) \cong \Gamma(M, j(V))\) intertwines the
  restriction maps of the two presheaves, so any two such instances commute with the restriction
  between them.
\end{lemma}

\begin{proof}
  \uses{lem:modules_pullback_mathlib, lem:modules_restrictFunctor_mathlib,
    lem:modules_restrictFunctorIsoPullback_mathlib}
  For an open immersion \(j\), Mathlib's restriction functor \(j^{\mathrm{res}}\) satisfies
  \(\Gamma(j^{\mathrm{res}} M, V) = \Gamma(M, j(V))\) \emph{definitionally}
  (\cref{lem:modules_restrictFunctor_mathlib}); in particular at \(V = \top\) the right-hand side is
  \(\Gamma(M, j(\top)) = \Gamma(M, \mathrm{range}\,j)\). The restriction functor is naturally
  isomorphic to the universal-property pullback,
  \(j^{\mathrm{res}} \cong j^{*} = \mathrm{pullback}\,j\)
  (\cref{lem:modules_restrictFunctorIsoPullback_mathlib},
  \cref{lem:modules_pullback_mathlib}). Applying the global-sections functor \(\Gamma(-, V)\) to the
  component of this natural isomorphism at \(M\) gives, for each open \(V\), an isomorphism
  \(\Gamma((\mathrm{pullback}\,j).\mathrm{obj}\,M, V) \cong \Gamma(j^{\mathrm{res}} M, V)
  = \Gamma(M, j(V))\); naturality of the functor isomorphism in \(V\) makes these commute with the
  presheaf restriction maps. Taking \(V = \top\) yields the displayed isomorphism with
  \(\Gamma(M, \mathrm{range}\,j)\).
\end{proof}

\begin{lemma}

\leanok
[Section transport along an open immersion, general in the open (coverage)]
  \label{lem:gamma_pullback_image_iso}
  \lean{AlgebraicGeometry.Scheme.Modules.gammaPullbackImageIso}
  \uses{lem:modules_pullback_mathlib, lem:modules_restrictFunctor_mathlib,
    lem:modules_restrictFunctorIsoPullback_mathlib}
  \textit{Project-bespoke; the general-in-\(U\) form of which \cref{lem:pullback_gamma_top_iso} is the
  \(U=\top\) instance.}
  For an open immersion \(f : X \hookrightarrow Y\), a sheaf of modules \(M\) on \(Y\), and any open
  \(U\) of \(X\), there is an isomorphism
  \(\Gamma((\mathrm{pullback}\,f).\mathrm{obj}\,M,\,U) \cong \Gamma(M,\,f\,''^{U}\,U)\) of abelian
  groups, obtained by applying the sections functor \(\Gamma(-,U)\) to the component at \(M\) of the
  inverse of Mathlib's \(\mathrm{restrictFunctorIsoPullback}\,f\). The codomain identity
  \(\Gamma((\mathrm{restrictFunctor}\,f).\mathrm{obj}\,M,U)=\Gamma(M,f\,''^{U}U)\) holds definitionally
  (\cref{lem:modules_restrictFunctor_mathlib}).
\end{lemma}
\begin{proof}
  \uses{lem:modules_restrictFunctorIsoPullback_mathlib, lem:modules_restrictFunctor_mathlib}
  Apply \(\mathrm{Functor.mapIso}\) along the sections functor
  \(\mathrm{toPresheaf}\,X \,\circ\, (\mathrm{evaluation}\,\cdot\,\cdot).\mathrm{obj}\,(\mathrm{op}\,U)\)
  to \(((\mathrm{restrictFunctorIsoPullback}\,f).\mathrm{symm}).\mathrm{app}\,M\); the codomain reads as
  \(\Gamma(M, f\,''^{U}U)\) by \cref{lem:modules_restrictFunctor_mathlib}. (Going through
  \(\mathrm{Functor.mapIso}\) sidesteps the \(\mathrm{IsIso}\,(\varphi.\mathrm{app}\,U)\) instance
  synthesis, which does not fire here.)
\end{proof}

\begin{lemma}

\leanok
[The section transport intertwines restriction maps (coverage)]
  \label{lem:gamma_pullback_image_iso_hom_naturality}
  \lean{AlgebraicGeometry.Scheme.Modules.gammaPullbackImageIso_hom_naturality}
  \uses{lem:gamma_pullback_image_iso, lem:modules_restrictFunctor_mathlib}
  \textit{Project-bespoke; the ``naturally in the open argument'' clause of
  \cref{lem:pullback_gamma_top_iso}.}
  For opens \(V \le U\) of \(X\), the section transports of \cref{lem:gamma_pullback_image_iso} at
  \(U\) and \(V\) commute with the presheaf restriction maps along \(V \le U\) on either side --- in
  particular along \(D(f)\sqcap U \le U\), the restriction the cover-form descent consumes.
\end{lemma}
\begin{proof}
  \uses{lem:gamma_pullback_image_iso}
  The forward map of \cref{lem:gamma_pullback_image_iso} reduces definitionally to
  \(((\mathrm{restrictFunctorIsoPullback}\,f).\mathrm{symm}.\mathrm{app}\,M).\mathrm{hom}.\mathrm{mapPresheaf}\),
  so the claim is the naturality square of that morphism of abelian presheaves, applied to the inclusion
  \(V \le U\).
\end{proof}

% NOTE: How lem:pullback_gamma_top_iso produces the gated Hfr (and thereby the named-form descent and
% gap1). For M : (Spec R).Modules quasi-coherent with datum q and a cover member q.X i ⊇ D(r), the P1
% transport (\cref{lem:isIso_fromTildeΓ_basicOpen_of_quasicoherent}) gives IsIso fromTildeΓ for the
% module M'_transported obtained by pulling M back along ι(q.X i), then along the basic-open ι, then
% across IsAffineOpen.isoSpec.inv. Chaining the section comparison \cref{lem:pullback_gamma_top_iso}
% through those same three pullbacks (and the pullback-form bridge
% \cref{lem:over_restrict_pullback_iso}) identifies Γ(M'_transported, ⊤) ≅ Γ(M, D(r)) and, at the
% f-locus open, Γ(M'_transported, D(f')) ≅ Γ(M, D(f) ⊓ D(r)), intertwining the restriction maps.
% Together with IsIso fromTildeΓ and \cref{lem:isLocalizedModule_restrict_of_isIso_fromTildeΓ} this
% yields IsLocalizedModule(powers f) on Γ(M, D(r)) → Γ(M, D(f) ⊓ D(r)) — exactly Hfr at U = D(r) (and
% at the overlaps U = D(r) ⊓ D(r')). Feeding Hfr to \cref{lem:section_localization_descent_of_cover}
% at the quasi-coherence cover \cref{lem:exists_finite_basicOpen_cover_le_quasicoherentData} closes
% the named form \cref{lem:section_localization_descent}, and gap1
% \cref{lem:qcoh_affine_isIso_fromTildeΓ} follows via
% \cref{lem:isIso_fromTildeΓ_of_isLocalizedModule_restrict}.

\begin{lemma}\leanok
[Ring-iso-semilinear localized-module transport (gap1 ingredient I)]
  \label{lem:isLocalizedModule_ringEquiv_semilinear}
  \lean{AlgebraicGeometry.Scheme.Modules.isLocalizedModule_of_ringEquiv_semilinear}
  Let \(\sigma : R \xrightarrow{\sim} R'\) be a ring isomorphism, \(S \le R\) a submonoid, and
  \(g : M_1 \to M_2\) an \(R\)-linear map satisfying \(\mathrm{IsLocalizedModule}\ S\ g\). Suppose
  \(e_1 : M_1 \simeq_+ N_1\) and \(e_2 : M_2 \simeq_+ N_2\) are additive isomorphisms that are
  \(\sigma\)-semilinear (\(e_i(a \cdot x) = \sigma(a) \cdot e_i(x)\)), and \(h : N_1 \to N_2\) is an
  \(R'\)-linear map intertwining them, \(h(e_1 x) = e_2(g x)\). Then \(h\) satisfies
  \(\mathrm{IsLocalizedModule}\ (S.\mathrm{map}\ \sigma)\ h\) over \(R'\).
\end{lemma}
\begin{proof}
  Transport the three \(\mathrm{IsLocalizedModule}\) fields of \(g\) along the semilinear pair
  \((e_1,e_2)\). \emph{map\_units:} the action of \(s' = \sigma(s)\) on \(N_2\) is conjugate via \(e_2\) to
  the (bijective, since a unit) action of \(s\) on \(M_2\), hence bijective, hence a unit
  (\(\mathrm{Module.End.isUnit\_iff}\)). \emph{surj:} given \(y \in N_2\), pull back \(e_2^{-1}y\), apply
  the surjectivity field of \(g\) to get \(g\,m = s\cdot e_2^{-1}y\), and push forward by \(e_2\); the
  image-submonoid witness is \(\langle \sigma s, s, s\in S, \mathrm{rfl}\rangle\). \emph{exists\_of\_eq:}
  symmetric, pulling the equality back through \(e_1\). Mathlib provides only the same-ring
  \(\mathrm{IsLocalizedModule.of\_linearEquiv}\); crossing a ring iso is the genuine gap.
\end{proof}

\begin{lemma}\leanok
[Base-change-of-localization descent (gap1 ingredient II)]
  \label{lem:isLocalizedModule_restrictScalars_powers_algebraMap}
  \lean{AlgebraicGeometry.Scheme.Modules.isLocalizedModule_restrictScalars_powers_algebraMap}
  Let \(R_r\) be an \(R\)-algebra, \(f \in R\), and \(g : M_1 \to M_2\) an \(R_r\)-linear map with
  \(\mathrm{IsScalarTower}\ R\ R_r\ M_i\). If \(g\) satisfies
  \(\mathrm{IsLocalizedModule}\ (\mathrm{powers}(\mathrm{algebraMap}\ R\ R_r\ f))\ g\) over \(R_r\), then
  the restriction of scalars of \(g\) to an \(R\)-linear map satisfies
  \(\mathrm{IsLocalizedModule}\ (\mathrm{powers}\ f)\) over \(R\).
\end{lemma}
\begin{proof}
  Each \(\mathrm{IsLocalizedModule}\) field transfers along the identification
  \(f^n \leftrightarrow (\mathrm{algebraMap}\ R\ R_r\ f)^n\) (\(\mathrm{map\_pow}\)) and the scalar
  compatibility \((\mathrm{algebraMap}\ R\ R_r\ a)\cdot m = a \cdot m\) (\(\mathrm{algebraMap\_smul}\) via
  the scalar tower); the two \(\mathrm{End}\) maps are the same underlying function, so the unit/surjectivity
  data transfer verbatim. This lets the \(R_r\)-level localization produced on the slice
  \(\Spec R_r\) be read back as the \(R\)-level data the cover-form descent consumes.
\end{proof}

\begin{lemma}\leanok
[Combined ring-iso + base-change localized-module transport (gap1 combiner)]
  \label{lem:isLocalizedModule_powers_transport}
  \lean{AlgebraicGeometry.Scheme.Modules.isLocalizedModule_powers_transport}
  \uses{lem:isLocalizedModule_ringEquiv_semilinear,
    lem:isLocalizedModule_restrictScalars_powers_algebraMap}
  % Role: the algebra/category-free CORE that the geometric Hfr producer
  % lem:section_localization_hfr_basicOpen feeds — the final combiner of bridges (I) and (II). The
  % rewrite (Submonoid.powers f').map σ = powers (algebraMap R A f) is Mathlib's Submonoid.map_powers.
  \textit{Project-bespoke: the combiner of bridges (I)~\cref{lem:isLocalizedModule_ringEquiv_semilinear}
  and (II)~\cref{lem:isLocalizedModule_restrictScalars_powers_algebraMap}.}
  Let \(\sigma : S \xrightarrow{\sim} A\) be a ring isomorphism with \(A\) an \(R\)-algebra, and let
  \(f \in R\), \(f' \in S\) satisfy \(\sigma(f') = \mathrm{algebraMap}\ R\ A\ f\). Let
  \(g : M_1 \to M_2\) be an \(S\)-linear map with \(\mathrm{IsLocalizedModule}(\mathrm{powers}\,f')\ g\),
  and let \(e_1 : M_1 \simeq_+ N_1\), \(e_2 : M_2 \simeq_+ N_2\) be \(\sigma\)-semilinear additive
  isomorphisms onto \(A\)-modules \(N_i\) that are also \(R\)-modules with
  \(\mathrm{IsScalarTower}\ R\ A\ N_i\). If \(h : N_1 \to N_2\) is \(A\)-linear and intertwines them,
  \(h(e_1 x) = e_2(g x)\), then the restriction of scalars of \(h\) to an \(R\)-linear map satisfies
  \(\mathrm{IsLocalizedModule}(\mathrm{powers}\,f)\) over \(R\).
\end{lemma}
\begin{proof}
  \uses{lem:isLocalizedModule_ringEquiv_semilinear,
    lem:isLocalizedModule_restrictScalars_powers_algebraMap}
  Bridge~(I) \cref{lem:isLocalizedModule_ringEquiv_semilinear} applied to the semilinear pair
  \((e_1, e_2)\) and the intertwiner \(h\) gives
  \(\mathrm{IsLocalizedModule}\bigl((\mathrm{powers}\,f').\mathrm{map}\,\sigma\bigr)\ h\) over \(A\). The
  submonoid rewrites: \((\mathrm{powers}\,f').\mathrm{map}\,\sigma = \mathrm{powers}(\sigma f') =
  \mathrm{powers}(\mathrm{algebraMap}\ R\ A\ f)\) (Mathlib's \(\mathrm{Submonoid.map\_powers}\) together
  with \(\sigma(f') = \mathrm{algebraMap}\ R\ A\ f\)). Then bridge~(II)
  \cref{lem:isLocalizedModule_restrictScalars_powers_algebraMap} descends the localization to
  \(\mathrm{powers}\,f\) over \(R\) after restriction of scalars.
\end{proof}

% NOTE: the remaining wall between the two bridges above and Hfr is the SEMILINEARITY of the
% section transport gammaPullbackImageIso over the open-immersion structure-sheaf ring iso. The next two
% blocks state that sub-build.

\begin{definition}\leanok
[Open-immersion structure-sheaf ring iso on an image open]
  \label{def:gamma_image_ring_equiv}
  \lean{AlgebraicGeometry.Scheme.Modules.gammaImageRingEquiv}
  \uses{lem:gamma_pullback_image_iso}
  % LEAN TYPE. The image open is the project operator `j ''ᵁ V` (image of opens under a scheme morphism,
  % as already used in the signature of `gammaPullbackImageIso`: `Γ(M, f ''ᵁ U)`, QuotScheme.lean:1783).
  % Intended signature, for `j : U ⟶ Y` `[IsOpenImmersion j]` and `V : Y.Opens` with `V ≤ ...`:
  %   `gammaImageRingEquiv (j) (V) : Γ(U, V) ≃+* Γ(Y, j ''ᵁ V)`
  % The Lean built the iso in direction SOURCE → IMAGE, `Γ(U, V) ≃+* Γ(Y, j ''ᵁ V)`; the displayed
  % `\sigma_V` below now matches this orientation.
  % where `Γ(Y, W) := Y.presheaf.obj (op W)` is the structure-sheaf ring of sections. The iso is the
  % section transport along the canonical iso of the open immersion onto its image (`IsOpenImmersion`
  % gives `U ≅ (j ''ᵁ ⊤)` as schemes; restrict to `V`).
  For an open immersion \(j : U \to Y\) of schemes and an open \(V \le U\), the immersion is an isomorphism
  onto its image \(j\,''^{U}\,V\); the structure-sheaf comparison of the open immersion gives a ring
  isomorphism in the source~\(\to\)~image direction
  \[
    \sigma_{V} : \Gamma(\mathcal{O}_U,\ V) \xrightarrow{\ \sim\ } \Gamma(\mathcal{O}_Y,\ j\,''^{U}\,V).
  \]
  (This source~\(\to\)~image orientation is the load-bearing one: it makes
  \cref{lem:gamma_pullback_image_iso_hom_semilinear} typecheck with \(a : \Gamma(\mathcal{O}_U, V)\) and
  feeds bridge~(I) \cref{lem:isLocalizedModule_ringEquiv_semilinear} with \(\sigma : R_{\mathrm{dom}} \simeq
  R_{\mathrm{cod}}\) verbatim.)
\end{definition}

\begin{lemma}\leanok
[Semilinearity of the pullback section transport (gap1 semilinearity wall)]
  \label{lem:gamma_pullback_image_iso_hom_semilinear}
  \lean{AlgebraicGeometry.Scheme.Modules.gammaPullbackImageIso_hom_semilinear}
  \uses{lem:gamma_pullback_image_iso, def:gamma_image_ring_equiv}
  Let \(j : U \to Y\) be an open immersion, \(M\) a sheaf of modules on \(Y\), and \(V \le U\) an open. The
  forward map of the section transport \cref{lem:gamma_pullback_image_iso},
  \[
    (\mathrm{gammaPullbackImageIso}\ j\ M\ V).\mathrm{hom} :
      \Gamma\bigl((\mathrm{pullback}\ j).\mathrm{obj}\ M,\ V\bigr) \longrightarrow \Gamma(M,\ j\,''^{U}\,V),
  \]
  is \(\sigma_V\)-semilinear (\cref{def:gamma_image_ring_equiv}): for \(a\) a section of the structure sheaf
  and \(x\) a section of the pullback module, the action on the pullback side and the action on the \(M\)
  side agree under \(\sigma_V\),
  \[
    \mathrm{hom}(a \cdot x) \;=\; \sigma_V(a) \cdot \mathrm{hom}(x).
  \]
\end{lemma}
\begin{proof}
  \uses{lem:gamma_pullback_image_iso, def:gamma_image_ring_equiv}
  The transport is \(\mathrm{Functor.mapIso}\) of \(((\mathrm{restrictFunctorIsoPullback}\ j).\mathrm{symm}.
  \mathrm{app}\ M)\), whose underlying map is the \(\mathrm{mapPresheaf}\) of a morphism of sheaves of
  modules; the module action on each side is the structure-sheaf action through the pullback's
  \(\mathrm{mapPresheaf}\), so the compatibility is the module-structure naturality of that morphism, read
  through the structure-sheaf identification \(\sigma_V\). The specific compatibility to invoke is the
  \(\mathrm{map\_smul}\) field of the underlying presheaf-of-modules morphism (its \(\mathrm{mapPresheaf}\)
  component is a module hom over the ring map), applied at the section \(x\) and scalar \(a\). Argue at the
  level of presheaf sections (term-mode applied \(\mathrm{map\_smul}\) / \(\mathrm{mapPresheaf}\)-naturality),
  \emph{not} by keyed rewriting under the modules instance diamond.
\end{proof}

\begin{lemma}\leanok
[$\mathrm{IsIso}\,\mathrm{fromTildeΓ}$ is invariant under isomorphism of modules]
  \label{lem:isIso_fromTildeΓ_of_iso}
  \lean{AlgebraicGeometry.Scheme.Modules.isIso_fromTildeΓ_of_iso}
  \uses{lem:isIso_fromTildeΓ_iff_isLocalizedModule_restrict}
  \textit{Project-bespoke transport, feeding the geometric Hfr producer.}
  Let \(M, M'\) be sheaves of modules on \(\Spec R\) with \(M \cong M'\). If
  \(M.\mathrm{fromTildeΓ}\) is an isomorphism, then so is \(M'.\mathrm{fromTildeΓ}\).
\end{lemma}
\begin{proof}
  \uses{lem:isIso_fromTildeΓ_iff_isLocalizedModule_restrict}
  By the criterion \(\mathrm{IsIso}\,M.\mathrm{fromTildeΓ} \iff M \in \mathrm{essImage}(\widetilde{(-)})\)
  (\cref{lem:isIso_fromTildeΓ_iff_isLocalizedModule_restrict}, Mathlib's \(\mathrm{isIso\_fromTildeΓ\_iff}\)),
  the hypothesis says \(M\) lies in the essential image of the tilde functor; the essential image is
  closed under isomorphism (\(\mathrm{Functor.essImage.ofIso}\)), so \(M'\) lies in it too, whence
  \(M'.\mathrm{fromTildeΓ}\) is an isomorphism.
\end{proof}

\begin{lemma}\leanok
[Section-localization descent from a local-tilde cover (gap1 keystone, D)]
  \label{lem:section_localization_descent}
  \lean{AlgebraicGeometry.Scheme.Modules.isLocalizedModule_basicOpen_descent}
  \uses{lem:section_localization_descent_of_basicOpen_cover,
    lem:section_localization_hfr_basicOpen,
    lem:isLocalizedModule_powers_transport, lem:pullback_gamma_top_iso,
    lem:isLocalizedModule_ringEquiv_semilinear,
    lem:isLocalizedModule_restrictScalars_powers_algebraMap,
    lem:gamma_pullback_image_iso_hom_semilinear,
    lem:isIso_fromTildeΓ_basicOpen_of_quasicoherent,
    lem:exists_finite_basicOpen_cover_le_quasicoherentData,
    lem:isLocalizedModule_restrict_of_isIso_fromTildeΓ,
    lem:isLocalizedModule_tilde_restrict}
  % NOTE: This NAMED keystone is the
  % basic-open cover form \cref{lem:section_localization_descent_of_basicOpen_cover} instantiated at the
  % quasi-coherence cover \cref{lem:exists_finite_basicOpen_cover_le_quasicoherentData}, with the
  % per-overlap Hfr supplied by the geometric producer \cref{lem:section_localization_hfr_basicOpen}.
  % NOTE: the general-U cover form \cref{lem:section_localization_descent_of_cover} is itself
  % axiom-clean, but its Hfr (a localization on EVERY open U ≤ D(r)) is NOT instantiable for an
  % arbitrary quasi-coherent M — a general open of the affine slice Spec R_r need not be
  % quasi-compact, so the section localization is unavailable there. The basic-open form is the one
  % gap1 uses: the sheaf-gluing engines only ever consult Hfr at basic opens D(s) and overlaps
  % D(r)⊓D(r') = D(r·r'), all quasi-compact.
  % isLocalizedModule_basicOpen_of_presentation does NOT close it: that lemma needs a GLOBAL
  % presentation of M, whereas here only per-D(r) tilde data is available.
  % SOURCE: [Stacks Project], "Properties of Schemes", \S "Sections over principal opens",
  % Lemma \texttt{lemma-invert-f-sections} (read from references/stacks-properties.tex, L2152-L2170).
  % Cf. [Hartshorne], II.5.3.
  % SOURCE QUOTE:
  % "Let $X$ be a scheme. Let $f \in \Gamma(X, \mathcal{O}_X)$.
  %  Denote $X_f \subset X$ the open where $f$ is invertible, see
  %  Schemes, Lemma \ref{schemes-lemma-f-open}.
  %  If $X$ is quasi-compact and quasi-separated, the canonical map
  %  $$
  %  \Gamma(X, \mathcal{O}_X)_f \longrightarrow \Gamma(X_f, \mathcal{O}_X)
  %  $$
  %  is an isomorphism. Moreover, if $\mathcal{F}$ is a quasi-coherent
  %  sheaf of $\mathcal{O}_X$-modules the map
  %  $$
  %  \Gamma(X, \mathcal{F})_f \longrightarrow \Gamma(X_f, \mathcal{F})
  %  $$
  %  is an isomorphism."
  \textit{Source: [Stacks Project], "Properties of Schemes", \S ``Sections over principal opens'',
  Lemma \texttt{lemma-invert-f-sections}; cf.\ [Hartshorne], II.5.3.}
  Let \(M\) be a quasi-coherent sheaf of modules on \(\Spec R\) (equivalently, by
  \cref{lem:isIso_fromTildeΓ_basicOpen_of_quasicoherent} and
  \cref{lem:exists_finite_basicOpen_cover_le_quasicoherentData}, \(M\) is locally isomorphic to a
  tilde \(\widetilde{N_j}\) on the members of a finite basic-open cover \(\{D(r_j)\}\) of
  \(\Spec R\)). Then for every \(f \in R\) the section restriction
  \(\Gamma(M, \top) \to \Gamma(M, D(f))\) exhibits the target as the localized module
  \((\mathrm{powers}\,f)^{-1}\,\Gamma(M, \top)\); that is, it is
  \(\mathrm{IsLocalizedModule}(\mathrm{powers}\,f)\) over \(R\).
\end{lemma}
% NOTE: the lemma STATEMENT is cited to Stacks lemma-invert-f-sections / Hartshorne II.5.3; the proof
% below is a project-internal descent (finite sheaf equalizer over the cover + flatness of
% localization) that deliberately does NOT route through the global affine equivalence (which is
% gap1 itself), so there is no transcribed `% SOURCE QUOTE PROOF`.
\begin{proof}
  \uses{lem:section_localization_descent_of_basicOpen_cover,
    lem:section_localization_hfr_basicOpen, lem:isLocalizedModule_powers_transport,
    lem:pullback_gamma_top_iso,
    lem:isIso_fromTildeΓ_basicOpen_of_quasicoherent,
    lem:exists_finite_basicOpen_cover_le_quasicoherentData,
    lem:isLocalizedModule_restrict_of_isIso_fromTildeΓ,
    lem:isLocalizedModule_tilde_restrict}
  Fix \(f \in R\). Refine the quasi-coherence cover to a finite basic-open cover
  \(\{D(r)\}_{r \in t}\) of \(\Spec R\), \(\mathrm{span}(t) = (1)\), with each \(D(r) \le q.X_{i(r)}\)
  (\cref{lem:exists_finite_basicOpen_cover_le_quasicoherentData}). It suffices to produce the
  \emph{basic-open} per-cover-element localization data \(\mathrm{Hfr}\) of
  \cref{lem:section_localization_descent_of_basicOpen_cover}: for every \(s \in R\) with
  \(D(s) \le D(r)\) for some \(r \in t\), the restriction
  \(\Gamma(M, D(s)) \to \Gamma(M, D(f) \sqcap D(s))\) is
  \(\mathrm{IsLocalizedModule}(\mathrm{powers}\,f)\). The basic-open cover form then assembles the three
  fields of the conclusion for the global restriction \(\Gamma(M, \top) \to \Gamma(M, D(f))\) --- its
  sheaf-gluing engines only ever consult \(\mathrm{Hfr}\) at basic opens \(D(s)\) and overlaps
  \(D(r) \sqcap D(r') = D(r \cdot r')\), all quasi-compact, so no localization on an arbitrary
  (possibly non-quasi-compact) open \(U \le D(r)\) is ever required.

  The basic-open datum \(\mathrm{Hfr}\) at \(D(s) \le D(r)\) is exactly the geometric producer
  \cref{lem:section_localization_hfr_basicOpen}: it transports P1's isomorphism counit
  \cref{lem:isIso_fromTildeΓ_basicOpen_of_quasicoherent} for the restriction of \(M\) to
  \(D(s) \cong \Spec R_s\) (realised as the pullback of \(M\) along the composite open immersion
  \(\Spec R_s \hookrightarrow \Spec R\)) into a localization at \(\mathrm{powers}\,f'\) over \(R_s\)
  (\(f'\) the image of \(f\)), via \cref{lem:isLocalizedModule_restrict_of_isIso_fromTildeΓ}; packages
  the section comparisons \cref{lem:pullback_gamma_top_iso} and \cref{lem:gamma_pullback_image_iso} as
  the transport \(\mathrm{AddEquiv}\)s \(\Gamma(M', \top) \cong \Gamma(M, D(s))\),
  \(\Gamma(M', D(f')) \cong \Gamma(M, D(f) \sqcap D(s))\); and reads the localization back as
  \(\mathrm{powers}\,f\) over \(R\) through the combiner \cref{lem:isLocalizedModule_powers_transport}.

  Feeding \(\mathrm{Hfr}\) and the cover \(\{D(r)\}_{r\in t}\) to
  \cref{lem:section_localization_descent_of_basicOpen_cover} yields that
  \(\Gamma(M, \top) \to \Gamma(M, D(f))\) is \(\mathrm{IsLocalizedModule}(\mathrm{powers}\,f)\) over
  \(R\), i.e.\ \((\mathrm{powers}\,f)^{-1}\,\Gamma(M, \top) \cong \Gamma(M, D(f))\). This is the
  \(\Gamma(X, \mathcal{F})_f \cong \Gamma(X_f, \mathcal{F})\) of Stacks
  \texttt{lemma-invert-f-sections} on \(X = \Spec R\), \(X_f = D(f)\), obtained from local-tilde data
  over a finite cover without appeal to the global affine equivalence.
\end{proof}

\begin{lemma}\leanok
[Quasi-coherent on an affine has trivial $\widetilde{(-)}$-counit (gap1)]
  \label{lem:qcoh_affine_isIso_fromTildeΓ}
  \lean{AlgebraicGeometry.Scheme.Modules.isIso_fromTildeΓ_of_isQuasicoherent}
  \uses{lem:section_localization_descent, lem:isIso_fromTildeΓ_iff_isLocalizedModule_restrict}
  % NOTE: gap1 is the PRIMARY stated
  % gap of the Quot-foundations layer — the QCoh(Spec R) ≃ Mod R essential-image gap (Stacks
  % lemma-quasi-coherent-affine). The assembly is a one-liner: the genuine content lives in the
  % section-localization descent
  % (\cref{lem:section_localization_descent}, keystone D), which is fed into the in-file equivalence
  % \cref{lem:isIso_fromTildeΓ_iff_isLocalizedModule_restrict} (already axiom-clean). The descent (D)
  % in turn rests on the slice-to-geometric bridge (\cref{lem:over_restrict_iso}, C) and the
  % per-element local-tilde transport (\cref{lem:isIso_fromTildeΓ_basicOpen_of_quasicoherent}, P1).
  % Route: analogies/quot-gap1-transport.md.
  % SOURCE: [Stacks Project], "Schemes", Lemma \texttt{lemma-quasi-coherent-affine}
  % (read from references/stacks-schemes.tex, L1279--1284).
  % SOURCE QUOTE: "Let $(X, \mathcal{O}_X) = (\Spec(R), \mathcal{O}_{\Spec(R)})$
  % be an affine scheme. Let $\mathcal{F}$ be a quasi-coherent $\mathcal{O}_X$-module. Then
  % $\mathcal{F}$ is isomorphic to the sheaf associated to the $R$-module $\Gamma(X,
  % \mathcal{F})$."
  \textit{The affine equivalence QCoh\((\Spec R)\) \(\simeq\) Mod\((R)\): a quasi-coherent module
  sheaf on \(\Spec R\) is \(\widetilde{\Gamma(M)}\) up to canonical isomorphism.}
  Let \(M\) be a sheaf of modules on \(\Spec R\) that is quasi-coherent
  (\(\mathrm{IsQuasicoherent}\,M\)). Then the counit \(M.\mathrm{fromTildeΓ}\) is an isomorphism;
  equivalently \(M\) lies in the essential image of \(\widetilde{(-)} : \mathrm{Mod}(R) \to (\Spec
  R).\mathrm{Modules}\).
\end{lemma}
% NOTE: the proof is a one-line assembly — feed the keystone descent (D) into the in-file
% equivalence isIso_fromTildeΓ_iff_isLocalizedModule_restrict — so it routes through no transcribed
% Stacks gluing proof; the SOURCE QUOTE above cites the lemma STATEMENT (Stacks
% lemma-quasi-coherent-affine).
\begin{proof}
  \uses{lem:section_localization_descent, lem:isIso_fromTildeΓ_iff_isLocalizedModule_restrict}
  Since \(M\) is quasi-coherent, the section-localization descent
  (\cref{lem:section_localization_descent}) gives that for every \(f \in R\) the section restriction
  \(\Gamma(M, \top) \to \Gamma(M, D(f))\) is \(\mathrm{IsLocalizedModule}(\mathrm{powers}\,f)\) over
  \(R\). This is exactly the right-hand side of the equivalence
  \cref{lem:isIso_fromTildeΓ_iff_isLocalizedModule_restrict} (already established in-file), whose
  left-hand side is \(\mathrm{IsIso}\,M.\mathrm{fromTildeΓ}\). Hence the counit is an isomorphism and
  \(M\) is identified with \(\widetilde{\Gamma(M, \top)}\), placing \(M\) in the essential image of
  \(\widetilde{(-)}\).
\end{proof}

\subsection{Section-transport producer for the basic-open Hfr}

This subsection builds the sole remaining ingredient of the gap1 keystone descent
\cref{lem:section_localization_descent}: the geometric \emph{producer}
\cref{lem:section_localization_hfr_basicOpen} that manufactures, for a quasi-coherent \(M\) on
\(\Spec R\) and a basic open \(D(s)\) inside a quasi-coherence cover member, the per-cover-element
localization datum \(\mathrm{Hfr}\) the basic-open cover form
\cref{lem:section_localization_descent_of_basicOpen_cover} consumes. The construction transports the
\(R_s\)-level localization living on the affine slice \(D(s) \cong \Spec R_s\) back to an
\(R\)-level statement through the combiner \cref{lem:isLocalizedModule_powers_transport}; the four
sub-lemmas below supply the geometric inputs (an iterated-pullback \(\mathrm{fromTildeΓ}\) iso, the
image/range computations, the \(\top\)-level semilinearity, and the scalar-tower instances).

\begin{definition}\leanok
[Composite basic-open immersion $j$]
  \label{def:composite_basic_open_immersion}
  \lean{AlgebraicGeometry.Scheme.Modules.compositeBasicOpenImmersion}
  \textit{Project-bespoke auxiliary: the composite open immersion identifying \(\Spec R_s\) with \(D(s)\).}
  For a cover member \(D(r) \le q.X_i\) and \(s : R\) with \(D(s) \le D(r)\), writing
  \(W = \iota_{q.X_i}^{-1}\,D(s)\), define
  \(j := \mathrm{isoSpec}.\mathrm{inv} \,\circ\, \iota_W \,\circ\, \iota_{q.X_i} :
  \Spec\,\Gamma(q.X_i, W) \to \Spec R\), where \(\mathrm{isoSpec}\) is the affine identification of the
  affine open \(W\) of \(q.X_i\). This is the morphism the producer sub-lemmas below are stated over.
\end{definition}

\begin{lemma}\leanok
[$j$ is an open immersion]
  \label{lem:composite_basic_open_immersion_isOpenImmersion}
  \lean{AlgebraicGeometry.Scheme.Modules.compositeBasicOpenImmersion_isOpenImmersion}
  \uses{def:composite_basic_open_immersion}
  \textit{Project-bespoke auxiliary.}
  The composite \(j\) of \cref{def:composite_basic_open_immersion} is an open immersion.
\end{lemma}
\begin{proof}
  \uses{def:composite_basic_open_immersion}
  \(j\) is the composite of the isomorphism \(\mathrm{isoSpec}.\mathrm{inv}\) (hence an open immersion) with
  the two open immersions \(\iota_W\) and \(\iota_{q.X_i}\); open immersions are closed under composition.
\end{proof}

\begin{lemma}\leanok
[Iterated-pullback $\mathrm{fromTildeΓ}$ iso along a composite immersion (producer, a)]
  \label{lem:pullback_composite_immersion_isIso_fromTildeΓ}
  \lean{AlgebraicGeometry.Scheme.Modules.pullback_composite_immersion_isIso_fromTildeΓ}
  \uses{def:composite_basic_open_immersion, lem:isIso_fromTildeΓ_of_iso,
    lem:isIso_fromTildeΓ_basicOpen_of_quasicoherent, lem:modules_pullback_mathlib}
  \textit{Project-bespoke: object-level input of the geometric Hfr producer.}
  Let \(M\) be quasi-coherent on \(\Spec R\), \(D(r)\) a cover member with \(D(r) \le q.X_i\), and
  \(s : R\) with \(D(s) \le D(r)\). Write \(W = \iota_{q.X_i}^{-1}\,D(s)\) and let
  \(j := \mathrm{isoSpec}.\mathrm{inv} \,\circ\, \iota_W \,\circ\, \iota_{q.X_i} : \Spec R_s \to \Spec R\)
  be the composite open immersion identifying \(\Spec R_s\) with \(D(s)\). Then the pullback module
  \((\mathrm{pullback}\,j).\mathrm{obj}\,M\) is isomorphic to the iterated pullback
  \((\mathrm{pullback}\,\mathrm{isoSpec}.\mathrm{inv}).\mathrm{obj}\bigl((\mathrm{pullback}\,\iota_W).
  \mathrm{obj}\bigl((\mathrm{pullback}\,\iota_{q.X_i}).\mathrm{obj}\,M\bigr)\bigr)\), and
  \(\mathrm{IsIso}\bigl(((\mathrm{pullback}\,j).\mathrm{obj}\,M).\mathrm{fromTildeΓ}\bigr)\).
\end{lemma}
\begin{proof}
  \uses{def:composite_basic_open_immersion, lem:isIso_fromTildeΓ_of_iso,
    lem:isIso_fromTildeΓ_basicOpen_of_quasicoherent, lem:modules_pullback_mathlib}
  The iso of pullback modules is the \(\mathrm{pullbackComp}\) coherence of the
  \(\mathrm{Scheme.Modules.pullback}\) pseudofunctor (\cref{lem:modules_pullback_mathlib}) applied to the
  three-fold composite \(j = \mathrm{isoSpec}.\mathrm{inv} \circ \iota_W \circ \iota_{q.X_i}\). The
  P1 keystone \cref{lem:isIso_fromTildeΓ_basicOpen_of_quasicoherent} supplies
  \(\mathrm{IsIso}\,\mathrm{fromTildeΓ}\) for the iterated pullback (its very construction restricts
  \(M\) along these three stages); transporting that across the composite iso by
  \cref{lem:isIso_fromTildeΓ_of_iso} gives
  \(\mathrm{IsIso}\bigl(((\mathrm{pullback}\,j).\mathrm{obj}\,M).\mathrm{fromTildeΓ}\bigr)\).
\end{proof}

\begin{lemma}\leanok
[Range of the composite immersion (producer, b-range)]
  \label{lem:composite_immersion_range_basicOpen}
  \lean{AlgebraicGeometry.Scheme.Modules.compositeBasicOpenImmersion_opensRange}
  \uses{def:composite_basic_open_immersion, lem:composite_basic_open_immersion_isOpenImmersion}
  \textit{Project-bespoke: range bookkeeping of the geometric Hfr producer.}
  With \(j\) the composite immersion of \cref{def:composite_basic_open_immersion},
  \(j.\mathrm{opensRange} = D(s)\).
\end{lemma}
\begin{proof}
  \uses{def:composite_basic_open_immersion}
  The range is the open-immersion range of the composite of the basic-open immersion \(\iota_W\), the
  affine-chart immersion \(\iota_{q.X_i}\), and the affine identification \(\mathrm{isoSpec}.\mathrm{inv}\)
  (each an open immersion onto a basic / affine open); its composite image is the intersection
  \((q.X_i) \sqcap D(s) = D(s)\) (using \(D(s) \le q.X_i\)).
\end{proof}

\begin{lemma}[$f$-locus image of the composite immersion (producer, b-flocus)]
  \label{lem:composite_immersion_flocus_basicOpen}
  % Deliberate duplicate pin: this narrative/grouping block has no standalone Lean decl (its two
  % claims are absorbed inline into section_localization_hfr_basicOpen, per the NOTE below). Pinned
  % to the realizing decl so criterion-3 (\lean{}) holds and the node stays wired; leandag matches
  % by name, so a duplicate pin keeps every \uses{} edge intact (iter-027 precedent).
  \lean{AlgebraicGeometry.Scheme.Modules.section_localization_hfr_basicOpen}
  \uses{lem:composite_immersion_range_basicOpen, def:gamma_image_ring_equiv,
    lem:compositeBasicOpenImmersion_image_basicOpen, lem:image_basicOpen_eq_inf}
  % NOTE: no standalone Lean decl pins this block; its content is absorbed inline into
  % section_localization_hfr_basicOpen. It bundles TWO claims, each now carried by a named decl /
  % inline step: (a) σ(f') = algebraMap R R_s f is discharged inline as σ.apply_symm_apply (the choice
  % f' := σ⁻¹(algebraMap R R_s f)); (b) j ''ᵁ D(f') = D(f) ⊓ D(s) is computed via
  % \cref{lem:compositeBasicOpenImmersion_image_basicOpen} + \cref{lem:image_basicOpen_eq_inf}. The
  % block is retained as the mathematical record of the two claims; the real decls live in the two
  % cited blocks.
  \textit{Project-bespoke: $f$-locus image / $\sigma$ bookkeeping of the geometric Hfr producer.}
  With \(j\) the composite immersion of \cref{def:composite_basic_open_immersion}, let
  \(\sigma = \mathrm{gammaImageRingEquiv}\,j\,\top : \Gamma(\mathcal{O}_{\Spec R_s}, \top) \xrightarrow{\sim}
  \Gamma(\mathcal{O}_{\Spec R}, j\,''^{U}\top)\) be the \(\top\)-level structure-sheaf ring iso
  (\cref{def:gamma_image_ring_equiv}), and put \(f' := \sigma^{-1}(\mathrm{algebraMap}\,R\,R_s\,f)\). Then
  \[
    \sigma(f') = \mathrm{algebraMap}\,R\,R_s\,f, \qquad j\,''^{U}\,D(f') = D(f) \sqcap D(s) .
  \]
\end{lemma}
\begin{proof}
  \uses{lem:composite_immersion_range_basicOpen, def:gamma_image_ring_equiv,
    lem:compositeBasicOpenImmersion_image_basicOpen, lem:image_basicOpen_eq_inf}
  The identity \(\sigma(f') = \mathrm{algebraMap}\,R\,R_s\,f\) is definitional from
  \(f' = \sigma^{-1}(\mathrm{algebraMap}\,R\,R_s\,f)\). For the image: by
  \cref{lem:compositeBasicOpenImmersion_image_basicOpen} the immersion \(j\) carries the basic open
  \(D(f')\) to the \(\Spec R\) basic open of the transported section, and
  \cref{lem:image_basicOpen_eq_inf} rewrites that image as
  \((j\,''^{U}\,\top) \sqcap D(\sigma(f'))\); using the range
  \(j.\mathrm{opensRange} = D(s)\) (\cref{lem:composite_immersion_range_basicOpen}) and
  \(\sigma(f') = \mathrm{algebraMap}\,R\,R_s\,f\) (whose non-vanishing locus is \(D(f)\)) gives
  \(D(f) \sqcap D(s)\).
\end{proof}

\begin{lemma}\leanok
[$\top$-level semilinearity of the section transport (producer, c)]
  \label{lem:gamma_image_iso_semilinear_top}
  % Deliberate duplicate pin: absorbed inline into section_localization_hfr_aux (no standalone Lean
  % decl, per the NOTE below). Pinned to the realizing decl for criterion-3; leandag matches by name.
  \lean{AlgebraicGeometry.Scheme.Modules.section_localization_hfr_aux}
  \uses{lem:gamma_pullback_image_iso_hom_semilinear, lem:gamma_pullback_image_iso_hom_naturality,
    def:gamma_image_ring_equiv}
  % NOTE: absorbed inline into section_localization_hfr_aux (\cref{lem:section_localization_hfr_aux});
  % no standalone Lean decl. Its content is the `he₁`/`he₂` semilinearity hypotheses inside that core,
  % proved via gamma_pullback_image_iso_hom_semilinear + appIso_inv_naturality. Retained as the
  % mathematical record; the opaque core cites it through \cref{lem:section_localization_hfr_aux}.
  \textit{Project-bespoke: upgrades the $D(f')$-level semilinearity to the $\top$-level $\sigma$.}
  The section transports \cref{lem:gamma_pullback_image_iso} (forward maps), which are semilinear over
  the \(D(f')\)-level ring iso \(\mathrm{gammaImageRingEquiv}\,j\,D(f')\)
  (\cref{lem:gamma_pullback_image_iso_hom_semilinear}), are also semilinear over the \(\top\)-level ring
  iso \(\sigma = \mathrm{gammaImageRingEquiv}\,j\,\top\) of
  \cref{lem:composite_immersion_flocus_basicOpen}, as required for the \(R_s\)-actions fed to the
  combiner.
\end{lemma}
\begin{proof}
  \uses{lem:gamma_pullback_image_iso_hom_semilinear, lem:gamma_pullback_image_iso_hom_naturality,
    def:gamma_image_ring_equiv}
  The structure-sheaf restriction \(\Gamma(\mathcal{O}_{\Spec R_s}, \top) \to \Gamma(\mathcal{O}_{\Spec
  R_s}, D(f'))\) intertwines the two ring isos \(\sigma\) and \(\mathrm{gammaImageRingEquiv}\,j\,D(f')\)
  (naturality of \(\mathrm{gammaImageRingEquiv}\) in the open argument), and the section transports at
  \(\top\) and \(D(f')\) commute with the module restriction maps
  (\cref{lem:gamma_pullback_image_iso_hom_naturality}); transporting the \(D(f')\)-level semilinearity
  \cref{lem:gamma_pullback_image_iso_hom_semilinear} along this restriction yields the \(\top\)-level
  semilinearity over \(\sigma\).
\end{proof}

\begin{lemma}\leanok
[$A$-module and scalar-tower instances on the $f$-locus sections (producer, d)]
  \label{lem:flocus_section_scalar_tower}
  % Deliberate duplicate pin: absorbed inline into section_localization_hfr_aux (no standalone Lean
  % decl, per the NOTE below). Pinned to the realizing decl for criterion-3; leandag matches by name.
  \lean{AlgebraicGeometry.Scheme.Modules.section_localization_hfr_aux}
  \uses{lem:composite_immersion_flocus_basicOpen}
  % NOTE: absorbed inline into section_localization_hfr_aux (\cref{lem:section_localization_hfr_aux});
  % no standalone Lean decl. The A-module and IsScalarTower instances are constructed inline in that
  % core via Module.compHom + IsScalarTower.of_algebraMap_smul. Retained as the mathematical record;
  % the opaque core cites it through \cref{lem:section_localization_hfr_aux}.
  \textit{Project-bespoke: the $A$-module / $\mathrm{IsScalarTower}$ instances feeding the combiner.}
  Let \(A = \Gamma(M, D(s))\) and consider the \(f\)-locus sections \(\Gamma(M, j\,''^{U}\,D(f')) =
  \Gamma(M, D(f) \sqcap D(s))\). These carry an \(A\)-module structure and an \(R\)-module structure with
  \(\mathrm{IsScalarTower}\,R\,A\,\Gamma(M, D(f) \sqcap D(s))\), where the \(A\)-action is by restriction
  of the structure-sheaf action along \(\Gamma(\mathcal{O}_{\Spec R}, D(s)) \to \Gamma(\mathcal{O}_{\Spec
  R}, D(f) \sqcap D(s))\), and likewise for the source \(\Gamma(M, D(s))\).
\end{lemma}
\begin{proof}
  \uses{lem:composite_immersion_flocus_basicOpen}
  The module \(\Gamma(M, D(f) \sqcap D(s))\) is a module over the section ring
  \(\Gamma(\mathcal{O}_{\Spec R}, D(f) \sqcap D(s))\); restriction of scalars along the structure-sheaf
  restriction \(\Gamma(\mathcal{O}_{\Spec R}, D(s)) \to \Gamma(\mathcal{O}_{\Spec R}, D(f) \sqcap D(s))\)
  makes it an \(A\)-module, and along \(R \to \Gamma(\mathcal{O}_{\Spec R}, D(s)) = A\) an \(R\)-module;
  the scalar tower \(\mathrm{IsScalarTower}\,R\,A\,\Gamma(M, D(f) \sqcap D(s))\) is the factorization of
  the two ring maps. The identification of the locus \(j\,''^{U}\,D(f') = D(f) \sqcap D(s)\) is
  \cref{lem:composite_immersion_flocus_basicOpen}.
\end{proof}

\begin{lemma}\leanok
[Image of an affine basic open under an open immersion of affines]
  \label{lem:image_basicOpen_of_affine}
  \lean{AlgebraicGeometry.Scheme.Modules.image_basicOpen_of_affine}
  \textit{Project-bespoke: pure-geometry image helper, stated with the immersion opaque.}
  Let \(j : \Spec S \to \Spec R\) be an open immersion and \(f' \in S\). Then the image of the basic
  open \(D(f')\) under \(j\) is the \(\Spec R\) scheme basic open of the transported global section
  \((j.\mathrm{appIso}\,\top)^{-1}\bigl((\Gamma\Spec\text{-iso}_S)^{-1}\,f'\bigr)\):
  \[
    j\,''^{U}\,D(f') = (\Spec R).\mathrm{basicOpen}\,
      \bigl((j.\mathrm{appIso}\,\top)^{-1}\,((\Gamma\Spec\text{-iso}_S)^{-1}\,f')\bigr).
  \]
\end{lemma}
\begin{proof}
  Rewrite \(D(f')\) as the \(\Spec S\) scheme basic open of the global structure section
  \((\Gamma\Spec\text{-iso}_S)^{-1}\,f'\) (Mathlib's \(\mathrm{basicOpen\_eq\_of\_affine}\)), then push
  it across \(j\) by the scheme-image-of-a-basic-open identity (Mathlib's
  \(\mathrm{Scheme.image\_basicOpen}\)), which sends a basic open to the basic open of the
  \(\mathrm{appIso}\)-transported section. The immersion \(j\) is kept abstract so the rewrite does not
  unfold a concrete composite.
\end{proof}

\begin{lemma}\leanok
[Image of the basic open under the composite immersion (instantiation)]
  \label{lem:compositeBasicOpenImmersion_image_basicOpen}
  \lean{AlgebraicGeometry.Scheme.Modules.compositeBasicOpenImmersion_image_basicOpen}
  \uses{lem:image_basicOpen_of_affine, def:composite_basic_open_immersion}
  \textit{Project-bespoke: \cref{lem:image_basicOpen_of_affine} at the concrete composite immersion.}
  With \(j = \) the composite immersion of \cref{def:composite_basic_open_immersion} and \(f'\) a slice
  element, the image \(j\,''^{U}\,D(f')\) is the \(\Spec R\) basic open of the transported section
  \((j.\mathrm{appIso}\,\top)^{-1}\,((\Gamma\Spec\text{-iso})^{-1}\,f')\). This is
  \cref{lem:image_basicOpen_of_affine} specialised to \(j = \mathrm{compositeBasicOpenImmersion}\).
\end{lemma}
\begin{proof}
  \uses{lem:image_basicOpen_of_affine}
  Immediate instantiation of \cref{lem:image_basicOpen_of_affine} at the open immersion
  \(j = \mathrm{compositeBasicOpenImmersion}\).
\end{proof}

\begin{lemma}\leanok
[Image of an affine basic open as an intersection with the range]
  \label{lem:image_basicOpen_eq_inf}
  \lean{AlgebraicGeometry.Scheme.Modules.image_basicOpen_eq_inf}
  \uses{lem:image_basicOpen_of_affine}
  \textit{Project-bespoke: rewrites the $f'$-locus image as range $\sqcap$ basic open.}
  Let \(j : \Spec S \to \Spec R\) be an open immersion, \(f' \in S\), and \(g \in \Gamma(\Spec R, \top)\)
  a global section whose restriction to \(j\,''^{U}\,\top\) equals the \(\mathrm{appIso}\)-transport of
  \(f'\). Then
  \[
    j\,''^{U}\,D(f') = (j\,''^{U}\,\top) \sqcap (\Spec R).\mathrm{basicOpen}\,g .
  \]
\end{lemma}
\begin{proof}
  \uses{lem:image_basicOpen_of_affine}
  By \cref{lem:image_basicOpen_of_affine}, \(j\,''^{U}\,D(f')\) is the basic open of the
  \(\mathrm{appIso}\)-transport of \(f'\); the transport hypothesis identifies that transport with the
  restriction of \(g\) to \(j\,''^{U}\,\top\), and the structure-sheaf basic-open-of-a-restriction
  identity (Mathlib's \(\mathrm{Scheme.basicOpen\_res}\)) turns the basic open of a restricted section
  into the intersection of the restriction domain \(j\,''^{U}\,\top\) with the basic open of \(g\).
\end{proof}

\begin{lemma}\leanok
[Basic-open Hfr along an abstract affine open immersion (opaque-$j$ core)]
  \label{lem:section_localization_hfr_aux}
  \lean{AlgebraicGeometry.Scheme.Modules.section_localization_hfr_aux}
  \uses{lem:gamma_image_iso_semilinear_top, lem:flocus_section_scalar_tower,
    lem:isLocalizedModule_powers_transport, lem:isLocalizedModule_restrict_of_isIso_fromTildeΓ,
    lem:gamma_pullback_image_iso, lem:gamma_pullback_image_iso_hom_semilinear,
    lem:gamma_pullback_image_iso_hom_naturality, def:gamma_image_ring_equiv}
  \textit{Project-bespoke: the abstract proof engine of the geometric Hfr producer; it absorbs the
  $f$-locus/semilinearity/scalar-tower content of the three preceding producer lemmas.}
  Let \(M\) be a sheaf of modules on \(\Spec R\) and \(j : \Spec S \to \Spec R\) an open immersion with
  \(\mathrm{IsIso}\bigl(((\mathrm{pullback}\,j).\mathrm{obj}\,M).\mathrm{fromTildeΓ}\bigr)\) (the P1
  datum). Fix \(f \in R\), a slice element \(f' \in S\), and target opens \(V \le U\) with
  \(U = j\,''^{U}\,\top\) and \(V = j\,''^{U}\,D(f')\), such that the \(\mathrm{appIso}\)-transport of
  \(f'\) is the restriction of \(\mathrm{algebraMap}\,R\,(\cdot)\) of \(f\) to \(j\,''^{U}\,\top\). Then
  the section restriction \(\Gamma(M, U) \to \Gamma(M, V)\) is
  \(\mathrm{IsLocalizedModule}(\mathrm{powers}\,f)\) over \(R\).
\end{lemma}
\begin{proof}
  \uses{lem:gamma_image_iso_semilinear_top, lem:flocus_section_scalar_tower,
    lem:isLocalizedModule_powers_transport, lem:isLocalizedModule_restrict_of_isIso_fromTildeΓ,
    lem:gamma_pullback_image_iso, lem:gamma_pullback_image_iso_hom_semilinear,
    lem:gamma_pullback_image_iso_hom_naturality, def:gamma_image_ring_equiv}
  Set \(M' = (\mathrm{pullback}\,j).\mathrm{obj}\,M\); the P1 hypothesis gives
  \(\mathrm{IsIso}\,M'.\mathrm{fromTildeΓ}\). The affine engine
  \cref{lem:isLocalizedModule_restrict_of_isIso_fromTildeΓ} localises the slice restriction
  \(\Gamma(M', \top) \to \Gamma(M', D(f'))\) at \(\mathrm{powers}\,f'\) over the slice ring \(S\). The
  \(\sigma\)-semilinear additive isomorphisms
  \(e_1 = \mathrm{gammaPullbackImageIso}\,j\,M\,\top\) and
  \(e_2 = \mathrm{gammaPullbackImageIso}\,j\,M\,D(f')\) (\cref{lem:gamma_pullback_image_iso}), over the
  \(\top\)-level structure ring iso
  \(\sigma = (\Gamma\Spec\text{-iso}_S)^{-1} \circ \mathrm{gammaImageRingEquiv}\,j\,\top\)
  (\cref{def:gamma_image_ring_equiv}), are semilinear by
  \cref{lem:gamma_image_iso_semilinear_top} (the \(\mathrm{he}_1/\mathrm{he}_2\) hypotheses, proved here
  inline from \cref{lem:gamma_pullback_image_iso_hom_semilinear} and the \(\mathrm{appIso}\) naturality
  square), and the section restriction intertwines them with the slice localization map by
  \cref{lem:gamma_pullback_image_iso_hom_naturality}. The target sections carry the \(A\)-module and
  \(R\)-module structures with scalar tower of \cref{lem:flocus_section_scalar_tower} (built here inline
  via \(\mathrm{Module.compHom}\) and \(\mathrm{IsScalarTower.of\_algebraMap\_smul}\)). Feeding
  \(\sigma\), \(f\), \(f'\), the slice localization, the semilinear pair \((e_1, e_2)\) and the
  intertwining restriction to the algebra combiner \cref{lem:isLocalizedModule_powers_transport}
  descends the localization to \(\mathrm{powers}\,f\) over \(R\). Finally the \(j\,''^{U}\)-form opens
  are transported to \(U\) and \(V\) by \(\mathrm{eqToHom}\) open isomorphisms and
  \(\mathrm{IsLocalizedModule.of\_linearEquiv}\).

  \emph{Load-bearing implementation lesson.} The immersion \(j\) must be kept \emph{opaque} in this
  helper: instantiating it at the concrete composite immersion before the final form-coercion makes the
  closing \(\mathrm{whnf}\) unfold the concrete section rings and blows past \(3.2\)M heartbeats. Stated
  with \(j\) abstract, the section-ring defeqs stay cheap and the proof closes within
  \(1.6\)M heartbeats; the concrete immersion is supplied only by the thin wrapper
  \cref{lem:section_localization_hfr_basicOpen}.
\end{proof}

\begin{lemma}\leanok
[The basic-open Hfr datum from section transport (geometric producer, TOP)]
  \label{lem:section_localization_hfr_basicOpen}
  \lean{AlgebraicGeometry.Scheme.Modules.section_localization_hfr_basicOpen}
  \uses{lem:section_localization_hfr_aux, lem:pullback_composite_immersion_isIso_fromTildeΓ,
    lem:composite_immersion_range_basicOpen, lem:image_basicOpen_eq_inf,
    lem:composite_immersion_flocus_basicOpen, def:gamma_image_ring_equiv,
    def:composite_basic_open_immersion}
  % NOTE: Implemented as a thin wrapper over the opaque-immersion core
  % \cref{lem:section_localization_hfr_aux} instantiated at
  % j = compositeBasicOpenImmersion; the heavy assembly (semilinearity, scalar tower, combiner) lives
  % in that core, so this wrapper only supplies the concrete immersion, the P1 datum, f', and the
  % opens identifications.
  \textit{Project-bespoke: the producer manufacturing the basic-open \(\mathrm{Hfr}\) datum from the
  per-element P1 transport; the sole remaining build of \cref{lem:section_localization_descent}.}
  Let \(M\) be a quasi-coherent sheaf of modules on \(\Spec R\), \(D(r)\) a cover member with
  \(D(r) \le q.X_i\), and \(s : R\) with \(D(s) \le D(r)\). Then the section restriction
  \(\Gamma(M, D(s)) \to \Gamma(M, D(f) \sqcap D(s))\) is
  \(\mathrm{IsLocalizedModule}(\mathrm{powers}\,f)\) over \(R\) --- i.e.\ the basic-open
  \(\mathrm{Hfr}\) datum demanded by \cref{lem:section_localization_descent_of_basicOpen_cover}.
\end{lemma}
\begin{proof}
  \uses{lem:section_localization_hfr_aux, lem:pullback_composite_immersion_isIso_fromTildeΓ,
    lem:composite_immersion_range_basicOpen, lem:image_basicOpen_eq_inf,
    lem:composite_immersion_flocus_basicOpen, def:gamma_image_ring_equiv}
  This is the thin instantiation of the opaque-\(j\) core
  \cref{lem:section_localization_hfr_aux} at the concrete composite immersion
  \(j = \) the immersion of \cref{def:composite_basic_open_immersion} identifying \(\Spec R_s\) with
  \(D(s)\). We discharge the four hypotheses of the core as follows.
  \begin{itemize}
    \item \emph{P1 datum.} \cref{lem:pullback_composite_immersion_isIso_fromTildeΓ} supplies
      \(\mathrm{IsIso}\bigl(((\mathrm{pullback}\,j).\mathrm{obj}\,M).\mathrm{fromTildeΓ}\bigr)\).
    \item \emph{The element \(f'\).} Take
      \(f' := \sigma^{-1}(\mathrm{algebraMap}\,R\,R_s\,f)\) with
      \(\sigma = (\Gamma\Spec\text{-iso})^{-1} \circ \mathrm{gammaImageRingEquiv}\,j\,\top\)
      (\cref{def:gamma_image_ring_equiv}); then the \(\mathrm{appIso}\)-transport hypothesis is exactly
      \(\sigma(f') = \mathrm{algebraMap}\,R\,R_s\,f\), which holds by \(\sigma.\mathrm{apply\_symm\_apply}\)
      (claim (a) of \cref{lem:composite_immersion_flocus_basicOpen}).
    \item \emph{Identification of the opens.} \(j\,''^{U}\,\top = D(s)\) by
      \cref{lem:composite_immersion_range_basicOpen}; and
      \(j\,''^{U}\,D(f') = D(f) \sqcap D(s)\) by \cref{lem:image_basicOpen_eq_inf} (claim (b) of
      \cref{lem:composite_immersion_flocus_basicOpen}).
  \end{itemize}
  Feeding these to \cref{lem:section_localization_hfr_aux} yields that the section restriction
  \(\Gamma(M, D(s)) \to \Gamma(M, D(f) \sqcap D(s))\) is
  \(\mathrm{IsLocalizedModule}(\mathrm{powers}\,f)\) over \(R\), which is the basic-open
  \(\mathrm{Hfr}\) datum. (The heavy assembly --- the \(\sigma\)-semilinear section transports, the
  scalar-tower instances, and the localization combiner --- is carried entirely by the core, which is
  why the immersion is kept opaque there; see \cref{lem:section_localization_hfr_aux}.)
\end{proof}

\subsection{G1-assemble: from per-basic-open section localization to the
\texorpdfstring{$\widetilde{(-)}$}{tilde}-counit}

This subsection records the bridge \(\text{section-localization on every basic open} \Rightarrow
\mathrm{IsIso}\,M.\mathrm{fromTildeΓ}\): given that the section restriction localizes on every basic
open, the tilde-Gamma counit is an isomorphism. Together with its converse
\cref{lem:isLocalizedModule_restrict_of_isIso_fromTildeΓ} it assembles the equivalence
\cref{lem:isIso_fromTildeΓ_iff_isLocalizedModule_restrict}, which is exactly what makes G1-core
(\cref{lem:qcoh_affine_section_localization}) and gap1
(\cref{lem:qcoh_affine_isIso_fromTildeΓ}) interderivable -- the route by which G1-core is obtained as
a corollary of gap1. The supporting Mathlib facts are recorded as dependency anchors.

% --- Mathlib dependency anchors for the G1-assemble step ---

\begin{lemma}[Basic opens form a basis of $\Spec R$]
  \label{lem:isBasis_basic_opens_mathlib}
  \lean{PrimeSpectrum.isBasis_basic_opens}
  \mathlibok
  \textit{Provided by Mathlib.}
  The basic opens \(D(f)\), \(f \in R\), form a basis for the Zariski topology of \(\Spec R\).
\end{lemma}

\begin{lemma}[Stalkwise injectivity from a basis]
  \label{lem:stalkFunctor_map_injective_of_isBasis_mathlib}
  \lean{TopCat.Presheaf.stalkFunctor_map_injective_of_isBasis}
  \mathlibok
  \textit{Provided by Mathlib.}
  If a morphism of presheaves is injective on the sections over every open in a basis, then the
  induced map on each stalk is injective.
\end{lemma}

\begin{lemma}[Isomorphism from stalkwise isomorphism]
  \label{lem:isIso_of_stalkFunctor_map_iso_mathlib}
  \lean{TopCat.Presheaf.isIso_of_stalkFunctor_map_iso}
  \mathlibok
  \textit{Provided by Mathlib.}
  A morphism of sheaves that induces an isomorphism on every stalk is an isomorphism.
\end{lemma}

\begin{lemma}[The canonical isomorphism of two localizations]
  \label{lem:isLocalizedModule_linearEquiv_mathlib}
  \lean{IsLocalizedModule.linearEquiv}
  \mathlibok
  \textit{Provided by Mathlib.}
  Two localizations \(g : M \to M'\) and \(g' : M \to M''\) of the same module at the same submonoid
  \(S\) are canonically isomorphic by a linear equivalence \(M' \cong M''\) compatible with \(g\)
  and \(g'\).
\end{lemma}

\begin{lemma}[Extensionality for maps out of a localized module]
  \label{lem:isLocalizedModule_linearMap_ext_mathlib}
  \lean{IsLocalizedModule.linearMap_ext}
  \mathlibok
  \textit{Provided by Mathlib.}
  Two linear maps out of a localized module agree as soon as they agree after precomposition with
  the localization map.
\end{lemma}

% --- Project-original G1-assemble nodes (Lean realizations) ---

\begin{lemma}\leanok
[A map intertwining two localizations is bijective]
  \label{lem:bijective_comp_of_localizations}
  \lean{AlgebraicGeometry.bijective_comp_of_localizations}
  % NOTE: the pinned Lean decl is `private` in the source (a project-internal helper for the
  % G1-assemble step); the \lean{} pin is recorded for the DAG but resolves only project-internally.
  \uses{lem:isLocalizedModule_linearEquiv_mathlib, lem:isLocalizedModule_linearMap_ext_mathlib}
  \textit{Linear-algebra glue for the G1-assemble step.}
  Let \(S \subseteq R\) be a submonoid and let \(g : M \to M'\), \(g' : M \to M''\) both exhibit a
  localization of \(M\) at \(S\). If a linear map \(h : M' \to M''\) satisfies \(h \circ g = g'\),
  then \(h\) is bijective.
\end{lemma}
\begin{proof}
  \uses{lem:isLocalizedModule_linearEquiv_mathlib, lem:isLocalizedModule_linearMap_ext_mathlib}
  The map \(h\) agrees with the canonical localization isomorphism
  \(\mathrm{IsLocalizedModule.linearEquiv}\,S\,g\,g'\)
  (\cref{lem:isLocalizedModule_linearEquiv_mathlib}): both compose with \(g\) to give \(g'\), so they
  coincide by uniqueness of maps out of a localization
  (\cref{lem:isLocalizedModule_linearMap_ext_mathlib}). A linear equivalence is bijective, hence so
  is \(h\).
\end{proof}

\begin{lemma}\leanok
[Basis-wise isomorphism of sheaves on $\Spec R$]
  \label{lem:isIso_sheaf_of_isIso_app_basicOpen}
  \lean{AlgebraicGeometry.isIso_sheaf_of_isIso_app_basicOpen}
  % NOTE: the pinned Lean decl is `private` in the source (a project-internal helper for the
  % G1-assemble step); the \lean{} pin is recorded for the DAG but resolves only project-internally.
  \uses{lem:isBasis_basic_opens_mathlib, lem:stalkFunctor_map_injective_of_isBasis_mathlib,
    lem:isIso_of_stalkFunctor_map_iso_mathlib}
  \textit{Upgrades a per-basic-open isomorphism to an isomorphism of sheaves.}
  Let \(\alpha : F \to G\) be a morphism of sheaves of \(R\)-modules on \(\Spec R\) whose component
  on every basic open \(D(f)\) is an isomorphism. Then \(\alpha\) is an isomorphism.
\end{lemma}
\begin{proof}
  \uses{lem:isBasis_basic_opens_mathlib, lem:stalkFunctor_map_injective_of_isBasis_mathlib,
    lem:isIso_of_stalkFunctor_map_iso_mathlib}
  The basic opens form a basis of \(\Spec R\) (\cref{lem:isBasis_basic_opens_mathlib}). Bijectivity
  of every component \(\alpha_{D(f)}\) gives, in particular, injectivity on each basis open, hence
  injectivity on every stalk (\cref{lem:stalkFunctor_map_injective_of_isBasis_mathlib});
  surjectivity on stalks follows by lifting germs through the bijective basic-open components. A
  morphism that is an isomorphism on every stalk is an isomorphism
  (\cref{lem:isIso_of_stalkFunctor_map_iso_mathlib}).
\end{proof}

\begin{lemma}\leanok
[$\mathrm{IsIso}\,M.\mathrm{fromTildeΓ}$ from per-basic-open section localization
(G1-assemble)]
  \label{lem:isIso_fromTildeΓ_of_isLocalizedModule_restrict}
  \lean{AlgebraicGeometry.isIso_fromTildeΓ_of_isLocalizedModule_restrict}
  \uses{lem:isIso_sheaf_of_isIso_app_basicOpen, lem:bijective_comp_of_localizations,
    lem:isLocalizedModule_tilde_restrict}
  \textit{The G1-assemble reduction: section-localization on every basic open forces the
  \(\widetilde{(-)}\)-counit to be an isomorphism.}
  Let \(M\) be a sheaf of modules on \(\Spec R\), and suppose that for every \(f \in R\) the section
  restriction \(\Gamma(M, \top) \to \Gamma(M, D(f))\) exhibits the target as
  \(\mathrm{IsLocalizedModule}(\mathrm{powers}\,f)\) over \(R\) -- exactly the conclusion of G1-core
  (\cref{lem:qcoh_affine_section_localization}). Then the counit \(M.\mathrm{fromTildeΓ}\) is an
  isomorphism.
\end{lemma}
\begin{proof}
  \uses{lem:isIso_sheaf_of_isIso_app_basicOpen, lem:bijective_comp_of_localizations,
    lem:isLocalizedModule_tilde_restrict}
  Write \(N = \Gamma(M, \top)\). On each basic open \(D(f)\), the component of
  \(\mathrm{modulesSpecToSheaf}(M.\mathrm{fromTildeΓ})\) is a map between two localizations of \(N\)
  at \(\mathrm{powers}\,f\): the source \(\Gamma(\widetilde N, D(f))\) localizes \(N\) via the affine
  tilde map (\cref{lem:isLocalizedModule_tilde_restrict}) and the target \(\Gamma(M, D(f))\)
  localizes \(N\) by hypothesis, the two being intertwined by the compatibility
  \(\mathrm{toOpen\_fromTildeΓ\_app}\). Hence the component is bijective
  (\cref{lem:bijective_comp_of_localizations}), so an isomorphism. By
  \cref{lem:isIso_sheaf_of_isIso_app_basicOpen} the sheaf morphism
  \(\mathrm{modulesSpecToSheaf}(M.\mathrm{fromTildeΓ})\) is then an isomorphism, and since
  \(\mathrm{modulesSpecToSheaf}\) is fully faithful it reflects isomorphisms, giving
  \(\mathrm{IsIso}\,M.\mathrm{fromTildeΓ}\).
\end{proof}

\begin{lemma}\leanok
[$\mathrm{IsIso}\,M.\mathrm{fromTildeΓ}$ characterized by section localization]
  \label{lem:isIso_fromTildeΓ_iff_isLocalizedModule_restrict}
  \lean{AlgebraicGeometry.isIso_fromTildeΓ_iff_isLocalizedModule_restrict}
  \uses{lem:isIso_fromTildeΓ_of_isLocalizedModule_restrict,
    lem:isLocalizedModule_restrict_of_isIso_fromTildeΓ}
  \textit{The two engines packaged as an equivalence.}
  For a sheaf of modules \(M\) on \(\Spec R\), the counit \(M.\mathrm{fromTildeΓ}\) is an
  isomorphism if and only if for every \(f \in R\) the section restriction
  \(\Gamma(M, \top) \to \Gamma(M, D(f))\) exhibits the target as
  \(\mathrm{IsLocalizedModule}(\mathrm{powers}\,f)\) over \(R\).
\end{lemma}
\begin{proof}
  \uses{lem:isIso_fromTildeΓ_of_isLocalizedModule_restrict,
    lem:isLocalizedModule_restrict_of_isIso_fromTildeΓ}
  The forward direction is \cref{lem:isLocalizedModule_restrict_of_isIso_fromTildeΓ} (the affine
  engine, applied for each \(f\)); the reverse direction is
  \cref{lem:isIso_fromTildeΓ_of_isLocalizedModule_restrict}.
\end{proof}

\begin{definition}

\leanok
[Schematic support of a sheaf of modules]
  \label{def:schematic_support}
  \lean{AlgebraicGeometry.Scheme.Modules.schematicSupport}
  \uses{def:modules_annihilator, lem:modules_annihilator_ideal_le,
    def:schematic_support_immersion}
  % SOURCE: [Nitsure], \S 1, sub-section "Stratification by Hilbert
  % Polynomials" (read from
  % references/nitsure-hilbert-quot-src/nitsure-hilbert-quot.tex, L468-L471).
  % SOURCE QUOTE:
  % "Let $X\to S$ be a finite type morphism of noetherian schemes,
  %  and let $L$ be a line bundle on $X$. Let $\F$
  %  be any coherent sheaf on
  %  $X$ whose schematic support is proper over $S$."
  \textit{Source: [Nitsure], \S 1 (FGA Explained Ch.~5).}
  Let \(\mathcal{F}\) be a sheaf of modules on a scheme \(X\). The
  \emph{schematic support} of \(\mathcal{F}\) is the closed subscheme of \(X\)
  cut out by its annihilator ideal sheaf (\cref{def:modules_annihilator}):
  \[
    \mathrm{schematicSupport}(\mathcal{F})
    \;=\; V\bigl(\mathrm{Ann}(\mathcal{F})\bigr) \;\hookrightarrow\; X,
  \]
  equipped with its canonical closed immersion into \(X\). This realises
  Nitsure's ``schematic support of \(\mathcal{F}\)'' as a closed subscheme of
  \(X\).
\end{definition}

\begin{definition}\leanok
[Schematic-support immersion]
  \label{def:schematic_support_immersion}
  \lean{AlgebraicGeometry.Scheme.Modules.schematicSupportι}
  \uses{def:schematic_support, def:modules_annihilator}
  Let \(\mathcal{F}\) be a sheaf of modules on a scheme \(X\). The
  \emph{schematic-support immersion} is the canonical closed immersion
  \[
    \mathrm{schematicSupport}(\mathcal{F}) \;\hookrightarrow\; X
  \]
  of the schematic support (\cref{def:schematic_support}) into the ambient
  scheme, namely the subscheme immersion of the annihilator ideal sheaf
  \(\mathrm{Ann}(\mathcal{F})\) (\cref{def:modules_annihilator}). It is a
  quasi-compact preimmersion onto the support, and it is the morphism whose
  composite with a base morphism \(f : X \to S\) defines proper support
  (\cref{def:has_proper_support}). Proved directly in Lean.
\end{definition}

\begin{lemma}[Properness as a morphism property]
  \label{lem:isProper_mathlib}
  \lean{AlgebraicGeometry.IsProper}
  \mathlibok
  \textit{Provided by Mathlib
  (\texttt{Mathlib.AlgebraicGeometry.Morphisms.Proper}).}
  A morphism \(f : X \to S\) of schemes is \emph{proper} iff it is separated,
  universally closed, and locally of finite type. Properness is stable under base
  change: if \(f : X \to S\) is proper and \(g : T \to S\) is any morphism then
  the base change \(X \times_S T \to T\) is proper. Consequently the base-change
  clause that Nitsure invokes (``properness \ldots\ preserved by base-change'')
  applies directly to the proper-support condition below.
\end{lemma}

\begin{definition}

\leanok
[Proper support over the base]
  \label{def:has_proper_support}
  \lean{AlgebraicGeometry.Scheme.Modules.HasProperSupport}
  \uses{def:schematic_support, def:schematic_support_immersion, lem:isProper_mathlib}
  % SOURCE: [Nitsure], \S 1, sub-section "The Functors $\Hilb_{X/S}$ and
  % $\Quot_{E/X/S}$" (read from
  % references/nitsure-hilbert-quot-src/nitsure-hilbert-quot.tex, L407-L409
  % + L418-L419).
  % SOURCE QUOTE:
  % "(i) a coherent sheaf $\F$ on $X_T = X\times_ST$
  %  such that the schematic support of $\F$ is proper over $T$
  %  and $\F$ is flat over $T$, together with
  %  [...]
  %  Then as properness and flatness are preserved
  %  by base-change, and as tensor-product is right exact,"
  \textit{Source: [Nitsure], \S 1 (FGA Explained Ch.~5).}
  Let \(\mathcal{F}\) be a sheaf of modules on a scheme \(X\) and let
  \(f : X \to S\) be a morphism. We say \(\mathcal{F}\) \emph{has proper support
  over \(S\) (along \(f\))} when the composite morphism
  \[
    \mathrm{schematicSupport}(\mathcal{F}) \hookrightarrow X \xrightarrow{f} S
  \]
  is proper (\cref{lem:isProper_mathlib}). Since properness is stable under base
  change (\cref{lem:isProper_mathlib}), this condition is preserved by pull-back
  along any \(S'\to S\), which is exactly what the well-definedness of the Quot
  functor's pullback action requires.
\end{definition}

\begin{definition}

\leanok
[Locally free of rank \(d\)]
  \label{def:is_locally_free_of_rank}
  \lean{AlgebraicGeometry.SheafOfModules.IsLocallyFreeOfRank}
  % SOURCE: [Nitsure], \S 1, "Elementary Examples, Exercises (2)" and
  % "Grassmannian of a vector bundle" (read from
  % references/nitsure-hilbert-quot-src/nitsure-hilbert-quot.tex, L562-L564
  % + L584-L585).
  % SOURCE QUOTE:
  % "$\langle \F, q\rangle$ of quotients $q: \oplus^r\,{\mathcal O}_T \to \F$
  %  where $\F$ is a locally free sheaf on $T$ of rank $d$.
  %  [...]
  %  the set of all equivalence classes $\langle \F, q\rangle$ of quotients
  %  $q: E_T \to \F$ where $\F$ is a locally free sheaf on $T$ of rank $d$,"
  \textit{Source: [Nitsure], \S 1, Exercise (2) (FGA Explained Ch.~5).}
  Let \(\mathcal{M}\) be a sheaf of modules on a scheme \(X\) and let
  \(d \in \mathbb{N}\). We say \(\mathcal{M}\) is \emph{locally free of rank
  \(d\)} when there exists an open cover \(\{U_i\}\) of \(X\) together with
  isomorphisms \(\mathcal{M}|_{U_i} \cong \mathcal{O}_{U_i}^{\oplus d}\) for
  each \(i\).

  \emph{Note.} Mathlib does not provide this predicate.
\end{definition}

\section{The Quot functor}
\label{sec:quot_functor}

\begin{definition}\leanok
[Quot functor]
  \label{def:quot_functor}
  \lean{AlgebraicGeometry.Scheme.QuotFunctor}
  \uses{def:hilbert_polynomial, def:has_proper_support}

  % SOURCE: [Nitsure], \S 1, sub-section "The Functors \(\Hilb_{X/S}\) and
  % \(\Quot_{E/X/S}\)" (read from
  % references/nitsure-hilbert-quot-src/nitsure-hilbert-quot.tex, L394-L433
  % + L481-L498).
  % SOURCE QUOTE:
  % "Let \(S\) be a noetherian scheme and let \(X\to S\) be a finite type
  %  scheme over it. Let \(E\) be a coherent sheaf on \(X\).
  %  Let \(Sch_S\) denote the category of all locally noetherian schemes
  %  over \(S\). For any \(T\to S\) in \(Sch_S\), a \textbf{family of quotients of \(E\)
  %  parametrised by \(T\)} will mean
  %  a pair \((\F,q)\) consisting of
  %
  %  (i) a coherent sheaf \(\F\) on \(X_T = X\times_ST\)
  %  such that the schematic support of \(\F\) is proper over \(T\)
  %  and \(\F\) is flat over \(T\), together with
  %
  %  (ii) a surjective \({\mathcal O}_{X_T}\)-linear homomorphism of sheaves
  %  \(q:E_T  \to \F\) where \(E_T\) is the pull-back of \(E\) under
  %  the projection \(X_T \to X\).
  %
  %  Two such families \((\F,q)\) and \((\F,q)\) parametrised by \(T\) will
  %  be regarded as equivalent if \(\ker(q) = \ker(q')\)),
  %  and \(\langle \F, q \rangle\) will denote
  %  an equivalence class. Then as properness and flatness are preserved
  %  by base-change, and as tensor-product is right exact,
  %  the pull-back of \(\langle \F, q \rangle\) under an \(S\)-morphism \(T'\to T\) is
  %  well-defined, which gives
  %  a set-valued contravariant functor
  %  \(\Quot_{E/X/S} : Sch_S \to Sets\) under which
  %  \(\)T\mapsto  \{\mbox{ All } \langle \F, q \rangle \mbox{ parametrised by }T\}\(\)
  %  [...]
  %  This shows that the functor \(\Quot_{E/X/S}\) naturally decomposes
  %  as a co-product
  %  \(\)\Quot_{E/X/S} = \coprod_{{\Phi}\in \Q[\lambda]}\, \Quot_{E/X/S}^{\Phi,L}\(\)
  %  where for any polynomial \({\Phi}\in \Q[\lambda]\), the functor
  %  \(\Quot_{E/X/S}^{\Phi,L}\) associates
  %  to any \(T\) the set of all equivalence classes of
  %  families \(\langle \F,q \rangle\) such that at each \(t\in T\) the
  %  Hilbert polynomial of the restriction \(\F_t\), calculated
  %  using the pull-back of \(L\), is \({\Phi}\)."
  \textit{Source: [Nitsure], \S 1 (FGA Explained Ch.~5).}
  Let \(S\) be a noetherian scheme, \(\pi : X \to S\) a finite-type morphism,
  \(L\) a line bundle on \(X\), \(E\) a coherent sheaf on \(X\), and
  \(\Phi \in \mathbb{Q}[\lambda]\). The \emph{Quot functor with Hilbert
  polynomial \(\Phi\)}
  \[
    \Quot^{\Phi,L}_{E/X/S} : (\mathrm{Sch}/S)^{op} \longrightarrow \mathbf{Set}
  \]
  sends an \(S\)-scheme \(T \to S\) to the set of equivalence classes
  \(\langle \mathcal{F}, q\rangle\) of pairs \((\mathcal{F}, q)\) where
  \begin{enumerate}
    \item \(\mathcal{F}\) is a coherent sheaf on \(X_T = X \times_S T\) whose
      schematic support is proper over \(T\) (\cref{def:has_proper_support}) and
      which is flat over \(T\);
    \item \(q : E_T \twoheadrightarrow \mathcal{F}\) is a surjective
      \(\mathcal{O}_{X_T}\)-linear homomorphism, where \(E_T\) is the pull-back
      of \(E\) along \(X_T \to X\);
    \item for every \(t \in T\) the Hilbert polynomial of \(\mathcal{F}|_{X_t}\)
      with respect to \(L|_{X_t}\) (\cref{def:hilbert_polynomial}) equals
      \(\Phi\).
  \end{enumerate}
  Two pairs \((\mathcal{F}, q)\) and \((\mathcal{F}', q')\) are equivalent iff
  \(\ker(q) = \ker(q')\). On morphisms \(f : T' \to T\) over \(S\), the functor
  acts by pulling back the quotient along \(X_{T'} \to X_T\) (well-defined
  since tensor product is right-exact and preserves both flatness and
  properness). The Hilbert scheme is the special case
  \(\Hilb^{\Phi,L}_{X/S} := \Quot^{\Phi,L}_{\mathcal{O}_X/X/S}\).
\end{definition}

\section{The Grassmannian scheme}
\label{sec:grassmannian_scheme}

The Quot construction proceeds by embedding \(\Quot^{\Phi,L}_{E/X/S}\) as a
locally closed subfunctor of a Grassmannian over \(S\). Since Mathlib carries
no Grassmannian \emph{scheme} (only a finite-rank-quotient variety in
linear algebra), we package the construction here.

\begin{definition}\leanok
[Grassmannian scheme]
  \label{def:grassmannian_scheme}
  \lean{AlgebraicGeometry.Scheme.Grassmannian}
  \uses{def:quot_functor, def:is_locally_free_of_rank}
  % SOURCE: [Nitsure], \S 1, "Elementary Examples, Exercises~(2)" and
  % "Construction of Grassmannian" (read from
  % references/nitsure-hilbert-quot-src/nitsure-hilbert-quot.tex, L546-L599
  % + L773-L850).
  % SOURCE QUOTE:
  % "For any integers \(r\ge d\ge 1\), an explicit construction
  %  the Grassmannian scheme
  %  \(\grass(r,d)\) over \(\Z\), together with the tautological
  %  quotient \(u:\oplus^r\,{\mathcal O}_{\grass(r,d)} \to \U\) where \
  %  \(\U\) is a rank \(d\) locally free sheaf on \(\grass(r,d)\),
  %  has been given at the end of this section. [...]
  %  Show that \(\grass(r,d)\) together with the quotient
  %  \(u:\oplus^r\,{\mathcal O}_{\grass(r,d)} \to \U\) represents the
  %  contravariant functor
  %  \(\)\Grass(r,d) = \Quot_{\oplus^r\,{\mathcal O}_{\Z}/\Z/\Z}^{d, {\mathcal O}_{\Z}}\(\)
  %  from schemes to sets, which associates to any \(T\) the
  %  set of all equivalence classes
  %  \(\langle \F, q\rangle\) of quotients \(q: \oplus^r\,{\mathcal O}_T \to \F\)
  %  where \(\F\) is a locally free sheaf on \(T\) of rank \(d\).
  %  [...]
  %  Using the above show that, more generally, if \(S\) is a scheme and
  %  \(E\) is a locally free \({\mathcal O}_S\)-module of rank \(r\),
  %  the functor \(\Grass(E,d)=\Quot_{E/S/S}^{d,{\mathcal O}_S}\)
  %  on all \(S\)-schemes which
  %  by definition associates to any \(T\)
  %  the set of all equivalence classes \(\langle \F, q\rangle\) of quotients
  %  \(q: E_T \to \F\) where \(\F\) is a locally free sheaf on \(T\) of rank \(d\),
  %  is representable. The representing scheme
  %  is denoted by \(\grass(E,d)\) and is called the rank \(d\)
  %  relative Grassmannian of \(E\) over \(S\)."
  \textit{Source: [Nitsure], \S 1, Exercise (2) and "Construction of
  Grassmannian" (FGA Explained Ch.~5).}
  Let \(S\) be a noetherian scheme and let \(V\) be a locally free
  \(\mathcal{O}_S\)-module of rank \(r\). For an integer \(1 \le d \le r\), the
  \emph{Grassmannian of rank-\(d\) quotients of \(V\)} is the functor
  \[
    \mathrm{Grass}(V, d) : (\mathrm{Sch}/S)^{op} \longrightarrow \mathbf{Set},
    \qquad
    T \longmapsto \{ \langle \mathcal{F}, q\rangle :
    q : V_T \twoheadrightarrow \mathcal{F},\; \mathcal{F}
    \text{ locally free of rank } d \}
  \]
  with equivalence \(\ker(q) = \ker(q')\); the rank-\(d\) local freeness of
  \(\mathcal{F}\) is as in \cref{def:is_locally_free_of_rank}. Concretely,
  \(\mathrm{Grass}(V, d) = \Quot^{d,\mathcal{O}_S}_{V/S/S}\): the Quot functor
  (\cref{def:quot_functor}) in the special case \(X = S\), \(E = V\), constant
  Hilbert polynomial \(\Phi = d\).
\end{definition}

\begin{lemma}[Representability predicate]
  \label{lem:functor_is_representable_mathlib}
  \lean{CategoryTheory.Functor.IsRepresentable}
  \mathlibok
  \textit{Provided by Mathlib.}
  Let \(\mathcal{C}\) be a category and \(F : \mathcal{C}^{op} \to \mathbf{Set}\)
  a presheaf of sets. The predicate \(F.\mathrm{IsRepresentable}\) asserts that
  \(F\) is representable, i.e.\ that there exists an object \(Y\) of
  \(\mathcal{C}\) together with a natural bijection
  \(\mathrm{Hom}_{\mathcal{C}}(-, Y) \cong F\) (the Yoneda bijection).
\end{lemma}

\begin{theorem}\leanok
[Representability of the Grassmannian]
  \label{thm:grassmannian_representable}
  \lean{AlgebraicGeometry.Scheme.Grassmannian.representable}
  \uses{def:grassmannian_scheme, thm:relative_spec_exists, thm:relative_spec_univ,
    lem:functor_is_representable_mathlib, def:gr_glued_scheme, lem:gr_separated,
    lem:gr_proper}

  % NOTE (representability proof blocked): The proof sketch recovers the representing object
  % through the RepresentableBy Yoneda form via \cref{thm:relative_spec_univ}. However,
  % \cref{thm:relative_spec_univ} is currently established in Lean in a form weaker than a full
  % \(\mathrm{Functor.RepresentableBy}\) witness. The proof is blocked on either (a)
  % strengthening \cref{thm:relative_spec_univ} to deliver a RepresentableBy witness, or (b)
  % finding a RepresentableBy-free argument via direct gluing of chart-wise classifying morphisms.
  % This is a deferred open question.
  % NOTE: The Lean statement \lean{AlgebraicGeometry.Scheme.Grassmannian.representable} is
  % currently a weakened existence skeleton that omits smoothness, properness, relative dimension
  % \(d(r-d)\), the tautological rank-\(d\) quotient, and the Pl\"ucker embedding. The
  % \verb|\lean{}| pin above points at a declaration that under-delivers the prose statement; it
  % should be strengthened or split into a separate skeleton label.

  % SOURCE: [Nitsure], \S 1, sub-section "Construction by gluing together
  % affine patches" through "Projective embedding"
  % (read from references/nitsure-hilbert-quot-src/nitsure-hilbert-quot.tex,
  % L798-L930).
  % SOURCE QUOTE:
  % "For any integers \(r\ge d\ge 1\), the Grassmannian scheme
  %  \(\grass(r,d)\) over \(\Z\), together with the tautological
  %  quotient \(u:\oplus^r\,{\mathcal O}_{\grass(r,d)} \to \U\) where \
  %  \(\U\) is a rank \(d\) locally free sheaf on \(\grass(r,d)\),
  %  can be explicitly constructed as follows. [...]
  %  the schemes
  %  \(U^I\), as \(I\) varies over all the
  %  \({r\choose d}\) different
  %  subsets of  \(\{ 1,\ldots, r\}\) of cardinality \(d\),
  %  can be glued together by the co-cycle \((\theta_{I,J})\) to form a
  %  finite-type scheme \(\grass(r,d)\) over \(\Z\). As each \(U^I\) is
  %  isomorphic to \(\AA^{d(r-d)}_{\Z}\), it follows that
  %  \(\grass(r,d) \to \spec \Z\) is smooth of relative dimension \(d(r-d)\).
  %  [...]
  %  We now show that \(\pi : \grass(r,d) \to \spec \Z\) is proper. It is enough to
  %  verify the valuative criterion of properness for discrete valuation rings.
  %  [...]
  %  As \(\U\) is given by the transition functions \(g_{I,J}\) described above,
  %  the determinant line bundle
  %  \(\det(\U)\) is given by the transition functions
  %  \(\det(g_{I,J}) = 1/P^I_J \in GL_1(U^I_J)\).
  %  [...]
  %  We leave it to the reader to verify that the sections \(\sigma_I\) form
  %  a linear system which is base point free and separates points
  %  relative to \(\spec \Z\), and so gives an embedding of
  %  \(\grass(r,d)\) into \(\PP^m_{\Z}\) where \(m = {r\choose d}-1\).
  %  This is a closed embedding by the properness of
  %  \(\pi:\grass(r,d) \to \spec \Z\). In particular, \(\det(\U)\)
  %  is a relatively very ample line bundle on \(\grass(r,d)\) over \(\Z\)."
  \textit{Source: [Nitsure], \S 1, "Construction of Grassmannian" (FGA
  Explained Ch.~5).}
  Let \(S\) be a noetherian scheme, \(V\) a locally free \(\mathcal{O}_S\)-module
  of rank \(r\), and \(1 \le d \le r\). The functor \(\mathrm{Grass}(V,d)\) of
  \cref{def:grassmannian_scheme} is representable by a smooth projective
  \(S\)-scheme \(\mathrm{Gr}_S(V, d) \to S\) of relative dimension \(d(r-d)\),
  equipped with a tautological quotient
  \(\pi^* V \twoheadrightarrow \mathcal{U}\) of rank \(d\). The determinant
  line bundle \(\det(\mathcal{U})\) is relatively very ample, and the
  associated linear system gives a closed embedding
  \(\mathrm{Gr}_S(V, d) \hookrightarrow
  \mathbb{P}_S\!\left(\bigwedge^d V\right)\) -- the \emph{Pl\"ucker embedding}.

\end{theorem}

% SOURCE QUOTE PROOF: see the "Construction of Grassmannian" sub-section of
% [Nitsure], \S 1 (references/nitsure-hilbert-quot-src/nitsure-hilbert-quot.tex,
% L773-L930), which gives the construction in five named steps
% ("Construction by gluing together affine patches", "Separatedness",
% "Properness", "Universal quotient", "Projective embedding"). The total
% verbatim text is ~150 lines; it is reproduced piecewise in the
% SOURCE QUOTE comment of the theorem statement above.
\begin{proof}
  
  Over \(S = \Spec \mathbb{Z}\) and \(V = \mathcal{O}_S^{\oplus r}\) the
  Grassmannian \(\mathrm{Gr}(r, d)\) is constructed by gluing
  \(\binom{r}{d}\) affine charts \(U^I = \Spec\mathbb{Z}[X^I]\) indexed by
  size-\(d\) subsets \(I \subseteq \{1, \dots, r\}\). The chart \(U^I\)
  carries the matrix \(X^I\) of size \(d \times r\) whose \(I\)-th \(d \times d\)
  minor is the identity, with the remaining \(d(r-d)\) entries free
  variables (so \(U^I \cong \mathbb{A}^{d(r-d)}_{\mathbb{Z}}\)). The
  transition maps \(\theta_{I,J} : U^I \cap U^J \to U^J \cap U^I\) defined
  by \(X^J \mapsto (X^I_J)^{-1} X^I\) on the locus where \(\det X^I_J\) is
  invertible satisfy the cocycle condition
  \(\theta_{I,K} = \theta_{I,J} \circ \theta_{J,K}\). The glued
  \(\mathrm{Gr}(r,d) \to \Spec\mathbb{Z}\) is smooth of relative dimension
  \(d(r-d)\). Separatedness is checked from the diagonal
  \(\Delta_{I,J} \subset U^I \times U^J\) cut out by \(X^J_I X^I - X^J = 0\).
  Properness follows from the valuative criterion for DVRs: given
  \(\varphi : \Spec K \to \mathrm{Gr}(r,d)\) over a DVR \(R\) with fraction
  field \(K\), factoring through some chart \(U^I\) via a ring map
  \(f : \mathbb{Z}[X^I] \to K\), choose \(J\) minimising the valuation
  \(\nu(f(P^I_J))\) of the minor; on the chart \(U^J\) all entries have
  non-negative valuation, so \(f\) extends uniquely to \(\Spec R\).

  The universal quotient \(u : \mathcal{O}_{\mathrm{Gr}(r,d)}^{\oplus r}
  \twoheadrightarrow \mathcal{U}\) is given on \(U^I\) by the matrix \(X^I\)
  and glued by \(g_{I,J} = (X^I_J)^{-1} \in GL_d(U^I \cap U^J)\). The
  determinant line bundle \(\det\mathcal{U}\) has transition functions
  \(\det g_{I,J} = 1/P^I_J\); the global sections
  \(\sigma_I \in \Gamma(\mathrm{Gr}(r,d), \det\mathcal{U})\) defined by
  \(\sigma_I|_{U^J} = P^J_I\) form a base-point-free linear system separating
  points over \(\Spec\mathbb{Z}\), giving an embedding into
  \(\mathbb{P}^{\binom{r}{d}-1}_{\mathbb{Z}}\), which is closed by
  properness. Representability of \(\mathrm{Grass}(r,d)\) is then verified
  on the universal pair \((\mathrm{Gr}(r,d), u)\): a quotient
  \(q : \mathcal{O}_T^{\oplus r} \twoheadrightarrow \mathcal{F}\) of
  rank \(d\) on \(T\) is everywhere locally trivialised; on the locus where
  the columns indexed by \(I\) are a basis the matrix of \(q\) has the form
  \(X^I\), defining the classifying morphism \(T \to U^I \subset
  \mathrm{Gr}(r,d)\). Base change from \(\mathbb{Z}\) to a general noetherian
  base \(S\) with a vector bundle \(V\) is direct: trivialise \(V\) on an open
  cover and glue. For non-trivial \(V\) the same construction produces
  \(\mathrm{Gr}_S(V, d)\) with Pl\"ucker embedding into
  \(\mathbb{P}_S(\bigwedge^d V)\).
\end{proof}

\section{Out of scope}
\label{sec:quot_out_of_scope}

This chapter packages \emph{only} the Quot-foundations layer: the Hilbert
polynomial, the Quot functor, and the Grassmannian scheme with its
representability. The following are explicitly out of scope here:
\begin{itemize}
  \item \textbf{Representability of the Quot scheme} -- the
    Grothendieck--Altman--Kleiman existence theorem and its proof skeleton
    (boundedness, Grassmannian embedding, flattening stratification,
    valuative criterion) are not constructed in this chapter.
  \item \textbf{Flat base change of cohomology} (Stacks 02KH) and its
    associated prerequisites -- not constructed in this chapter.
  \item \textbf{Hilbert scheme as a separate chapter}. The Hilbert scheme
    \(\Hilb_{X/S}\) is the special case
    \(\Hilb_{X/S} = \Quot_{\mathcal{O}_X/X/S}\) of the Quot functor and is
    introduced inline (\cref{def:quot_functor}); no separate chapter.
  \item \textbf{Snapper's Lemma} -- the polynomial-eventually property of
    Euler characteristics, used in \cref{def:hilbert_polynomial}. Cited
    via Nitsure \S 1 and (textbook reference) Hartshorne III.5.2;
    Lean formalization is a prerequisite sub-task.
  \item \textbf{Identification with \(\mathbb{P}^n\) and relative
    Grassmannian variants} (Nitsure \S 1, Exercises (1)--(4)) --
    optional applications, not consumed by Route A.
  \item \textbf{Group-scheme structure on \(\Pic_{C/k}\)} (A.2.c) -- the
    abelian-group-scheme refinement is downstream. Its blueprint home is the
    FGA Picard-representability assembly (the chapter that bundles the relative
    Picard functor, its representability, and the abelian-group-scheme
    structure on \(\Pic_{C/k}\)); that assembly lives in the parent project
    which consumes this Quot-foundations layer and is not part of the present
    subproject.
  \item \textbf{Identity-component \(\Pic^0\) and abelian-variety
    structure} (A.3) -- a separate Route A sub-row, fed by but not
    constructed in this chapter.
\end{itemize}
